% Lesson 56 — Grammar: The active and passive voices — Anglais Tle A — Cours signé OnBuch+
% Programme officiel MINESEC. Compiler avec : tectonic EN56-grammar-the-active-and-passive-voices.tex
\newcommand{\DOCMATIERE}{Anglais}
\newcommand{\DOCNIVEAU}{Tle A}
\newcommand{\DOCMODULE}{Chapter 5 : Media and Communication}
\newcommand{\DOCLECON}{EN56}
\newcommand{\DOCTITRE}{Lesson 56 — Grammar: The active and passive voices}
% OnBuch+ — préambule « cours signés » ENRICHI (compilé avec tectonic / XeLaTeX)
% Système visuel « Pop » d'OnBuch+ : fond crème (jamais blanc), bordures encre
% épaisses, ombres « dures » décalées, titres Archivo Black en capitales.
% Version MATHÉMATIQUES : graphes (pgfplots/tikz), boîtes pédagogiques
% (définition, propriété, méthode, attention, le savais-tu, exercice résolu,
% exercice, TP, bilan), page de garde et signature OnBuch+ ancrée en bas.
%
% Macros attendues dans main.tex AVANT \input{preamble} :
%   \DOCMATIERE  \DOCNIVEAU  \DOCMODULE  \DOCLECON (numéro)  \DOCTITRE
\documentclass[11pt]{article}

\usepackage{fontspec}
\usepackage[english,french]{babel} % français = langue principale ; l'anglais via \eng{..} / textebox
\usepackage[a4paper,margin=2cm,top=3.4cm,bottom=2.8cm,headheight=1.4cm,footskip=1.3cm]{geometry}
\usepackage{amsmath,amssymb,amsfonts}
\usepackage{mathtools} % \xmapsto, \coloneqq... souvent utilisés par les modèles
\usepackage{cancel} % simplifications barrées (bilans de forces)
% \abs{z} (module, valeur absolue) et \norm{u} (norme), utilisés par les modèles
\providecommand{\abs}{}\let\abs\relax\DeclarePairedDelimiter{\abs}{\lvert}{\rvert}
\providecommand{\norm}{}\let\norm\relax\DeclarePairedDelimiter{\norm}{\lVert}{\rVert}
\usepackage[table]{xcolor} % [table] : fournit aussi \rowcolors (alternance de lignes)
\usepackage{pagecolor}
\usepackage{tikz}
\usetikzlibrary{arrows.meta,positioning,calc,shapes.geometric,shapes.misc,shapes.arrows,decorations.pathmorphing,decorations.pathreplacing,decorations.markings,patterns,shadows,mindmap,trees,fit,backgrounds,matrix,intersections,angles,quotes}
\usepackage{pgfplots}
\usepackage[version=4]{mhchem} % équations nucléaires \ce{^{235}_{92}U}
\usepackage[european,siunitx]{circuitikz} % schémas électriques
\pgfplotsset{compat=1.18}
\usepgfplotslibrary{fillbetween,groupplots} % aires entre courbes (calcul intégral)
\usepackage{enumitem}
\usepackage{fancyhdr}
% [most] charge toutes les bibliothèques tcolorbox (listings, minted, raster,
% documentation...) et gonfle inutilement la mémoire TeX sur un document long
% (~300 boîtes) : on ne charge que ce qui est réellement utilisé.
\usepackage[skins,breakable]{tcolorbox}
\usepackage{graphicx}
\usepackage{colortbl}
\usepackage{booktabs}
\usepackage{tabularx}
\usepackage{diagbox} % case d'en-tête diagonale des tableaux à double entrée
% Colonnes extensibles centrée / gauche / droite (C, L, R), souvent utilisées par les modèles
\newcolumntype{C}{>{\centering\arraybackslash}X}
\newcolumntype{L}{>{\raggedright\arraybackslash}X}
\newcolumntype{R}{>{\raggedleft\arraybackslash}X}
\usepackage{array}
\usepackage{multicol}
\usepackage{multirow}
\usepackage{siunitx}
% parse-numbers=false : accepte aussi \SI{2,5 \times 10^{-3}}{...}
\sisetup{locale=FR,output-decimal-marker={,},group-minimum-digits=4,parse-numbers=false}
\DeclareSIUnit\ohm{\ensuremath{\symup{\Omega}}} % U+2126 absent de la police
\DeclareSIUnit\division{div} % réglages d'oscilloscope (ms/div, V/div)
\DeclareSIUnit\tr{tr} % tours (vitesse de rotation, ex. \SI{1500}{\tr\per\minute})
\DeclareSIUnit\molar{M} % concentration molaire (ex. \SI{3}{\molar}, \SI{1}{\micro\molar})
\DeclareSIUnit\an{an} % année (le modèle confond parfois avec \year, un compteur TeX) — ex. \SI{5}{\kilo\meter\per\an}
\DeclareSIUnit\FCFA{FCFA} % franc CFA (ex. \SI{500}{\FCFA\per\kilo\gram})
\DeclareSIUnit\degreeBrix{\textdegree Brix} % teneur en sucre d'un sirop/jus (réfractométrie)
\DeclareSIUnit\month{mois} % le modèle utilise parfois \month, un compteur TeX, comme unité de durée
\usepackage{qrcode}
% \degree : commande fréquemment utilisée par les modèles pour un angle ou une
% température en dehors de \SI{}{} (non fournie par les paquets chargés).
\providecommand{\degree}{^{\circ}}
% \qed (fin de démonstration), utilisé par les modèles sans amsthm
\providecommand{\qed}{\hfill$\square$}
% \barre : double barre des valeurs interdites dans un tableau de variations (tkz-tab)
\providecommand{\barre}{\ensuremath{\|}}
% environnement proof (amsthm non chargé) utilisé par les modèles
\ifdefined\proof\else\newenvironment{proof}[1][Démonstration]{\par\noindent\textit{#1.}\ }{\qed\par}\fi
\usepackage{lastpage}
\usepackage{hyperref}
\hypersetup{hidelinks}

% ============================================================
% Palette « Pop » OnBuch+ — mêmes hex que la classe P (plus_ui.dart)
% ============================================================
\definecolor{poporange}{HTML}{F59321}
\definecolor{popdark}{HTML}{E07A0C}
\definecolor{popcream}{HTML}{FFF6E8}
\definecolor{popink}{HTML}{1C1714}
\definecolor{poppurple}{HTML}{7A5AE0}
\definecolor{popgreen}{HTML}{2E9E6A}
\definecolor{poppink}{HTML}{E0457B}
\definecolor{popblue}{HTML}{2F7DE1}
\definecolor{popgold}{HTML}{C98A12}
\definecolor{popmuted}{HTML}{8A7F76}
\definecolor{popline}{HTML}{EADFD2}
\definecolor{popgray}{HTML}{8A7F76}
\colorlet{popgrey}{popgray}
% Alias tolérants (le modèle nomme parfois les couleurs différemment)
\colorlet{popwhite}{white}
\colorlet{popblack}{popink}
\colorlet{popred}{poppink}
\colorlet{popyellow}{popgold}
% Teintes claires (fonds de tableaux/encadrés) — le modèle nomme parfois
% les couleurs avec un suffixe L (ex. popblueL) sans les avoir définies.
\colorlet{poporangeL}{poporange!20}
\colorlet{popdarkL}{popdark!20}
\colorlet{poppurpleL}{poppurple!20}
\colorlet{popgreenL}{popgreen!20}
\colorlet{poppinkL}{poppink!20}
\colorlet{popblueL}{popblue!20}
\colorlet{popgoldL}{popgold!20}
\colorlet{popredL}{popred!20}
\colorlet{popgrayL}{popgray!20}
% Teintes claires pour fonds de tableaux / zones
\colorlet{popblueL}{popblue!10!white}
\colorlet{poporangeL}{poporange!14!white}
\colorlet{popgreenL}{popgreen!12!white}
\colorlet{poppurpleL}{poppurple!12!white}
\colorlet{poppinkL}{poppink!12!white}
\colorlet{popgoldL}{popgold!14!white}

\pagecolor{popcream}

% ============================================================
% Typographie
% ============================================================
\defaultfontfeatures{Ligatures=TeX}
\setmainfont{PlusJakartaSans-Medium.ttf}[Path=../../../fonts/,BoldFont=PlusJakartaSans-Bold.ttf]
% Symboles, API et emojis absents de la police principale : repli sur DejaVu Sans (caractères actifs)
\newfontface\symfont{DejaVuSans.ttf}[Path=../../../fonts/]
\makeatletter
\newcommand\symactive[1]{\catcode#1=13 \begingroup\lccode`\~=#1 \lowercase{\endgroup\def~}{{\symfont\char#1}}}
\newcommand\symblank[1]{\catcode#1=13 \begingroup\lccode`\~=#1 \lowercase{\endgroup\def~}{}}
\newcommand\symhyph[1]{\catcode#1=13 \begingroup\lccode`\~=#1 \lowercase{\endgroup\def~}{\mbox{-}}}
\newcount\symc
\newcommand\symrange[2]{\symc=#1 \@whilenum\symc<#2 \do{\expandafter\symactive\expandafter{\the\symc}\advance\symc by 1 }}
\newcommand\symblankrange[2]{\symc=#1 \@whilenum\symc<#2 \do{\expandafter\symblank\expandafter{\the\symc}\advance\symc by 1 }}
\makeatother
\symrange{"0250}{"02FF}\symactive{"2713}\symactive{"2714}\symactive{"2717}\symactive{"2718}\symactive{"2610}\symactive{"2611}\symactive{"2612}
\symhyph{"2011}\symhyph{"2E1A}\symblankrange{"1F300}{"1FAFF}\symblank{"FE0F}
% Maths en sans-serif (Fira Math) avec chiffres et lettres droites de Plus
% Jakarta, pour que les formules se fondent dans le texte.
\usepackage[math-style=ISO,bold-style=ISO]{unicode-math}
\setmathfont{FiraMath-Regular.otf}[Path=../../../fonts/]
\setmathfont{PlusJakartaSans-Medium.ttf}[Path=../../../fonts/,range={up/num,up/{Latin,latin},\mathup}]
\usepackage{icomma} % virgule décimale sans espace en mode math
% Caractères mathématiques Unicode absents des polices de texte (Plus Jakarta,
% Archivo Black) : rendus par leur équivalent mathématique, en texte comme en
% formule — sinon ils s'affichent en carré vide (ex. « Arithmétique dans ℤ »).
\usepackage{newunicodechar}
\newunicodechar{ℝ}{\ensuremath{\mathbb{R}}}
\newunicodechar{ℤ}{\ensuremath{\mathbb{Z}}}
\newunicodechar{ℂ}{\ensuremath{\mathbb{C}}}
\newunicodechar{ℕ}{\ensuremath{\mathbb{N}}}
\newunicodechar{ℚ}{\ensuremath{\mathbb{Q}}}
\newunicodechar{𝔻}{\ensuremath{\mathbb{D}}}
\newunicodechar{α}{\ensuremath{\alpha}}
\newunicodechar{β}{\ensuremath{\beta}}
\newunicodechar{γ}{\ensuremath{\gamma}}
\newunicodechar{δ}{\ensuremath{\delta}}
\newunicodechar{ε}{\ensuremath{\varepsilon}}
\newunicodechar{θ}{\ensuremath{\theta}}
\newunicodechar{λ}{\ensuremath{\lambda}}
\newunicodechar{μ}{\ensuremath{\mu}}
\newunicodechar{σ}{\ensuremath{\sigma}}
\newunicodechar{τ}{\ensuremath{\tau}}
\newunicodechar{φ}{\ensuremath{\varphi}}
\newunicodechar{ψ}{\ensuremath{\psi}}
\newunicodechar{ω}{\ensuremath{\omega}}
\newunicodechar{Σ}{\ensuremath{\Sigma}}
% majuscules produites par \MakeUppercase dans les titres (α-aminés -> Α)
\newunicodechar{Α}{\ensuremath{\alpha}}
\newunicodechar{Β}{\ensuremath{\beta}}
\newunicodechar{Γ}{\ensuremath{\Gamma}}
\newunicodechar{Δ}{\ensuremath{\Delta}}
\newunicodechar{Ε}{\ensuremath{\varepsilon}}
\newunicodechar{Θ}{\ensuremath{\Theta}}
\newunicodechar{Λ}{\ensuremath{\Lambda}}
\newunicodechar{Μ}{\ensuremath{\mu}}
\newunicodechar{Τ}{\ensuremath{\tau}}
\newunicodechar{Φ}{\ensuremath{\Phi}}
\newunicodechar{Ψ}{\ensuremath{\Psi}}
\newunicodechar{Ω}{\ensuremath{\Omega}}
\newunicodechar{↦}{\ensuremath{\mapsto}}
\newunicodechar{∈}{\ensuremath{\in}}
\newunicodechar{∉}{\ensuremath{\notin}}
\newunicodechar{∧}{\ensuremath{\wedge}}
\newunicodechar{⇔}{\ensuremath{\Leftrightarrow}}
\newunicodechar{⟺}{\ensuremath{\Longleftrightarrow}}
\newunicodechar{⇒}{\ensuremath{\Rightarrow}}
\newunicodechar{⟹}{\ensuremath{\Longrightarrow}}
\newunicodechar{∘}{\ensuremath{\circ}}
\newunicodechar{∪}{\ensuremath{\cup}}
\newunicodechar{∩}{\ensuremath{\cap}}
\newunicodechar{≡}{\ensuremath{\equiv}}
\newunicodechar{⊂}{\ensuremath{\subset}}
\newunicodechar{⊥}{\ensuremath{\perp}}
\newunicodechar{∃}{\ensuremath{\exists}}
\newunicodechar{∀}{\ensuremath{\forall}}
\newunicodechar{∛}{\ensuremath{\sqrt[3]{\,}}}
\newunicodechar{∜}{\ensuremath{\sqrt[4]{\,}}}
\newunicodechar{⃗}{\ensuremath{{}^{\to}}}
\newunicodechar{ⁿ}{\ifmmode^{n}\else\textsuperscript{\ensuremath{n}}\fi}
\newunicodechar{ˣ}{\ifmmode^{x}\else\textsuperscript{\ensuremath{x}}\fi}
\newunicodechar{ᵇ}{\ifmmode^{b}\else\textsuperscript{\ensuremath{b}}\fi}
\newunicodechar{ʳ}{\ifmmode^{r}\else\textsuperscript{\ensuremath{r}}\fi}
\newunicodechar{ᵘ}{\ifmmode^{u}\else\textsuperscript{\ensuremath{u}}\fi}
\newunicodechar{ʸ}{\ifmmode^{y}\else\textsuperscript{\ensuremath{y}}\fi}
\newunicodechar{ᵃ}{\ifmmode^{a}\else\textsuperscript{\ensuremath{a}}\fi}
\newunicodechar{ᵉ}{\ifmmode^{e}\else\textsuperscript{\ensuremath{e}}\fi}
\newunicodechar{ᵏ}{\ifmmode^{k}\else\textsuperscript{\ensuremath{k}}\fi}
\newunicodechar{ᵖ}{\ifmmode^{p}\else\textsuperscript{\ensuremath{p}}\fi}
\newunicodechar{ᴺ}{\ifmmode^{N}\else\textsuperscript{\ensuremath{N}}\fi}
\newunicodechar{ᶜ}{\ifmmode^{c}\else\textsuperscript{\ensuremath{c}}\fi}
\newunicodechar{ʰ}{\ifmmode^{h}\else\textsuperscript{\ensuremath{h}}\fi}
\newunicodechar{ᵠ}{\ifmmode^{q}\else\textsuperscript{\ensuremath{q}}\fi}
\newunicodechar{ₙ}{\ifmmode_{n}\else\textsubscript{\ensuremath{n}}\fi}
\newunicodechar{ₐ}{\ifmmode_{a}\else\textsubscript{\ensuremath{a}}\fi}
\newunicodechar{ᵢ}{\ifmmode_{i}\else\textsubscript{\ensuremath{i}}\fi}
\newunicodechar{ᵣ}{\ifmmode_{r}\else\textsubscript{\ensuremath{r}}\fi}
\newunicodechar{ₖ}{\ifmmode_{k}\else\textsubscript{\ensuremath{k}}\fi}
\newunicodechar{ᵦ}{\ifmmode_{\beta}\else\textsubscript{\ensuremath{\beta}}\fi}
\newunicodechar{ₓ}{\ifmmode_{x}\else\textsubscript{\ensuremath{x}}\fi}
\newfontfamily\popdisplay{ArchivoBlack-Regular.ttf}[Path=../../../fonts/]
\newfontfamily\poplogo{SpaceGrotesk-Bold.ttf}[Path=../../../fonts/]
\newfontfamily\popsemi{PlusJakartaSans-SemiBold.ttf}[Path=../../../fonts/]
\newfontfamily\popxbold{PlusJakartaSans-ExtraBold.ttf}[Path=../../../fonts/]
\newfontfamily\popxboldls{PlusJakartaSans-ExtraBold.ttf}[Path=../../../fonts/,LetterSpace=3.0]

\setlist[enumerate]{leftmargin=1.7em,itemsep=5pt,topsep=4pt}
\setlist[itemize]{leftmargin=1.7em,itemsep=4pt,topsep=4pt,label={\color{poporange}\textbullet}}
\setlength{\parindent}{0pt}
\setlength{\parskip}{5pt}
\renewcommand{\arraystretch}{1.25}

% pgfplots : style Pop par défaut
\pgfplotsset{
  every axis/.append style={axis line style={popink,line width=1pt},tick style={popink},
    grid style={popline,line width=0.5pt},label style={font=\small},tick label style={font=\footnotesize}},
  popcurve/.style={poporange,line width=1.8pt,no markers,smooth},
}
% Même style disponible dans un \draw TikZ simple (hors pgfplots)
\tikzset{popcurve/.style={poporange,line width=1.8pt}}
\tikzset{popgrid/.style={popline,very thin}}
\colorlet{popcurve}{poporange} % le nom du style est parfois utilisé comme couleur

% ============================================================
% Pill / squiggle
% ============================================================
\newcommand{\poppill}[3]{%
  \tikz[baseline=-0.55ex]\node[draw=popink,line width=1.2pt,rounded corners=2.6pt,%
    fill=#2,inner xsep=6.5pt,inner ysep=3pt]{%
    \popxboldls\fontsize{7}{7}\selectfont\color{#3}\MakeUppercase{#1}};%
}
\newcommand{\popsquiggle}[1]{%
  \tikz[baseline=-0.4ex]{
    \draw[#1,line width=2.6pt,line cap=round]
      plot [smooth] coordinates {(0,0.09) (0.34,0.01) (0.68,0.21) (1.02,0.01) (1.36,0.10)};
  }%
}
% Mot-clé surligné (marqueur orange sous le texte)
\newcommand{\cle}[1]{{\popxbold\color{popdark}#1}}

% ============================================================
% En-tête / pied
% ============================================================
\pagestyle{fancy}
\fancyhf{}
\renewcommand{\headrulewidth}{0pt}
\renewcommand{\footrulewidth}{0pt}
\lhead{\raisebox{-0.15ex}{\poplogo\fontsize{13}{13}\selectfont\color{popink}OnBuch\color{poporange}+}}
\rhead{\poppill{\DOCMATIERE\ \textbullet\ \DOCNIVEAU\ \textbullet\ Leçon \DOCLECON}{white}{popink}}
\lfoot{\popxbold\fontsize{7.6}{7.6}\selectfont\color{popmuted}\MakeUppercase{Cours signé OnBuch+}\ \textbullet\ {\popsemi\DOCTITRE}}
\rfoot{\tikz[baseline=-0.6ex]\node[rounded corners=8.5pt,fill=popink,minimum height=17pt,minimum width=17pt,inner xsep=5pt,inner sep=0]{\popxbold\fontsize{8}{8}\selectfont\color{popcream}\thepage\,/\,\pageref*{LastPage}};}

% ============================================================
% Titres
% ============================================================
\newcommand{\chaptitre}[2]{%
  \begin{tcolorbox}[enhanced,colback=poporange,colframe=popink,boxrule=2.6pt,arc=11pt,
      drop shadow=popink,left=15pt,right=15pt,top=12pt,bottom=11pt,before skip=3pt,after skip=18pt]
    \poppill{#2}{popink}{popcream}\par\vspace{9pt}
    {\raggedright\hyphenpenalty=10000\exhyphenpenalty=10000\popdisplay\fontsize{22}{24}\selectfont\color{popcream}\MakeUppercase{#1}\par}
  \end{tcolorbox}
}
% Section numérotée : pastille orange avec le numéro + titre Archivo
\newcounter{popsec}
\newcommand{\coursec}[1]{%
  \stepcounter{popsec}\setcounter{popsub}{0}%
  \par\vspace{16pt}\noindent
  \begin{minipage}{\linewidth}
  \tikz[baseline=-0.7ex]\node[draw=popink,line width=1.6pt,fill=poporange,rounded corners=4pt,minimum size=22pt,inner sep=0]{\popdisplay\fontsize{12}{12}\selectfont\color{popcream}\thepopsec};%
  \hspace{8pt}{\popdisplay\fontsize{15}{17}\selectfont\color{popink}\MakeUppercase{#1}}\par
  \vspace{4pt}\noindent{\color{poporange}\rule{36pt}{3.6pt}}
  \end{minipage}\par\nopagebreak\vspace{8pt}
}
\newcounter{popsub}
\newcommand{\courssub}[1]{%
  \stepcounter{popsub}%
  \par\vspace{9pt}\noindent{\popxbold\fontsize{11.5}{13}\selectfont\color{popink}{\color{poporange}\thepopsec.\thepopsub}\hspace{6pt}#1}\par\nopagebreak\vspace{4pt}
}

% ============================================================
% Boîtes pédagogiques — même recette : fond blanc, bordure épaisse colorée,
% ombre dure encre, badge plein. Titre optionnel [#1].
% ============================================================
\tcbset{popbase/.style={breakable,enhanced,colback=white,boxrule=2.3pt,arc=9pt,
  shadow={3pt}{-3pt}{0pt}{popink},left=14pt,right=14pt,top=11pt,bottom=11pt,before skip=11pt,after skip=15pt,
  fontupper=\popsemi\fontsize{10.6}{14.5}\selectfont\color{popink},
  fontlower=\popsemi\fontsize{10.6}{14.5}\selectfont\color{popink}}}
% Coche dessinée (le glyphe ✓ n'existe pas dans les polices du document)
\newcommand{\popcheck}[1]{\tikz[baseline=-0.5ex]\draw[#1,line width=1.6pt,line cap=round,line join=round](-0.13,0.0)--(-0.04,-0.09)--(0.13,0.1);}
\providecommand{\coche}{\popcheck{popgreen}}
\providecommand{\resultat}[1]{\par\smallskip\noindent\fcolorbox{popgreen}{popgreenL}{\parbox{\dimexpr\linewidth-2\fboxsep-2\fboxrule\relax}{\textbf{Résultat.} #1}}\par\smallskip}
\newcommand{\popboxhead}[3]{\poppill{#1}{#2}{white}\ifx\relax#3\relax\else\hspace{6pt}{\popxbold\fontsize{10.6}{12}\selectfont\color{popink}#3}\fi\par\vspace{7pt}}

\newtcolorbox{aretenir}[1][]{popbase,colframe=popgreen,before upper={\popboxhead{À retenir}{popgreen}{#1}}}
% Texte en anglais (document de lecture, dialogue, extrait) : cadre
% d'étude. Un titre optionnel (source, auteur) s'affiche dans l'en-tête.
\newtcolorbox{textebox}[1][]{popbase,colframe=popblue,before upper={\popboxhead{Text}{popblue}{#1}\begin{otherlanguage}{english}},after upper={\end{otherlanguage}}}
% Marques ✓ / ✗ (forme correcte / fautive) souvent utilisées par les modèles
\providecommand{\cross}{\textcolor{poppink}{\ensuremath{\times}}}
\providecommand{\xmark}{\cross}\providecommand{\crossmark}{\cross}\providecommand{\cmark}{\ensuremath{\checkmark}}
% Ligne de réponse / justification à compléter (inventée par les modèles)
\providecommand{\justification}[1]{\par\noindent\textit{Justification~:} \dotfill\par}
% \textipa (package tipa non chargé) : la police ne couvre pas l'API -> simple passage
\providecommand{\textipa}[1]{#1}
% Soulignement (ulem non chargé) : \uline, \ul
\providecommand{\uline}[1]{\underline{#1}}\providecommand{\ul}[1]{\underline{#1}}
% \glossaire{\item ...} inventée par les modèles : liste de glossaire
\providecommand{\glossaire}[1]{\begin{itemize}#1\end{itemize}}
% \alert (beamer) inventée par les modèles : texte mis en évidence
\providecommand{\alert}[1]{\textbf{\textcolor{poppink}{#1}}}
% variantes ASCII d'ordinaux écrites en macro
\providecommand{\iere}{\textsuperscript{re}}\providecommand{\ieme}{\textsuperscript{e}}\providecommand{\ere}{\textsuperscript{re}}\providecommand{\eme}{\textsuperscript{e}}\providecommand{\er}{\textsuperscript{er}}\providecommand{\ie}{\textsuperscript{e}}
% Ordinaux français écrits en macro par les modèles : 1\ière, 3\ième
\newcommand{\ière}{\textsuperscript{re}}\newcommand{\ième}{\textsuperscript{e}}
% \exemple (inventée par les modèles) : étiquette « Exemple : »
\providecommand{\exemple}{\textbf{Exemple~:}\ }
% \square utilisable aussi hors mode math (cases à cocher)
\AtBeginDocument{\let\squareMath\square\renewcommand{\square}{\ensuremath{\squareMath}}}
% Mot ou phrase en anglais à l'intérieur d'un paragraphe français
\newcommand{\eng}[1]{\ifmmode\text{\textit{#1}}\else\begingroup\selectlanguage{english}\itshape #1\endgroup\fi}
\let\esp\eng % alias toléré
% flèches d'intonation inventées par les modèles
\providecommand{\textsearrow}{}\renewcommand{\textsearrow}{\ensuremath{\searrow}}\providecommand{\textnearrow}{}\renewcommand{\textnearrow}{\ensuremath{\nearrow}}
% \fr{..} : mot français (inventé par les modèles)
\providecommand{\fr}[1]{\textit{#1}}
\newtcolorbox{exemplebox}[1][]{popbase,colframe=poporange,before upper={\popboxhead{Exemple}{poporange}{#1}}}
\newtcolorbox{definition}[1][]{popbase,colframe=popblue,before upper={\popboxhead{Définition}{popblue}{#1}}}
\newtcolorbox{propriete}[1][]{popbase,colframe=poppurple,before upper={\popboxhead{Propriété}{poppurple}{#1}}}
\newtcolorbox{methode}[1][]{popbase,colframe=popgold,before upper={\popboxhead{Méthode}{popgold}{#1}}}
\newtcolorbox{attention}[1][]{popbase,colframe=poppink,before upper={\popboxhead{Attention}{poppink}{#1}}}
\newtcolorbox{experience}[1][]{popbase,colframe=popblue,colback=popblueL,before upper={\popboxhead{Expérience}{popblue}{#1}}}
% « Le savais-tu ? » : carte violette pleine (contexte, culture, Cameroun)
\newtcolorbox{savaistu}[1][]{popbase,colback=poppurple,colframe=popink,
  fontupper=\popsemi\fontsize{10.4}{14.2}\selectfont\color{white},
  before upper={\poppill{Le savais-tu\,?}{white}{poppurple}\ifx\relax#1\relax\else\hspace{6pt}{\popxbold\color{white}#1}\fi\par\vspace{7pt}}}
% Exercice résolu : énoncé en haut, \tcblower puis la solution sur fond crème
\newcounter{popexo}\newcounter{popexr}
\newtcolorbox{exoresolu}[1][]{popbase,colframe=poporange,colbacklower=popcream,
  segmentation style={popink,line width=1.2pt,dashed},
  before upper={\stepcounter{popexr}\popboxhead{Exercice résolu \thepopexr}{poporange}{#1}},
  before lower={\poppill{Solution}{popink}{popcream}\par\vspace{6pt}}}
% Exercice (non résolu en place) — la correction est dans la partie Corrigés
\newtcolorbox{exercice}[1][]{popbase,colframe=popink,
  before upper={\stepcounter{popexo}\popboxhead{Exercice \thepopexo}{popink}{#1}}}
% Corrigé
\newtcolorbox{corrige}[1][]{popbase,colframe=popgreen,colback=popgreenL,
  before upper={\popboxhead{Corrigé}{popgreen}{#1}}}
% Objectifs / prérequis / bilan
\newtcolorbox{objectifs}{popbase,colframe=popink,colback=poporangeL,
  before upper={\popboxhead{Objectifs de la leçon}{popink}{}}}
\newtcolorbox{prerequis}{popbase,colframe=popink,colback=popblueL,
  before upper={\popboxhead{Prérequis}{popblue}{}}}
\newtcolorbox{bilan}{popbase,colframe=popink,colback=white,boxrule=2.8pt,
  before upper={\popboxhead{Fiche bilan}{poporange}{L'essentiel à connaître pour l'examen}}}
% Figure encadrée avec légende
\newtcolorbox{popfigure}[1][]{enhanced,breakable=false,colback=white,colframe=popline,boxrule=1.4pt,arc=8pt,
  left=10pt,right=10pt,top=10pt,bottom=8pt,before skip=10pt,after skip=12pt,halign=center,
  lower separated=false,halign lower=center,
  fontlower=\popsemi\fontsize{9}{11.5}\selectfont\color{popmuted},
  #1}
\newcommand{\legende}[1]{\par\vspace{6pt}{\popsemi\fontsize{9}{11.5}\selectfont\color{popmuted}#1}}

% ============================================================
% Signature OnBuch+
% ============================================================
\newcommand{\poplogoinline}[1]{{\poplogo\fontsize{#1}{#1}\selectfont\color{popink}OnBuch\color{poporange}+}}
\newcommand{\signatureonbuch}{%
  \begin{tcolorbox}[enhanced,colback=popink,colframe=popink,boxrule=2.2pt,arc=10pt,
      drop shadow=poporange,left=14pt,right=14pt,top=10pt,bottom=10pt,before skip=0pt,after skip=0pt]
    \tikz[baseline=-0.7ex]\node[circle,fill=popgreen,draw=popcream,line width=1.3pt,minimum size=22pt,inner sep=0]{\popcheck{white}};%
    \hspace{9pt}\begin{minipage}[c]{0.62\linewidth}
      {\popxbold\fontsize{12}{13}\selectfont\color{popcream}Signé\ \textbullet\ {\poplogo OnBuch\color{poporange}+}}\par\vspace{2pt}
      {\raggedright\popsemi\fontsize{8}{10.5}\selectfont\color{popline}\MakeUppercase{Équipe pédagogique OnBuch+}\\\MakeUppercase{Conforme au programme officiel MINESEC}\par}
    \end{minipage}\hfill
    {\poplogo\fontsize{22}{22}\selectfont\color{popcream}OnBuch\color{poporange}+}
  \end{tcolorbox}%
}

% ============================================================
% Page de garde
% #1 titre · #2 matière · #3 niveau · #4 module · #5 n° leçon · #6 durée ·
% #7 description / objectifs (texte ou itemize)
% ============================================================
\newcommand{\pagegarde}[7]{%
  \thispagestyle{empty}
  \newgeometry{a4paper,margin=1.7cm,top=1.6cm,bottom=1.4cm}%
  \begin{tikzpicture}[remember picture,overlay]
    \fill[poporange!18] ([xshift=-3.2cm,yshift=-3.6cm]current page.north east) circle (4.6cm);
    \fill[poppurple!14] ([xshift=2.4cm,yshift=7.6cm]current page.south west) circle (3.8cm);
  \end{tikzpicture}%
  {\poplogo\fontsize{30}{30}\selectfont\color{popink}OnBuch\color{poporange}+}\hfill
  \raisebox{0.6ex}{\poppill{Cours signé \textbullet\ Premium}{popink}{popcream}}\par
  \vspace{1pt}\popsquiggle{poppurple}\par
  \vspace{8pt}
  \begin{tcolorbox}[enhanced,colback=poporange,colframe=popink,boxrule=2.8pt,arc=13pt,
      drop shadow=popink,left=18pt,right=18pt,top=12pt,bottom=13pt,after skip=10pt]
    \poppill{#2 \textbullet\ #3}{popink}{popcream}\hspace{5pt}\poppill{#4}{white}{popink}\par\vspace{8pt}
    {\popdisplay\fontsize{14}{14}\selectfont\color{popink}LEÇON #5}\par\vspace{4pt}
    {\raggedright\hyphenpenalty=10000\exhyphenpenalty=10000\popdisplay\fontsize{25}{27}\selectfont\color{popcream}\MakeUppercase{#1}\par}
    \vspace{8pt}
    {\popxbold\fontsize{9.5}{9.5}\selectfont\color{popink}\MakeUppercase{Durée conseillée : #6}}
  \end{tcolorbox}
  \begin{tcolorbox}[enhanced,colback=white,colframe=popink,boxrule=2.2pt,arc=10pt,
      drop shadow=popink,left=16pt,right=16pt,top=10pt,bottom=10pt,after skip=8pt]
    \poppill{Dans cette leçon}{poppurple}{white}\par\vspace{5pt}
    \popsemi\fontsize{9.3}{12.6}\selectfont\color{popink}#7
  \end{tcolorbox}
  \noindent\begin{minipage}[t]{0.487\linewidth}
    \begin{tcolorbox}[enhanced,colback=popgreen,colframe=popink,boxrule=2.2pt,arc=10pt,
        drop shadow=popink,left=11pt,right=11pt,top=10pt,bottom=10pt,height=3.1cm,valign=center]
      \tcbox[colback=white,colframe=popink,boxrule=1.3pt,arc=5pt,boxsep=0pt,nobeforeafter,
        left=3pt,right=3pt,top=3pt,bottom=3pt,valign=center]{\qrcode[height=1.9cm]{https://play.google.com/store/apps/details?id=com.luvvix.onbuch&hl=fr}}%
      \hspace{8pt}\begin{minipage}[c]{0.5\linewidth}
        {\popdisplay\fontsize{11}{12.5}\selectfont\color{white}\MakeUppercase{Télécharge OnBuch}}\par\vspace{4pt}
        {\raggedright\popsemi\fontsize{8.4}{11}\selectfont\color{white}Gratuit \textbullet\ Google Play\\Tuteur IA, annales, concours, résultats\par}
      \end{minipage}
    \end{tcolorbox}
  \end{minipage}\hfill
  \begin{minipage}[t]{0.487\linewidth}
    \begin{tcolorbox}[enhanced,colback=poppurple,colframe=popink,boxrule=2.2pt,arc=10pt,
        drop shadow=popink,left=11pt,right=11pt,top=10pt,bottom=10pt,height=3.1cm,valign=center]
      \tikz[baseline=-0.7ex]\node[circle,fill=white,draw=popink,line width=1.3pt,minimum size=1.6cm,inner sep=0]{\popdisplay\fontsize{24}{24}\selectfont\color{poppurple}+};%
      \hspace{8pt}\begin{minipage}[c]{0.6\linewidth}
        {\popdisplay\fontsize{11}{12.5}\selectfont\color{white}\MakeUppercase{Passe à OnBuch+}}\par\vspace{4pt}
        {\raggedright\popsemi\fontsize{8.4}{11}\selectfont\color{white}Tous les cours signés, Tuteur IA illimité, fiches PDF — onglet OnBuch+ de l'app\par}
      \end{minipage}
    \end{tcolorbox}
  \end{minipage}
  \vspace{8pt}
  \popcontenu
  \vfill
  \signatureonbuch
  \restoregeometry
  \newpage
}

% Rangée « Ce cours contient » (page de garde)
\newcommand{\poptile}[3]{% #1 couleur #2 gros texte #3 légende
  \begin{tcolorbox}[enhanced,colback=#1,colframe=popink,boxrule=2pt,arc=9pt,shadow={2.5pt}{-2.5pt}{0pt}{popink},
      left=6pt,right=6pt,top=9pt,bottom=9pt,halign=center,valign=center,height=1.75cm,width=\linewidth,nobeforeafter]
    {\popdisplay\fontsize{12.5}{13}\selectfont\color{white}\MakeUppercase{#2}}\par\vspace{3pt}
    {\centering\popsemi\fontsize{7.8}{9.5}\selectfont\color{white}#3\par}
  \end{tcolorbox}}
\newcommand{\popcontenu}{%
  \par\noindent{\popxboldls\fontsize{7.5}{7.5}\selectfont\color{popmuted}CE COURS SIGNÉ CONTIENT}\par\vspace{7pt}
  \noindent\begin{minipage}[t]{0.235\linewidth}\poptile{popblue}{Cours}{complet, illustré, pas à pas}\end{minipage}\hfill
  \begin{minipage}[t]{0.235\linewidth}\poptile{poporange}{Méthodes}{et exercices résolus}\end{minipage}\hfill
  \begin{minipage}[t]{0.235\linewidth}\poptile{poppink}{Exercices}{type examen + corrigés}\end{minipage}\hfill
  \begin{minipage}[t]{0.235\linewidth}\poptile{popgreen}{Fiche bilan}{l'essentiel à retenir}\end{minipage}\par}

% Sommaire
\newtcolorbox{sommaire}{enhanced,colback=white,colframe=popink,boxrule=2.2pt,arc=10pt,
  drop shadow=popink,left=18pt,right=18pt,top=13pt,bottom=13pt,before skip=6pt,after skip=20pt,
  before upper={\poppill{Sommaire}{poppurple}{white}\par\vspace{9pt}\popsemi\fontsize{10.6}{14.5}\selectfont\color{popink}}}

% Signature de fin de cours — poussée tout en bas de la dernière page
\newcommand{\findecours}{%
  \par\vfill\nopagebreak
  \begin{center}\popsquiggle{poporange}\end{center}\vspace{4pt}
  \signatureonbuch
}
\AtBeginDocument{\def\llbracket{\mathopen{[\mkern-3.5mu[}}\def\rrbracket{\mathclose{]\mkern-3.5mu]}}} % intervalles d'entiers (glyphe ⟦ absent de la police)
% Glyphes absents de Fira Math : redéfinis à partir de glyphes disponibles
\AtBeginDocument{%
  \def\setminus{\mathbin{\backslash}}\let\smallsetminus\setminus
  \def\vdots{\mathord{\vbox{\baselineskip3.2pt\lineskiplimit0pt\kern4pt\hbox{.}\hbox{.}\hbox{.}}}}%
  \def\ddots{\mathinner{\mkern1mu\raise6.5pt\hbox{.}\mkern2mu\raise3.6pt\hbox{.}\mkern2mu\raise0.7pt\hbox{.}\mkern1mu}}%
  \def\triangle{\mathord{\scalebox{0.9}{$\symup{\Delta}$}}}%
  \def\nabla{\mathord{\raisebox{\depth}{\scalebox{1}[-1]{$\symup{\Delta}$}}}}%
  \def\checkmark{\popcheck{popdark}}%
  \def\star{\mathbin{\ast}}\def\bigstar{\mathord{\ast}}%
  \def\diamond{\mathbin{\scalebox{0.6}{$\Diamond$}}}%
  \def\bigcup{\mathop{\mathchoice{\raisebox{-0.3em}{\scalebox{1.7}{$\cup$}}}{\scalebox{1.25}{$\cup$}}{\cup}{\cup}}\displaylimits}%
  \def\bigcap{\mathop{\mathchoice{\raisebox{-0.3em}{\scalebox{1.7}{$\cap$}}}{\scalebox{1.25}{$\cap$}}{\cap}{\cap}}\displaylimits}%
  \def\lesseqqgtr{\mathrel{\lessgtr}}\def\bigoplus{\mathop{\oplus}\displaylimits}%
}
\providecommand{\ssi}{si et seulement si} % abréviation fréquente
\AtBeginDocument{\providecommand{\wideparen}{\overparen}} % arc de cercle

%% FIGFIX-BEGIN v3 — correctifs globaux des figures (docs/onbuch-generation/tools/figfix)
\usepackage{adjustbox}\usepackage{etoolbox}
% toute figure TikZ/pgfplots plus large que la ligne est réduite à la largeur disponible
\BeforeBeginEnvironment{tikzpicture}{\begin{adjustbox}{max width=\linewidth}}
\AfterEndEnvironment{tikzpicture}{\end{adjustbox}}
% légendes pgfplots lisibles par-dessus les courbes
\pgfplotsset{every axis/.append style={legend style={fill=white,fill opacity=0.92,text opacity=1,draw=black!15,inner sep=3pt}}}
% étiquettes d'arêtes (syntaxe edge["..."]) avec fond blanc
\tikzset{every edge quotes/.append style={fill=white,inner sep=1.2pt}}
% bulles de cartes mentales : texte à la ligne dans la bulle, bulle un peu plus grande
\tikzset{every concept/.append style={text width=2.0cm,align=center,minimum size=2.3cm,inner sep=2pt}}
%% FIGFIX-END
\begin{document}

\pagegarde{Lesson 56 — Grammar: The active and passive voices}{Anglais}{Tle A}{Chapter 5 : Media and Communication}{EN56}{2 h}{Cette leçon de grammaire étudie la voix active et la voix passive en anglais : formation, emplois, transformation d'une voix à l'autre et pièges typiques des francophones. Elle s'appuie sur des exemples tirés des médias et de la communication, thème du chapitre 5.
\vspace{4pt}
\begin{multicols}{2}\small
\begin{itemize}
\item À la fin de cette leçon, je sais distinguer une phrase à la voix active d'une phrase à la voix passive.
\item À la fin de cette leçon, je sais former la voix passive à tous les temps principaux (présent, prétérit, futur, present perfect, modal).
\item À la fin de cette leçon, je sais transformer une phrase active en phrase passive et inversement.
\item À la fin de cette leçon, je sais utiliser l'agent introduit par 'by' et savoir quand l'omettre.
\item À la fin de cette leçon, je sais employer la voix passive dans les contextes appropriés (presse, science, procédures, faits objectifs).
\item À la fin de cette leçon, je sais éviter les erreurs fréquentes des francophones (oubli de 'be', mauvais participe, traduction mot à mot de 'on').
\end{itemize}\end{multicols}}

\begin{sommaire}
\begin{multicols}{2}
\begin{enumerate}
\item Observation et découverte
\item Formation de la voix passive
\item Emplois de la voix passive
\item Comparaison français/anglais et pièges
\item Application guidée et production
\item Activité d'intégration
\item Méthodes et exercices résolus
\item Exercices
\item Corrigés des exercices
\item Fiche bilan
\end{enumerate}
\end{multicols}
\end{sommaire}

\begin{objectifs}
\begin{itemize}
\item À la fin de cette leçon, je sais distinguer une phrase à la voix active d'une phrase à la voix passive.
\item À la fin de cette leçon, je sais former la voix passive à tous les temps principaux (présent, prétérit, futur, present perfect, modal).
\item À la fin de cette leçon, je sais transformer une phrase active en phrase passive et inversement.
\item À la fin de cette leçon, je sais utiliser l'agent introduit par 'by' et savoir quand l'omettre.
\item À la fin de cette leçon, je sais employer la voix passive dans les contextes appropriés (presse, science, procédures, faits objectifs).
\item À la fin de cette leçon, je sais éviter les erreurs fréquentes des francophones (oubli de 'be', mauvais participe, traduction mot à mot de 'on').
\item À la fin de cette leçon, je sais utiliser la voix passive dans des phrases et un court paragraphe de type article de presse.
\item À la fin de cette leçon, je sais repérer les verbes qui ne se mettent pas au passif (verbes intransitifs).
\end{itemize}
\end{objectifs}

\begin{prerequis}
\begin{itemize}
\item Les auxiliaires 'be' et 'have' et leurs conjugaisons (Première)
\item Le participe passé des verbes réguliers et irréguliers (Seconde/Première)
\item Les temps principaux : présent simple, prétérit, futur, present perfect (Première)
\item La structure de la phrase anglaise : sujet + verbe + complément (Seconde)
\item Les modaux can, must, should, will (Première)
\end{itemize}
\end{prerequis}

\newpage

\coursec{Observation et découverte}

Chaque matin à Yaoundé, dans les salles de classe comme dans les salons, on entend les nouvelles à la CRTV : \eng{A new road has been built in Bamenda}, \eng{The President was interviewed yesterday}, \eng{Two journalists were arrested in Lagos}. Au marché de Buea, un commerçant anglophone annonce fièrement : \eng{These tomatoes were grown in Santa} — là où un francophone dirait naturellement « On a cultivé ces tomates à Santa ». Partout, la même structure étrange revient : un sujet, une forme du verbe \eng{be}, puis un participe passé. Pourquoi les journalistes et les locuteurs anglophones l'emploient-ils si souvent ? Que signifie-t-elle exactement, et comment la construire sans faute ?

\textbf{Problématique :} comment l'anglais exprime-t-il qu'un sujet \emph{subit} une action au lieu de la faire, et en quoi cette structure — la voix passive — diffère-t-elle du français ? La maîtriser, c'est comprendre les médias anglophones (BBC, CNN, The Guardian) et réussir les épreuves de grammaire du Baccalauréat A.

\courssub{Lire des titres de presse : première observation}

\begin{textebox}[Read these headlines]
\textbf{A new radio station was opened in Bamenda.}

\textbf{The President was interviewed on CRTV last night.}

\textbf{English is spoken in many countries.}

\textbf{Two journalists were arrested in Lagos.}

\textbf{A new TV channel has been launched in Douala.}

\textbf{The match will be broadcast live.}

What do you notice about the verbs?
\end{textebox}

\begin{methode}[Observer des exemples de grammaire : la bonne démarche]
\begin{enumerate}
\item \textbf{Repérer le sujet grammatical} de chaque phrase (le mot ou le groupe avant le verbe).
\item \textbf{Se demander qui fait l'action} : est-ce le sujet, ou quelqu'un d'autre ?
\item \textbf{Entourer le verbe} et décomposer ses parties (auxiliaire + verbe principal).
\item \textbf{Comparer} avec une phrase « normale » qui exprime la même idée.
\item \textbf{Dégager la structure récurrente} et formuler une première règle.
\end{enumerate}
\end{methode}

Appliquons cette démarche au premier titre : \eng{A new radio station was opened in Bamenda.} Le sujet grammatical est \eng{a new radio station} (« une nouvelle station de radio »). Question guide : est-ce la station qui \emph{ouvre} ? Non ! Une station ne peut pas s'ouvrir toute seule : quelqu'un — une entreprise, le gouvernement, un propriétaire — l'a ouverte. Le sujet \emph{reçoit} donc l'action ; il ne la fait pas.

Même constat pour les autres titres :

\begin{center}
\begin{tabularx}{\linewidth}{|X|X|X|}
\hline
\rowcolor{popblueL} \textbf{Sujet grammatical} & \textbf{Verbe} & \textbf{Le sujet fait-il l'action ?} \\
\hline
A new radio station & was opened & Non : la station \emph{est ouverte} par quelqu'un \\
\hline
The President & was interviewed & Non : le Président \emph{est interviewé} par un journaliste \\
\hline
English & is spoken & Non : l'anglais \emph{est parlé} par des millions de personnes \\
\hline
Two journalists & were arrested & Non : les journalistes \emph{sont arrêtés} par la police \\
\hline
The match & will be broadcast & Non : le match \emph{est diffusé} par une chaîne de télévision \\
\hline
\end{tabularx}
\end{center}

\begin{aretenir}[Le constat fondamental]
Dans toutes ces phrases, le sujet grammatical n'est \textbf{pas} l'auteur de l'action : il en est la \emph{cible}, le \emph{résultat} ou la \emph{victime}. On dit qu'il \textbf{subit} l'action. C'est exactement le contraire d'une phrase ordinaire comme \eng{The police arrested two journalists}, où le sujet (\eng{the police}) fait l'action.
\end{aretenir}

\courssub{Actif ou passif : qui fait l'action ?}

Comparons maintenant chaque titre avec sa version « normale », dite \textbf{active} :

\begin{exemplebox}[Active voice vs passive voice]
\eng{The police arrested two journalists.} « La police a arrêté deux journalistes. » — \textbf{actif} : le sujet (\eng{the police}) fait l'action.

\eng{Two journalists were arrested (by the police).} « Deux journalistes ont été arrêtés (par la police). » — \textbf{passif} : le sujet (\eng{two journalists}) subit l'action.

\eng{A company has launched a new TV channel.} « Une entreprise a lancé une nouvelle chaîne de télévision. » — \textbf{actif}.

\eng{A new TV channel has been launched (by a company).} « Une nouvelle chaîne de télévision a été lancée. » — \textbf{passif}.

\eng{The journalist wrote the article.} « Le journaliste a écrit l'article. » — \textbf{actif}.

\eng{The article was written by the journalist.} « L'article a été écrit par le journaliste. » — \textbf{passif}.
\end{exemplebox}

\begin{popfigure}
\begin{tikzpicture}[font=\small]
% ACTIVE
\node[draw=popblue, fill=popblueL, rounded corners, align=center, minimum width=2.6cm] (s1) at (0,2.2) {\textbf{Sujet}\\ (agent)};
\node[draw=popdark, fill=popcream, rounded corners, align=center, minimum width=2.2cm] (v1) at (4.2,2.2) {\textbf{Verbe}\\ (action)};
\node[draw=popgreen, fill=popgreenL, rounded corners, align=center, minimum width=2.6cm] (o1) at (8.4,2.2) {\textbf{COD}\\ (patient)};
\draw[-{Stealth[length=3mm]}, very thick, popblue] (s1) -- (v1);
\draw[-{Stealth[length=3mm]}, very thick, popblue] (v1) -- (o1);
\node[font=\bfseries\small, popblue] at (4.2,3.1) {ACTIVE VOICE};
\node[font=\footnotesize\itshape, align=center] at (4.2,1.35) {The police arrested two journalists.};
% PASSIVE
\node[draw=popgreen, fill=popgreenL, rounded corners, align=center, minimum width=2.6cm] (s2) at (0,-0.6) {\textbf{Sujet}\\ (patient)};
\node[draw=popdark, fill=popcream, rounded corners, align=center, minimum width=2.2cm] (v2) at (4.2,-0.6) {\textbf{be + V-en}\\ (action)};
\node[draw=poporange, fill=poporangeL, rounded corners, align=center, minimum width=2.6cm] (o2) at (8.4,-0.6) {\textbf{by + agent}\\ (souvent omis)};
\draw[-{Stealth[length=3mm]}, very thick, poporange] (o2) -- (v2);
\draw[-{Stealth[length=3mm]}, very thick, poporange] (v2) -- (s2);
\node[font=\bfseries\small, poporange] at (4.2,0.3) {PASSIVE VOICE};
\node[font=\footnotesize\itshape, align=center] at (4.2,-1.45) {Two journalists were arrested (by the police).};
\end{tikzpicture}
\legende{À la voix active, l'action part du sujet vers le COD ; à la voix passive, le COD de la phrase active devient le sujet, et l'agent est introduit par \eng{by} (ou disparaît).}
\end{popfigure}

\begin{definition}[La voix passive — passive voice]
A sentence is in the \cle{passive voice} when the subject \emph{receives} the action instead of \emph{doing} it. Une phrase est à la \cle{voix passive} lorsque son sujet grammatical \emph{subit} l'action au lieu de l'accomplir. À la \cle{voix active} (\cle{active voice}), au contraire, le sujet fait l'action : c'est l'agent.
\end{definition}

Deux termes techniques à retenir dès maintenant :

\begin{itemize}
\item l'\cle{agent} : celui qui fait réellement l'action. À la voix passive, il est introduit par la préposition \eng{by} : \eng{The article was written \textbf{by the journalist}.} ;
\item le \cle{patient} : celui ou ce qui subit l'action. À la voix passive, c'est lui qui devient le sujet grammatical.
\end{itemize}

\courssub{Repérer la structure récurrente : be + participe passé}

Revenons à nos titres de presse et décomposons chaque verbe passif :

\begin{center}
\begin{tabularx}{\linewidth}{|X|X|X|}
\hline
\rowcolor{popblueL} \textbf{Titre} & \textbf{be (conjugué)} & \textbf{Participe passé} \\
\hline
English \textbf{is spoken} in many countries. & is & spoken \\
\hline
A new radio station \textbf{was opened} in Bamenda. & was & opened \\
\hline
Two journalists \textbf{were arrested} in Lagos. & were & arrested \\
\hline
A new TV channel \textbf{has been launched} in Douala. & has been & launched \\
\hline
The match \textbf{will be broadcast} live. & will be & broadcast \\
\hline
\end{tabularx}
\end{center}

\begin{propriete}[La signature de la voix passive]
Toute forme passive en anglais contient deux éléments obligatoires :
\begin{enumerate}
\item le verbe \cle{be}, conjugué au temps voulu et accordé avec le sujet (\eng{is}, \eng{was}, \eng{were}, \eng{has been}, \eng{will be}...) ;
\item le \cle{participe passé} du verbe principal (\eng{opened}, \eng{arrested}, \eng{spoken}, \eng{written}, \eng{broadcast}...).
\end{enumerate}
C'est \eng{be} qui porte le temps et l'accord ; le participe passé, lui, ne change jamais. La formation détaillée à tous les temps sera étudiée dans la section suivante ; à ce stade, sachez simplement \emph{reconnaître} cette signature : \textbf{be + participe passé}.
\end{propriete}

\begin{attention}[Un participe passé seul ne suffit pas]
L'erreur classique du francophone est d'oublier l'auxiliaire \eng{be} : \eng{The match broadcast live} ne veut rien dire. En français, « le match diffusé en direct » peut fonctionner comme un titre de rubrique ; en anglais, il faut toujours la forme complète : \eng{The match \textbf{will be} broadcast live.} Pas de \eng{be}, pas de passif !
\end{attention}

\begin{popfigure}
\begin{tikzpicture}[mindmap, grow cyclic, every node/.style={concept, align=center, font=\small},
concept color=popblue!40, root concept/.append style={font=\bfseries\small, minimum size=2.6cm},
level 1/.append style={level distance=3.4cm, sibling angle=90, minimum size=2.2cm}]
\node[root concept, align=center]{PASSIVE\\ VOICE}
  child[concept color=popgreen!45] { node[align=center]{Forme\\ be + participe\\ passé} }
  child[concept color=poporange!45] { node[align=center]{Emplois\\ agent inconnu\\ ou événement\\ mis en avant} }
  child[concept color=poppink!45] { node[align=center]{Agent\\ by + agent\\ (souvent omis)} }
  child[concept color=poppurple!40] { node[align=center]{Temps\\ c'est be qui\\ se conjugue} };
\end{tikzpicture}
\legende{Carte mentale de la voix passive : une forme (be + participe passé), des emplois précis, un agent introduit par \eng{by}, et un auxiliaire \eng{be} qui porte tous les temps.}
\end{popfigure}

\courssub{Première comparaison avec le français}

Bonne nouvelle : le français possède aussi une voix passive, construite de la même façon — « être » + participe passé. \eng{Two journalists were arrested} se traduit mot à mot : « Deux journalistes \textbf{ont été arrêtés} ». La structure est donc familière.

Mais il y a une différence d'\emph{emploi} essentielle : l'anglais utilise la voix passive \textbf{là où le français préfère souvent le pronom « on »} ou une tournure active. C'est la grande découverte de cette leçon.

\begin{center}
\begin{tabularx}{\linewidth}{|X|X|X|}
\hline
\rowcolor{popblueL} \textbf{Français (naturel)} & \textbf{English (natural)} & \textbf{Remarque} \\
\hline
On parle anglais dans de nombreux pays. & English is spoken in many countries. & « on » français = passif anglais \\
\hline
On a cultivé ces tomates à Santa. & These tomatoes were grown in Santa. & passif préférable \\
\hline
On a ouvert une nouvelle station de radio à Bamenda. & A new radio station was opened in Bamenda. & passif préférable \\
\hline
On a volé ma moto hier ! & My motorbike was stolen yesterday! & le voleur est inconnu \\
\hline
On fabrique ces chaussures au Cameroun. & These shoes are made in Cameroon. & \eng{made in} : expression passive figée \\
\hline
On réparera la route l'année prochaine. & The road will be repaired next year. & futur passif : \eng{will be + V-en} \\
\hline
\end{tabularx}
\end{center}

\begin{attention}[Piège majeur : ne pas traduire « on » par un sujet imaginaire]
Face à « On a ouvert une nouvelle station », beaucoup d'élèves francophones cherchent un sujet qui n'existe pas (\eng{they opened}, \eng{people opened}...) alors que la voix passive est exactement faite pour cela : \eng{A new station was opened.} Règle pratique : quand le français dit « on » sans préciser qui agit, l'anglais préfère presque toujours le passif. Inversement, ne traduisez jamais un passif anglais par « être » + infinitif : \eng{The match will be broadcast} ne signifie pas « le match sera diffuser » mais « le match sera \emph{diffusé} ».
\end{attention}

\begin{savaistu}[Pourquoi la presse anglophone adore le passif]
Ouvrez le site de la BBC, de CNN ou du quotidien britannique \emph{The Guardian} : les titres et les dépêches regorgent de voix passives — \eng{Three people were injured in the accident}, \eng{A new law has been passed}, \eng{The suspect is being questioned}. Pourquoi ? Parce que la voix passive met l'accent sur \textbf{l'événement} et ses victimes plutôt que sur l'auteur de l'action, souvent inconnu, évident ou sans importance au moment où la nouvelle « tombe ». Elle permet aussi au journaliste de rester neutre et factuel. Au Cameroun, les journaux anglophones comme \emph{The Guardian Post} ou les bulletins de la CRTV en anglais suivent exactement la même tradition. Lire la presse anglophone sans connaître le passif, c'est se priver de la moitié des informations !
\end{savaistu}

\courssub{Application guidée : je repère, je compare, je conclus}

\begin{exoresolu}[Repérer la voix et l'agent]
\textbf{Énoncé.} Pour chaque phrase, indiquez si elle est à la voix active ou passive, soulignez le sujet grammatical et, s'il y a lieu, encadrez l'agent.

1. \eng{The students wrote a letter to the headmaster.}

2. \eng{A letter was written to the headmaster.}

3. \eng{The match will be broadcast live on CRTV.}

4. \eng{Farmers grow cocoa in the South-West region.}

5. \eng{Cocoa is grown in the South-West region.}

\tcblower
\textbf{Solution détaillée.}

1. \eng{The students wrote a letter to the headmaster.} — \textbf{Voix active}. Sujet : \eng{the students}. Le verbe \eng{wrote} est un prétérit simple, sans \eng{be} : les élèves font l'action d'écrire.

2. \eng{A letter was written to the headmaster.} — \textbf{Voix passive}. Sujet : \eng{a letter}. On reconnaît la signature \eng{was} (be au prétérit) + \eng{written} (participe passé de \eng{write}). La lettre subit l'action ; l'agent n'est pas mentionné.

3. \eng{The match will be broadcast live on CRTV.} — \textbf{Voix passive}. Sujet : \eng{the match}. Signature : \eng{will be} + \eng{broadcast} (participe passé identique à la base verbale). L'agent (la chaîne) est sous-entendu.

4. \eng{Farmers grow cocoa in the South-West region.} — \textbf{Voix active}. Sujet : \eng{farmers}. Présent simple, pas de \eng{be} : les agriculteurs font l'action.

5. \eng{Cocoa is grown in the South-West region.} — \textbf{Voix passive}. Sujet : \eng{cocoa}. Signature : \eng{is} + \eng{grown} (participe passé de \eng{grow}). Remarquez que cette phrase passive est la traduction naturelle de « On cultive le cacao dans la région du Sud-Ouest ».
\end{exoresolu}

\begin{exercice}[À vous de jouer : active ou passive ?]
Indiquez la voix de chaque phrase, puis traduisez-la en français naturel (attention au « on » !).

1. \eng{A new bridge has been built over the Wouri river.}

2. \eng{The government built a new bridge over the Wouri river.}

3. \eng{French and English are taught in Cameroonian schools.}

4. \eng{The referee blew the final whistle.}

5. \eng{The final whistle was blown by the referee.}
\end{exercice}

\begin{corrige}[Exercice : active ou passive ?]
1. \textbf{Passive} (\eng{has been} + \eng{built}). « On a construit un nouveau pont sur le Wouri » ou « Un nouveau pont a été construit sur le Wouri. »

2. \textbf{Active} : le sujet \eng{the government} fait l'action. « Le gouvernement a construit un nouveau pont sur le Wouri. »

3. \textbf{Passive} (\eng{are} + \eng{taught}). « On enseigne le français et l'anglais dans les écoles camerounaises. »

4. \textbf{Active} : \eng{the referee} fait l'action. « L'arbitre a sifflé le coup de sifflet final. »

5. \textbf{Passive} avec agent exprimé (\eng{by the referee}). « Le coup de sifflet final a été donné par l'arbitre. »
\end{corrige}

\begin{aretenir}[Ce qu'il faut retenir de cette découverte]
\begin{itemize}
\item À la \textbf{voix active}, le sujet \emph{fait} l'action ; à la \textbf{voix passive}, le sujet \emph{subit} l'action.
\item La signature du passif est : \textbf{be (conjugué) + participe passé}. Sans \eng{be}, pas de passif.
\item L'agent, quand il est mentionné, est introduit par \eng{by} ; il est très souvent omis.
\item Le passif anglais correspond souvent au « on » français : \eng{English is spoken here} = « On parle anglais ici ».
\item La presse anglophone l'utilise massivement pour mettre l'événement en avant.
\end{itemize}
\end{aretenir}

Maintenant que nous savons \emph{reconnaître} la voix passive et comprendre sa logique, la prochaine étape consiste à apprendre à la \emph{construire} à tous les temps : c'est l'objet de la section suivante, « Formation de la voix passive ».

\coursec{Formation de la voix passive}

Après avoir observé dans la section précédente comment la voix passive permet de déplacer le focus de l'agent vers l'action ou l'objet concerné, nous allons maintenant maîtriser sa construction mécanique. La règle est d'une régularité exemplaire : une fois le principe compris, vous pourrez mettre n'importe quel temps à la voix passive sans hésitation.

\courssub{La règle fondamentale : la formule magique}

\begin{propriete}[Règle de formation de la voix passive]
Pour transformer une phrase active en phrase passive, on applique la formule immuable :
\medskip

\centerline{\textbf{Sujet (ancien COD) + BE (conjugué au temps du verbe actif) + Participe Passé (+ by + Agent)}}
\medskip

\noindent Points clés à retenir :
\begin{itemize}
    \item L'auxiliaire \cle{BE} est le \textbf{seul} verbe qui porte la marque du temps (présent, passé, futur, modalité) et de la personne (singulier/pluriel).
    \item Le \cle{participe passé} (past participle) \textbf{ne change jamais de forme} : il reste invariable, quel que soit le temps ou le sujet.
    \item L'agent (l'ancien sujet) est introduit par la préposition \cle{by} et n'est mentionné que s'il apporte une information nécessaire (s'il est inconnu, évident ou sans importance, on l'omet).
\end{itemize}
\end{propriete}

\courssub{Observation sur exemples authentiques}

Avant de dresser le tableau complet des temps, observons la mécanique à l'œuvre sur des phrases tirées de l'actualité médiatique (thème du Chapter 5).

\begin{textebox}[Scripts d'actualité — Voix active vs Voix passive]
\textbf{Active:} \eng{The national radio \textbf{broadcasts} the news at 8 p.m. every evening.} (La radio nationale diffuse les informations à 20 h chaque soir.)

\textbf{Passive:} \eng{The news \textbf{is broadcast} at 8 p.m. every evening (by the national radio).} (Les informations sont diffusées à 20 h chaque soir (par la radio nationale).)
\medskip

\textbf{Active:} \eng{A journalist \textbf{interviewed} the Minister of Communication yesterday.} (Un journaliste a interviewé le Ministre de la Communication hier.)

\textbf{Passive:} \eng{The Minister of Communication \textbf{was interviewed} yesterday (by a journalist).} (Le Ministre de la Communication a été interviewé hier (par un journaliste).)
\medskip

\textbf{Active:} \eng{The government \textbf{will launch} a new digital platform next month.} (Le gouvernement lancera une nouvelle plateforme numérique le mois prochain.)

\textbf{Passive:} \eng{A new digital platform \textbf{will be launched} next month (by the government).} (Une nouvelle plateforme numérique sera lancée le mois prochain (par le gouvernement).)
\end{textebox}

\begin{exemplebox}[Analyse pas à pas : \eng{The news is broadcast at 8 p.m.}]
\begin{enumerate}
    \item \textbf{Identification du COD actif} : \eng{the news} (receveur de l'action \eng{broadcast}).
    \item \textbf{Nouveau sujet passif} : \eng{The news} (remonte en tête de phrase). Accord : singulier $\rightarrow$ \eng{is}.
    \item \textbf{Auxiliaire BE au temps de l'actif} : L'actif est au \textbf{present simple} (\eng{broadcasts}) $\rightarrow$ \eng{BE} = \eng{is} (3\textsuperscript{e} pers. sing.).
    \item \textbf{Participe passé} : \eng{broadcast} (verbe irrégulier : broadcast / broadcast / broadcast). Forme invariée.
    \item \textbf{Complément d'agent (optionnel)} : \eng{by the national radio}.
\end{enumerate}
\eng{The news is broadcast at 8 p.m.} $\rightarrow$ « Les informations sont diffusées à 20 h. »
\end{exemplebox}

\courssub{Tableau récapitulatif des formes (Le « Moteur » BE)}

Le tableau ci-dessous synthétise les formes de \cle{BE} aux principaux temps. C'est la « colonne vertébrale » de la voix passive. Apprenez-le par cœur.

\begin{popfigure}
\begin{center}
\begin{tabularx}{\linewidth}{|X|X|X|}
\hline
\rowcolor{popblue} \textbf{Temps / Modalité} & \textbf{Voix Active (to write)} & \textbf{Voix Passive (to be written)} \\
\hline
Present Simple & writes & \textbf{is written} \\
\hline
Present Continuous & is writing & \textbf{is being written} \\
\hline
Past Simple & wrote & \textbf{was written} \\
\hline
Past Continuous & was writing & \textbf{was being written} \\
\hline
Present Perfect & has written & \textbf{has been written} \\
\hline
Past Perfect & had written & \textbf{had been written} \\
\hline
Future (will) & will write & \textbf{will be written} \\
\hline
Future Perfect & will have written & \textbf{will have been written} \\
\hline
Modal (can/must) & can write & \textbf{can be written} \\
\hline
Modal Perfect & must have written & \textbf{must have been written} \\
\hline
\end{tabularx}
\end{center}
\legende{Tableau des formes de la voix passive. Notez la régularité : BE prend toutes les marques, V3 (participe passé) reste fixe.}
\end{popfigure}

\begin{aretenir}[Le point crucial : \eng{BE} + \eng{BEEN} aux temps composés]
Aux temps parfaits (Present Perfect, Past Perfect, Future Perfect, Modal Perfect), l'auxiliaire \cle{BE} prend la forme \cle{BEEN} (participe passé de BE) suivi du participe passé du verbe lexical.
\medskip
\centerline{Active: \eng{has written} $\rightarrow$ Passive: \eng{has \textbf{been} written}}
\centerline{Active: \eng{will have written} $\rightarrow$ Passive: \eng{will have \textbf{been} written}}
\medskip
Ne jamais oublier le \eng{been} ! C'est l'erreur n°1 aux temps composés.
\end{aretenir}

\courssub{Les participes passés irréguliers : la liste de survie}

Comme le participe passé ne change jamais, vous devez connaître par cœur les formes irrégulières des verbes fréquents dans les médias et la communication.

\begin{center}
\begin{tabularx}{\linewidth}{|l|l|l|X|}
\hline
\rowcolor{popblue} \textbf{Infinitif} & \textbf{Prétérit} & \textbf{Participe Passé (V3)} & \textbf{Traduction / Contexte média} \\
\hline
write & wrote & \textbf{written} & écrit (article, report) \\
take & took & \textbf{taken} & pris (photo, interview, note) \\
make & made & \textbf{made} & fait (film, announcement, decision) \\
build & built & \textbf{built} & construit (studio, network) \\
speak & spoke & \textbf{spoken} & parlé (langue, interview) \\
steal & stole & \textbf{stolen} & volé (data, equipment) \\
break & broke & \textbf{broken} & cassé (news, record, silence) \\
choose & chose & \textbf{chosen} & choisi (candidate, topic) \\
send & sent & \textbf{sent} & envoyé (email, reporter) \\
broadcast & broadcast & \textbf{broadcast} & diffusé (news, programme) \\
\hline
\end{tabularx}
\end{center}

\begin{attention}[Piège majeur : \eng{was wrote} vs \eng{was written}]
\emph{Erreur typique de francophone :} On confond le prétérit (\eng{wrote}) et le participe passé (\eng{written}) car en français, le participe passé « écrit » ressemble souvent au passé simple « il écrivit » ou au passé composé « il a écrit ».
\medskip
\noindent \textbf{FAUX} : \eng{The article *was wrote by John.} \\
\noindent \textbf{VRAI} : \eng{The article \textbf{was written} by John.}
\medskip
\noindent \textbf{Règle d'or} : Après \cle{BE} (was, were, is, are, been, be), il \textbf{n'y a jamais} de forme en \cle{-ed} (prétérit régulier) ni de forme de prétérit irrégulier (wrote, took, made). Il y a \textbf{uniquement le participe passé (V3)}.
\end{attention}

\courssub{Cas particulier : Les verbes à deux compléments (Ditransitifs)}

Certains verbes (give, send, offer, tell, show, promise, pay, teach...) prennent deux objets : un COI (à qui ?) et un COD (quoi ?). À la voix passive, \textbf{les deux} peuvent devenir sujet. Cela donne deux passifs possibles.

\begin{propriete}[Double passif : Two Passive Possibilities]
Active : \eng{Someone \textbf{gave} the reporter \textbf{a microphone}.} (Quelqu'un a donné un micro au reporter.)
\medskip
\begin{itemize}
    \item \textbf{Passif 1 (Focus sur la personne / COI devient sujet) — Le plus fréquent en anglais :}
    \eng{The reporter \textbf{was given} a microphone.} (Le reporter s'est vu donner un micro / Un micro a été donné au reporter.)
    \item \textbf{Passif 2 (Focus sur l'objet / COD devient sujet) :}
    \eng{A microphone \textbf{was given} \textbf{to} the reporter.} (Un micro a été donné au reporter.) $\rightarrow$ Notez l'ajout de \cle{to} devant le COI devenu complément.
\end{itemize}
\end{propriete}

\begin{exemplebox}[Autres exemples de double passif]
\begin{center}
\begin{tabularx}{\linewidth}{|X|X|}
\hline
\rowcolor{popblue} \textbf{Active} & \textbf{Passives possibles} \\
\hline
\eng{They \textbf{offered} him a job.} & 1. \eng{He \textbf{was offered} a job.} (Usuel) \\
& 2. \eng{A job \textbf{was offered} \textbf{to} him.} \\
\hline
\eng{She \textbf{sent} the editor the photos.} & 1. \eng{The editor \textbf{was sent} the photos.} \\
& 2. \eng{The photos \textbf{were sent} \textbf{to} the editor.} \\
\hline
\eng{The station \textbf{promised} the listeners a debate.} & 1. \eng{The listeners \textbf{were promised} a debate.} \\
& 2. \eng{A debate \textbf{was promised} \textbf{to} the listeners.} \\
\hline
\end{tabularx}
\end{center}
\end{exemplebox}

\begin{attention}[Oublier \eng{to} au passif 2]
\emph{Erreur fréquente :} \eng{*A microphone was given the reporter.}
\medskip
\noindent Quand le COD devient sujet, l'ancien COI (la personne) \textbf{doit} être introduit par \cle{to} (ou \cle{for} selon le verbe : \eng{bought for}, \eng{sent to}).
\end{attention}

\courssub{Exception absolue : Les verbes intransitifs n'ont pas de passif}

C'est une règle logique : la voix passive transforme le COD (Objet Direct) en Sujet. Si le verbe n'a \textbf{pas de COD} (verbe intransitif), la transformation est impossible.

\begin{definition}[Verbes intransitifs — Pas de voix passive]
Les verbes d'état, de mouvement, ou d'action sans objet direct ne peuvent pas se mettre à la voix passive.
\medskip
\noindent \textbf{Exemples courants :} \cle{arrive, happen, occur, die, sleep, go, come, live, stay, wait, exist, disappear, rise, fall, laugh, cry}.
\end{definition}

\begin{exemplebox}[Impossible de passiver]
\begin{itemize}
    \item \eng{The accident \textbf{happened} yesterday.} $\rightarrow$ \eng{*The accident *was happened yesterday.} (\textbf{FAUX})
    \item \eng{Many people \textbf{died} in the flood.} $\rightarrow$ \eng{*Many people *were died in the flood.} (\textbf{FAUX})
    \item \eng{The train \textbf{arrived} late.} $\rightarrow$ \eng{*The train *was arrived late.} (\textbf{FAUX})
    \item \eng{The journalist \textbf{lives} in Yaoundé.} $\rightarrow$ \eng{*The journalist *is lived in Yaoundé.} (\textbf{FAUX})
\end{itemize}
\medskip
\noindent \textbf{Astuce de vérification} : Posez la question « \textbf{What/Whom} + verbe ? ». Si la question n'a pas de sens (\eng{*What happened the accident?}), le verbe est intransitif $\rightarrow$ Pas de passif.
\end{exemplebox}

\begin{savaistu}[Verbes « transitifs en français, intransitifs en anglais »]
Attention aux faux amis de transitivité !
\begin{itemize}
    \item \eng{\textbf{Listen to}} (écouter) : Transitif indirect en anglais. \eng{He listened to the radio.} $\rightarrow$ Passif possible : \eng{The radio \textbf{was listened to}.} (Verbe à particule/préposition).
    \item \eng{\textbf{Look at}} (regarder) : \eng{The picture \textbf{was looked at}.}
    \item \eng{\textbf{Phone / Call}} (téléphoner à) : Transitif direct en anglais ! \eng{He phoned me.} $\rightarrow$ \eng{I \textbf{was phoned}.}
    \item \eng{\textbf{Enter}} (entrer dans) : Transitif direct en anglais. \eng{He entered the room.} $\rightarrow$ \eng{The room \textbf{was entered}.} (Pas de \eng{into} au passif).
\end{itemize}
\end{savaistu}

\courssub{Schéma de transformation visuel}

Pour ancrer la mécanique, visualisons la « chirurgie » de la phrase.

\begin{popfigure}
\begin{tikzpicture}[
    node distance=8mm and 12mm,
    box/.style={draw=popdark, thick, rounded corners=4pt, fill=popcream, minimum width=2.5cm, minimum height=1.2cm, align=center, font=\small},
    boxact/.style={box, fill=popgreenL!60},
    boxpass/.style={box, fill=poporangeL!60},
    arrow/.style={-Stealth, thick, popblue, shorten >=2pt, shorten <=2pt},
    labelstyle/.style={font=\tiny\bfseries, text=popdark},
    cross/.style={draw=popink, thick, -Stealth, shorten >=2pt, shorten <=2pt}
]
% Ligne Active
\node[boxact, align=center] (s1){Sujet Agent\\ \eng{John}};
\node[boxact, right=of s1, align=center] (v1){Verbe Actif\\ \eng{eats} \\ \scriptsize(Present Simple)};
\node[boxact, right=of v1, align=center] (o1){COD\\ \eng{the cake}};

% Flèche transformation
\node[draw=poporange, thick, dashed, rounded corners, fit=(s1)(o1), label={[labelstyle]above:Transformation Passive}] (fit1) {};

% Ligne Passive (décalée vers le bas)
\node[boxpass, below=2.5cm of s1, align=center] (s2){Nouveau Sujet\\ \eng{The cake}};
\node[boxpass, right=of s2, align=center] (v2){BE + V3\\ \eng{is eaten} \\ \scriptsize(be au présent + V3)};
\node[boxpass, right=of v2, align=center] (a2){By-agent (optionnel)\\ \eng{by John}};

% Flèches de déplacement
\draw[arrow, bend left=30] (o1.south) to node[left, labelstyle, pos=0.5] {Devient Sujet} (s2.north);
\draw[arrow, bend right=30] (s1.south) to node[right, labelstyle, pos=0.5] {Devient By-agent} (a2.north);
\draw[cross] (v1.south) -- node[left, labelstyle] {Verbe $\rightarrow$ BE + V3} (v2.north);

% Légendes internes
\node[labelstyle, above=2mm of s1] {Active Voice};
\node[labelstyle, above=2mm of s2] {Passive Voice};
\end{tikzpicture}
\legende{Schéma de transformation : Le COD \eng{the cake} remonte en sujet ; le verbe \eng{eats} se décompose en \eng{is} (temps) + \eng{eaten} (V3) ; l'ancien sujet \eng{John} devient complément d'agent \eng{by John}.}
\end{popfigure}

\courssub{Méthode pas à pas : Transformer l'actif en passif (4 étapes)}

\begin{methode}[De l'actif au passif en 4 étapes]
Appliquez cette méthode systématiquement pour ne jamais vous tromper.
\begin{enumerate}
    \item \textbf{Repérez le COD (Complément d'Objet Direct)} de la phrase active. Question : « \textbf{What/Whom} + sujet + verbe ? ».
    \item \textbf{Placez ce COD en tête de phrase} : il devient le \textbf{Nouveau Sujet}. Faites l'accord (singulier/pluriel) pour choisir la forme de BE.
    \item \textbf{Conjuguez l'auxiliaire BE} au \textbf{même temps} que le verbe actif, puis ajoutez le \textbf{Participe Passé (V3)} du verbe lexical. (Vérifiez l'irrégularité !).
    \item \textbf{Décidez pour l'agent (By-agent)} : L'ancien sujet est-il important ? Si OUI $\rightarrow$ ajoutez \eng{by} + ancien sujet. Si NON (inconnu, évident, général) $\rightarrow$ \textbf{omettez-le}.
\end{enumerate}
\end{methode}

\begin{exoresolu}[Application de la méthode]
\textbf{Énoncé :} Transformez la phrase active suivante en voix passive. Omettez l'agent s'il n'est pas nécessaire.
\medskip
\eng{The editor \textbf{has rejected} the manuscript.}
\medskip
\textbf{Solution détaillée :}
\begin{enumerate}
    \item \textbf{COD} : \eng{What has the editor rejected?} $\rightarrow$ \eng{the manuscript}.
    \item \textbf{Nouveau Sujet} : \eng{The manuscript} (Singulier).
    \item \textbf{Temps du verbe actif} : \eng{has rejected} = \textbf{Present Perfect}.
    \item \textbf{BE au Present Perfect (3e pers. sing.)} : \eng{has been}.
    \item \textbf{Participe passé de reject} : \eng{rejected} (régulier).
    \item \textbf{Agent} : \eng{The editor} (spécifique, peut être gardé ou omis selon le contexte). Ici, on le garde pour l'exemple.
\end{enumerate}
\medskip
\noindent \textbf{Résultat :} \eng{The manuscript \textbf{has been rejected} (by the editor).}
\medskip
\noindent \textit{Traduction :} « Le manuscrit a été rejeté (par le rédacteur en chef). »
\end{exoresolu}

\begin{exercice}[Entraînement immédiat — Transformation]
Mettez les phrases suivantes à la voix passive. Conservez le temps. Omettez l'agent quand il est générique (\eng{people, someone, they}).
\begin{enumerate}
    \item \eng{Journalists \textbf{write} news articles every day.} (Present Simple)
    \item \eng{The technician \textbf{is repairing} the microphone right now.} (Present Continuous)
    \item \eng{The station \textbf{broadcast} the match live yesterday.} (Past Simple — attention irrégulier)
    \item \eng{Someone \textbf{has stolen} the camera equipment.} (Present Perfect)
    \item \eng{The director \textbf{will announce} the results tomorrow.} (Future)
    \item \eng{The producer \textbf{can edit} the video tonight.} (Modal)
    \item \eng{They \textbf{offered} the anchor a new contract.} (Double passif : donnez les deux versions)
\end{enumerate}
\end{exercice}

\begin{corrige}[Exercice : Entraînement immédiat — Transformation]
\begin{enumerate}
    \item \eng{News articles \textbf{are written} every day.} (\eng{by journalists} omis = générique).
    \item \eng{The microphone \textbf{is being repaired} right now.}
    \item \eng{The match \textbf{was broadcast} live yesterday.} (\textbf{Piège} : \eng{broadcast} = V3 identique à l'infinitif/prétérit. Pas \eng{*was broadcasted}).
    \item \eng{The camera equipment \textbf{has been stolen}.}
    \item \eng{The results \textbf{will be announced} tomorrow.}
    \item \eng{The video \textbf{can be edited} tonight.}
    \item \textbf{Version 1 (Focus personne) :} \eng{The anchor \textbf{was offered} a new contract.} \\
          \textbf{Version 2 (Focus objet) :} \eng{A new contract \textbf{was offered} \textbf{to} the anchor.}
\end{enumerate}
\end{corrige}

\courssub{Synthèse visuelle : La carte mentale de la formation}

\begin{popfigure}
\begin{tikzpicture}[
    mindmap,
    concept color=popblue,
    text=white,
    level 1 concept/.append style={level distance=4.5cm, sibling angle=90, font=\small\bfseries},
    level 2 concept/.append style={level distance=3cm, sibling angle=45, font=\footnotesize},
    level 3 concept/.append style={level distance=2.5cm, font=\tiny},
    every node/.style={concept},
    root concept/.append style={concept color=popdark, font=\normalsize\bfseries, minimum size=2.5cm},
    c1/.style={concept color=popgreen},
    c2/.style={concept color=poporange},
    c3/.style={concept color=poppurple},
    c4/.style={concept color=poppink},
]
\node[root concept, align=center]{Formation\\ Passive}
    child[c1] {node {Sujet = Ancien COD}
        child {node {Accord Nombre}}
        child {node {Cas : Nominatif}}}
    child[c2] {node {Auxiliaire BE}
        child {node {Porte le TEMPS}}
        child {node {Porte le NOMBRE}}
        child {node {Formes : am/is/are/was/were/been/be}}}
    child[c3] {node {Participe Passé (V3)}
        child {node {INVARIABLE}}
        child {node {Régulier : -ed}}
        child {node {Irrégulier : Liste V3}}}
    child[c4] {node {By-Agent (Optionnel)}
        child {node {Préposition BY}}
        child {node {Ancien Sujet}}
        child {node {Omettre si générique}}};
\end{tikzpicture}
\legende{Carte mentale : Les 4 piliers de la construction passive. Le centre (BE) est la clé de voûte temporelle.}
\end{popfigure}

\begin{attention}[Check-list anti-erreurs avant de valider votre phrase passive]
\begin{itemize}
    \item[$\square$] Ai-je mis \textbf{BE} au bon temps (celui de l'actif) ?
    \item[$\square$] Le participe passé est-il la forme \textbf{V3} (colonne 3 du tableau irréguliers) et non le prétérit ?
    \item[$\square$] Le sujet (ancien COD) est-il bien au \textbf{singulier/pluriel} correct pour accorder BE ?
    \item[$\square$] Si j'ai mis \eng{by}, est-ce que l'agent est vraiment nécessaire ? (\eng{by someone} $\rightarrow$ à supprimer).
    \item[$\square$] Pour les verbes à 2 compléments : ai-je mis \eng{to} devant la personne si l'objet est sujet ? (\eng{was given to him}).
    \item[$\square$] Le verbe est-il \textbf{transitif} ? (Pas de passif pour \eng{happen, sleep, arrive, die}).
\end{itemize}
\end{attention}

Cette section vous a donné les outils mécaniques pour construire n'importe quelle forme passive. La prochaine section (\og Emplois de la voix passive \fg) expliquera \textbf{quand} et \textbf{pourquoi} choisir la voix passive plutôt que l'active dans la communication médiatique et journalistique.

\coursec{The Active and Passive Voices : Formation, Emplois et Pièges}

\courssub{Rappel : Formation de la voix passive}

\begin{propriete}[Formation de la voix passive]
La voix passive se forme avec l'auxiliaire \cle{be} conjugué au temps requis + le \cle{participe passé} (V-en) du verbe lexical. Seuls les verbes \textbf{transitifs directs} (suivis d'un COD) peuvent se mettre à la voix passive.
\begin{center}
\begin{tabular}{|l|l|l|}
\hline
\rowcolor{popblueL} \textbf{Temps / Mode} & \textbf{Voix active} & \textbf{Voix passive} \\
\hline
Present Simple & \eng{They clean the office.} & \eng{The office \textbf{is cleaned}.} \\
\hline
Past Simple & \eng{They cleaned the office.} & \eng{The office \textbf{was cleaned}.} \\
\hline
Present Perfect & \eng{They have cleaned the office.} & \eng{The office \textbf{has been cleaned}.} \\
\hline
Past Perfect & \eng{They had cleaned the office.} & \eng{The office \textbf{had been cleaned}.} \\
\hline
Future (will) & \eng{They will clean the office.} & \eng{The office \textbf{will be cleaned}.} \\
\hline
Modals & \eng{They must clean the office.} & \eng{The office \textbf{must be cleaned}.} \\
\hline
Present Continuous & \eng{They are cleaning the office.} & \eng{The office \textbf{is being cleaned}.} \\
\hline
Past Continuous & \eng{They were cleaning the office.} & \eng{The office \textbf{was being cleaned}.} \\
\hline
\end{tabular}
\end{center}
\emph{Attention :} Les verbes intransitifs (\eng{go, sleep, arrive, happen, occur}) n'ont \textbf{jamais} de forme passive.
\end{propriete}

\courssub{Observation et découverte à partir d'un texte authentique}

Lisons l'extrait suivant, tiré d'un article de presse fictif rapportant un incident sur l'axe routier Douala--Yaoundé. Repérez toutes les formes passives et demandez-vous : \emph{qui fait l'action ? Est-il mentionné ? Est-il important ?}

\begin{textebox}[Article de presse : « Major Accident on Douala-Yaoundé Highway »]
Early yesterday morning, a serious accident \textbf{was reported} near the Edea junction. A fuel tanker \textbf{was hit} by a speeding truck coming from the opposite direction. The driver of the tanker \textbf{was trapped} inside the cabin for over two hours before he \textbf{was rescued} by firefighters. He \textbf{was rushed} to the regional hospital where he \textbf{is being treated} for severe burns. 

The highway \textbf{was closed} for six hours while the wreckage \textbf{was cleared}. Motorists \textbf{were advised} to use the alternative route through Bafia. An investigation \textbf{has been opened} by the police to determine the exact causes of the crash. It \textbf{is believed} that brake failure \textbf{might have caused} the collision. No fatalities \textbf{have been reported} so far.
\end{textebox}

\begin{definition}[Vocabulaire clé : Article de presse (Accident routier)]
\begin{itemize}
    \item \textbf{fuel tanker} (n.) : camion-citerne (carburant)
    \item \textbf{junction} (n.) : carrefour, jonction, échangeur
    \item \textbf{speeding} (adj./part.) : qui roule trop vite, en excès de vitesse
    \item \textbf{wreckage} (n.) : épave, débris (d'un véhicule, avion)
    \item \textbf{brake failure} (n.) : défaillance des freins
    \item \textbf{fatalities} (n. pl.) : morts, victimes mortelles
    \item \textbf{so far} (adv.) : jusqu'à présent, pour l'instant
    \item \textbf{alternative route} (n.) : itinéraire de déviation, route alternative
\end{itemize}
\end{definition}

\begin{center}
\begin{tabularx}{\linewidth}{|X|X|X|}
\hline
\rowcolor{popblueL} \textbf{Forme passive repérée} & \textbf{Agent (si mentionné)} & \textbf{Pourquoi le passif ?} \\
\hline
was reported & -- (inconnu) & L'auteur de l'information (la police, un témoin) n'est pas identifié. \\
\hline
was hit & by a speeding truck & L'agent (le camion) est l'information nouvelle, mise en focus à la fin. \\
\hline
was trapped / was rescued / was rushed / is being treated & by firefighters (une fois) & L'accent est mis sur la victime (le chauffeur) et la chronologie des secours. \\
\hline
was closed / was cleared & -- (évident : les autorités/ouvriers) & L'agent est prévisible, l'important est l'état de la route. \\
\hline
were advised & -- (autorités) & Conseil impersonnel, style administratif. \\
\hline
has been opened & by the police & Information officielle, l'agent est pertinent pour la crédibilité. \\
\hline
is believed & -- (opinion générale) & Structure impersonnelle pour rapporter une rumeur/hypothèse sans s'engager. \\
\hline
have been reported & -- & Statistique, focus sur le chiffre (zéro). \\
\hline
\end{tabularx}
\end{center}

\begin{attention}[Analyse fine : \eng{might have caused}]
La forme \eng{might have caused} (dans \eng{brake failure might have caused the collision}) est à la voix \textbf{active} (modal perfect active). Le sujet \eng{brake failure} fait l'action. Elle ne figure pas dans le tableau des passifs ci-dessus. Le seul passif de cette phrase est \eng{is believed}.
\end{attention}

\courssub{Les quatre emplois principaux de la voix passive}

\begin{propriete}[Règle générale : quand choisir le passif ?]
On utilise la voix passive dans les cas suivants :
\begin{enumerate}
    \item \textbf{L'agent est inconnu, indéterminé ou impossible à identifier.} C'est l'emploi le plus fréquent à l'oral comme à l'écrit.
    \item \textbf{L'agent est évident, prévisible ou sans intérêt informationnel.} Le mentionner alourdirait la phrase (ex. : « par la police », « par les médecins », « par quelqu'un »).
    \item \textbf{On veut mettre l'accent sur l'action, le résultat ou la victime (le patient), et non sur l'auteur.} C'est la norme dans le style \textbf{journalistique}, \textbf{scientifique}, \textbf{technique} et \textbf{administratif}.
    \item \textbf{Par politesse, discrétion ou prudence,} pour éviter d'accuser quelqu'un directement ou pour adoucir une responsabilité (style diplomatique, « \eng{corporate speak} »).
\end{enumerate}
\end{propriete}

\courssub{Emploi 1 : L'agent est inconnu}

C'est le cas typique des vols, découvertes, accidents non élucidés. En français, on utilise souvent « on » + actif ou le pronominal passif (« s'est fait voler »).

\begin{exemplebox}[Agent inconnu]
\eng{My phone \textbf{was stolen} in the bendskin yesterday.} \\
« On a volé mon téléphone dans le bendskin hier. » (Je ne sais pas qui.)

\eng{A new species of frog \textbf{has been discovered} in the Korup Forest.} \\
« Une nouvelle espèce de grenouille a été découverte dans la forêt de Korup. » (Les scientifiques ne sont pas nommés, l'espèce est le sujet.)
\end{exemplebox}

\courssub{Emploi 2 : L'agent est évident ou sans importance}

Si tout le monde sait \emph{qui} fait l'action (la police arrête, les médecins soignent, les ouvriers réparent), le mentionner avec \eng{by} est redondant.

\begin{exemplebox}[Agent évident]
\eng{The thief \textbf{was arrested} last night.} \\
« Le voleur a été arrêté hier soir. » (Forcément par la police / les gendarmes. Inutile de dire \eng{by the police}.)

\eng{The road \textbf{is being repaired} this week.} \\
« La route est en cours de réparation cette semaine. » (Par des ouvriers, c'est évident.)

\eng{She \textbf{was operated on} immediately.} \\
« Elle a été opérée immédiatement. » (Par des chirurgiens, au bloc opératoire.)
\end{exemplebox}

\courssub{Emploi 3 : Focus sur l'action / le résultat (Style journalistique, scientifique, technique)}

Le sujet de la phrase passive est ce qui nous intéresse vraiment : le vaccin, la loi, le livre, l'expérience. L'agent (les chercheurs, le parlement, l'auteur) est secondaire ou connu.

\begin{exemplebox}[Focus sur le résultat / Style objectif]
\eng{A new vaccine \textbf{has been developed} against malaria.} \\
« Un nouveau vaccin a été mis au point contre le paludisme. » (L'accent est sur le vaccin, pas sur « by researchers ».)

\eng{The data \textbf{were analysed} using SPSS software.} \\
« Les données ont été analysées à l'aide du logiciel SPSS. » (Style scientifique : l'agent « we/les chercheurs » est gommé pour l'objectivité.)

\eng{The new budget \textbf{was approved} by Parliament yesterday.} \\
« Le nouveau budget a été adopté par le Parlement hier. » (Ici \eng{by Parliament} est gardé car c'est l'institution officielle, info importante.)
\end{exemplebox}

\courssub{Emploi 4 : Politesse, discrétion, évitement de la responsabilité}

Très utilisé dans les communiqués officiels, les emails professionnels ou les excuses pour ne pas pointer du doigt un individu.

\begin{exemplebox}[Discrétion / « Corporate speak »]
\eng{Mistakes \textbf{were made} in the management of the fund.} \\
« Des erreurs ont été commises dans la gestion du fonds. » (Personne n'est accusé nommément. Comparez : \eng{The manager made mistakes} — accusatoire.)

\eng{Your application \textbf{cannot be processed} at the moment.} \\
« Votre demande ne peut pas être traitée pour le moment. » (Plus poli que \eng{We cannot process your application}.)
\end{exemplebox}

\courssub{L'agent : « by » (auteur) vs « with » (instrument)}

Lorsque l'agent est mentionné, la préposition change selon sa nature.

\begin{propriete}[By vs With]
\begin{itemize}
    \item \textbf{BY + Agent (personne, animal, force naturelle, institution) :} \emph{Who did it?} \\
    \eng{The novel \textbf{was written by} Chinua Achebe.} (Auteur humain)
    \eng{The house \textbf{was destroyed by} the storm.} (Force naturelle)
    \item \textbf{WITH + Instrument / Outil / Matériau / Ingrédient :} \emph{What was used?} \\
    \eng{The letter \textbf{was written with} a fountain pen.} (Outil)
    \eng{The cake \textbf{was made with} local cocoa.} (Ingrédient/Matériau)
\end{itemize}
\emph{Attention :} Si l'instrument est utilisé \emph{par} un agent, les deux peuvent apparaître : \eng{The safe was opened \textbf{by the thief} \textbf{with} a blowtorch.}
\end{propriete}

\begin{center}
\begin{tabularx}{\linewidth}{|X|X|X|}
\hline
\rowcolor{popblueL} \textbf{Phrase active} & \textbf{Passif avec BY (Agent)} & \textbf{Passif avec WITH (Instrument)} \\
\hline
The carpenter cut the wood with a saw. & The wood was cut \textbf{by the carpenter}. & The wood was cut \textbf{with a saw}. \\
\hline
The chef prepared the dish with spices. & The dish was prepared \textbf{by the chef}. & The dish was prepared \textbf{with spices}. \\
\hline
The police caught the thief with a trap. & The thief was caught \textbf{by the police}. & The thief was caught \textbf{with a trap}. \\
\hline
\end{tabularx}
\end{center}

\courssub{Structures impersonnelles : « It is said that... » / « He is said to be... »}

\emph{Point crucial pour le Baccalauréat et la compréhension de la presse.} Ces structures permettent de rapporter des rumeurs, des ouï-dire, des croyances communes sans citer de source précise. Elles équivalent au français « On dit que... », « Il est dit que... », « Il paraît que... ».

\begin{propriete}[Deux constructions équivalentes]
Soit un verbe de parole ou de pensée au passif (\eng{say, think, believe, know, expect, report, consider, suppose, understand, allege}) + un sujet fictif \eng{It} + une proposition \eng{that} \textbf{OU} le sujet réel « remonté » + infinitif.

\begin{center}
\begin{tabular}{|l|l|l|}
\hline
\rowcolor{popblueL} \textbf{Structure A : It + Passif + That-clause} & \textbf{Structure B : Sujet réel + Passif + Infinitif} & \textbf{Traduction française} \\
\hline
\eng{It is said that he is rich.} & \eng{He is said to be rich.} & On dit qu'il est riche. \\
\hline
\eng{It was thought that he had left.} & \eng{He was thought to have left.} & On pensait qu'il était parti. \\
\hline
\eng{It is expected that prices will rise.} & \eng{Prices are expected to rise.} & On s'attend à ce que les prix augmentent. \\
\hline
\eng{It is reported that two soldiers died.} & \eng{Two soldiers are reported to have died.} & On signale que deux soldats sont morts. \\
\hline
\end{tabular}
\end{center}

\textbf{Règle de transformation (Structure A $\rightarrow$ B) :}
\begin{enumerate}
    \item Le sujet de la proposition \eng{that} devient le sujet de la phrase principale.
    \item Le verbe principal (passif) garde son temps.
    \item Le verbe de la proposition \eng{that} passe à l'infinitif (\eng{to be} / \eng{to have been} / \eng{to + V}).
\end{enumerate}

\end{propriete}

\begin{savaistu}[Le passif impersonnel au Bac et dans les journaux]
Au Cameroun comme au Nigeria ou au Ghana, les titres de journaux abusent de \eng{It is alleged that...} (« Il est allégué que... ») ou \eng{He is alleged to have...} pour parler d'affaires judiciaires en cours sans présumer de la culpabilité (présomption d'innocence).
Ex. : \eng{The Minister \textbf{is alleged to have} embezzled public funds.} = « Le Ministre aurait détourné des fonds publics. » (Le conditionnel français « aurait » rend bien la nuance de rumeur non prouvée).
\end{savaistu}

\courssub{Organigramme de décision : Faut-il mettre l'agent ?}

\begin{popfigure}
\begin{tikzpicture}[
    node distance=1.8cm and 2.5cm,
    decision/.style={diamond, draw=popblue, thick, fill=popblueL, text width=3.5cm, align=center, inner sep=4pt},
    process/.style={rectangle, draw=popdark, thick, fill=popcream, text width=3.5cm, align=center, inner sep=4pt, rounded corners=4pt},
    startstop/.style={ellipse, draw=popgreen, thick, fill=popgreenL, text width=2.5cm, align=center, inner sep=4pt},
    arrow/.style={-Stealth, thick, draw=popdark},
    yesno/.style={midway, fill=white, inner sep=2pt, font=\small\bfseries, text=popdark}
]
\node (start) [startstop] {Choix de la voix};
\node (q1) [decision, below of=start] {L'agent est-il \textbf{connu} ?};
\node (q2) [decision, below left of=q1, xshift=-1cm] {L'agent est-il \textbf{important / informatif} ?};
\node (p1) [process, below right of=q1, xshift=1cm, align=center]{\textbf{Voix passive SANS agent}\\(Agent omis)\\\eng{The car was stolen.}};
\node (p2) [process, below left of=q2, align=center]{\textbf{Voix passive AVEC \eng{BY}}\\(Agent gardé)\\\eng{By Chinua Achebe}};
\node (p3) [process, below right of=q2, align=center]{\textbf{Voix ACTIVE préférée}\\(Agent = sujet)\\\eng{My brother repaired it.}};
\node (inst) [decision, below of=q2, yshift=-0.5cm] {S'agit-il d'un \textbf{instrument} ?};
\node (p4) [process, below left of=inst, align=center]{\textbf{Passif final (BY + WITH)}\\\eng{... by the thief with a torch.}};
\node (end) [startstop, below of=inst, yshift=-1.5cm] {Phrase finale};

\draw [arrow] (start) -- (q1);
\draw [arrow] (q1) -- node [yesno, anchor=east] {NON} (p1);
\draw [arrow] (q1) -- node [yesno, anchor=west] {OUI} (q2);
\draw [arrow] (q2) -- node [yesno, anchor=north east] {OUI} (p2);
\draw [arrow] (q2) -- node [yesno, anchor=north west] {NON} (p3);
\draw [arrow] (p2) -- (inst);
\draw [arrow] (inst) -- node [yesno, anchor=east] {OUI} (p4);
\draw [arrow] (inst) -- node [yesno, anchor=west] {NON} (end);
\draw [arrow] (p1) |- (end);
\draw [arrow] (p3) |- (end);
\draw [arrow] (p4) -- (end);
\end{tikzpicture}
\legende{Organigramme de décision pour l'expression de l'agent (by), de l'instrument (with) ou le choix de l'actif.}
\end{popfigure}

\courssub{Exercice résolu : Transformation et justification}

\begin{exoresolu}[Transformation active $\leftrightarrow$ passif + choix de l'agent]
\textbf{Énoncé :} Transformez les phrases actives en passives. Décidez si l'agent doit être gardé (\eng{by}), remplacé par l'instrument (\eng{with}) ou omis. Justifiez votre choix en français.

\begin{enumerate}
    \item Someone has cleaned the office. (Le nettoyeur est inconnu / sans importance)
    \item Chimamanda Ngozi Adichie wrote \emph{Half of a Yellow Sun}. (Auteur célèbre, info cruciale)
    \item The mechanic repaired the engine with a new tool. (Focus sur le moteur ; mécanicien évident, outil nouveau)
    \item They say that the President will visit Buea next month. (Structure impersonnelle « On dit que... »)
    \item The storm destroyed several houses in Limbe. (Force naturelle = agent gardé avec \eng{by})
\end{enumerate}

\tcblower
\textbf{Solution détaillée :}

\begin{enumerate}
    \item \eng{The office \textbf{has been cleaned}.} \\
    \emph{Justification :} Agent « someone » indéterminé $\rightarrow$ omis (Emploi 1).
    \item \eng{\emph{Half of a Yellow Sun} \textbf{was written by} Chimamanda Ngozi Adichie.} \\
    \emph{Justification :} L'auteur est l'information principale, la raison d'être du livre $\rightarrow$ \eng{by} gardé (Emploi 3, focus sur l'œuvre mais agent prestigieux).
    \item \eng{The engine \textbf{was repaired with} a new tool.} \\
    \emph{Justification :} « The mechanic » est évident (Emploi 2 $\rightarrow$ omis). « A new tool » est l'info nouvelle, l'instrument $\rightarrow$ \eng{with} (Emploi instrument).
    \item \eng{The President \textbf{is said to visit} Buea next month.} \quad OU \quad \eng{It \textbf{is said that} the President will visit Buea next month.} \\
    \emph{Justification :} Verbe \eng{say} au passif + sujet réel « The President » remonté + infinitif \eng{to visit} (Structure B impersonnelle).
    \item \eng{Several houses in Limbe \textbf{were destroyed by} the storm.} \\
    \emph{Justification :} « The storm » est une force naturelle, cause de l'événement, non un humain « évident » comme un médecin. L'agent est informatif $\rightarrow$ \eng{by} gardé.
\end{enumerate}
\end{exoresolu}

\courssub{Tableau récapitulatif des emplois}

\begin{center}
\begin{tabularx}{\linewidth}{|X|X|X|}
\hline
\rowcolor{popblueL} \textbf{Contexte / Intention} & \textbf{Exemple type} & \textbf{Traduction / Équivalent français} \\
\hline
Agent \textbf{inconnu} (vol, découverte, accident) & \eng{My bag \textbf{has been stolen}.} & « On a volé mon sac. » / « Mon sac a été volé. » \\
\hline
Agent \textbf{évident} (métier, fonction) & \eng{The suspect \textbf{was arrested}.} & « Le suspect a été arrêté. » (par la police) \\
\hline
Focus sur \textbf{résultat / objet} (Science, Presse, Tech) & \eng{The vaccine \textbf{has been approved}.} & « Le vaccin a été homologué. » \\
\hline
\textbf{Politesse / Dépersonnalisation} (Admin, Pro) & \eng{The deadline \textbf{was missed}.} & « Le délai n'a pas été respecté. » (évite « You missed... ») \\
\hline
\textbf{Opinion / Rumeur} (Journal, Bac) & \eng{He \textbf{is believed to be} guilty.} & « On croit qu'il est coupable. » / « Il serait coupable. » \\
\hline
\textbf{Instrument / Moyen} & \eng{The door \textbf{was opened with} a key.} & « La porte a été ouverte \textbf{avec} une clé. » \\
\hline
\end{tabularx}
\end{center}

\begin{attention}[Piège majeur : L'abus du passif « à la française »]
En français, le passif est très fréquent et souvent neutre (\eng{« La décision a été prise »}). En anglais, \textbf{l'actif est la norme par défaut} dès que l'agent est connu et pertinent. Un élève francophone a tendance à systématiser le passif pour « faire anglais » ou traduire mot à mot le « a été + participe ».

\begin{center}
\begin{tabularx}{\linewidth}{|X|X|X|}
\hline
\rowcolor{poporangeL} \textbf{Phrase française (Passif naturel)} & \textbf{Erreur fréquente (Passif forcé)} & \textbf{Anglais naturel (Actif préféré)} \\
\hline
« Mon frère a réparé mon téléphone. » & \eng{My phone was repaired by my brother.} & \eng{\textbf{My brother repaired} my phone.} \\
\hline
« Quelqu'un a nettoyé la salle. » & \eng{The room was cleaned by someone.} & \eng{\textbf{Someone cleaned} the room.} (ou passif sans agent : \eng{The room was cleaned.}) \\
\hline
« Ils construisent un nouveau pont. » & \eng{A new bridge is being built by them.} & \eng{\textbf{They are building} a new bridge.} \\
\hline
\end{tabularx}
\end{center}

\cle{Règle d'or :} Si vous pouvez dire \eng{by my brother}, \eng{by them}, \eng{by someone} sans apporter d'information nouvelle $\rightarrow$ \textbf{Repassez à l'actif !} Le passif anglais est \textbf{marqué} (marked), l'actif est \textbf{non marqué} (unmarked).
\end{attention}

\courssub{Entraînement écrit et production orale}

\begin{exercice}[Entraînement écrit : Voix active / passive]
\textbf{Consigne :} Transformez les phrases suivantes en voix passive. Décidez si l'agent doit être exprimé (\eng{by}), l'instrument (\eng{with}) ou omis. Justifiez brièvement en français.
\begin{enumerate}
    \item The government has announced new measures for the school year. (Focus on measures, agent obvious/less important)
    \item A pickpocket stole my wallet in the market. (Agent unknown)
    \item The students are cleaning the classroom with brooms. (Focus on classroom, agent obvious, instrument mentioned)
    \item People believe that the new bridge will open next year. (Impersonal structure \eng{It is believed...} / \eng{The new bridge is believed...})
    \item The chef cooked the meal with fresh spices from the farm. (Focus on meal, agent obvious, instrument/new info)
    \item They cancelled the match because of heavy rain. (Focus on match, agent 'they' vague)
\end{enumerate}
\end{exercice}

\begin{corrige}[Exercice : Entraînement écrit]
\begin{enumerate}
    \item \eng{New measures for the school year \textbf{have been announced} (by the government).} \\
    \emph{Justification :} L'agent "the government" est prévisible (Emploi 2), on peut l'omettre. Le focus est sur les mesures.
    \item \eng{My wallet \textbf{was stolen} in the market.} \\
    \emph{Justification :} L'agent (un pickpocket) est inconnu (Emploi 1).
    \item \eng{The classroom \textbf{is being cleaned with} brooms.} \\
    \emph{Justification :} "The students" évident (omission). "Brooms" = instrument/new info $\rightarrow$ \eng{with}.
    \item \eng{The new bridge \textbf{is believed to open} next year.} \quad OU \quad \eng{It \textbf{is believed that} the new bridge will open next year.} \\
    \emph{Justification :} Structure impersonnelle (Emploi 5 / Rumeur). Transformation : Sujet "The new bridge" remonté, "will open" $\rightarrow$ "to open".
    \item \eng{The meal \textbf{was cooked with} fresh spices from the farm.} \\
    \emph{Justification :} "The chef" évident (omis). "Fresh spices" = instrument/ingrédient mis en avant $\rightarrow$ \eng{with}.
    \item \eng{The match \textbf{was cancelled} because of heavy rain.} \\
    \emph{Justification :} "They" vague/indéterminé (omis). La cause (pluie) est exprimée par la préposition "because of", pas par un agent \eng{by}.
\end{enumerate}
\end{corrige}

\begin{experience}[Activité orale : « Passive Detective »]
\textbf{Objectif :} Identifier la raison du choix du passif dans des phrases authentiques.
\textbf{Consigne :} En binôme. L'élève A lit une phrase (ci-dessous). L'élève B doit dire : « \emph{Passive used because...} » (Agent unknown / Obvious / Focus on result / Politeness / Impersonal).
\begin{enumerate}
    \item \eng{English is spoken in many countries.} (Focus on language / Agent obvious: people)
    \item \eng{The meeting has been cancelled.} (Politeness / Agent obvious: boss/organizer)
    \item \eng{The painting was stolen in 1990.} (Agent unknown)
    \item \eng{It is expected that inflation will drop.} (Impersonal / Rumour)
    \item \eng{The goal was scored by Mbappé in the 90th minute.} (Focus on event, Agent important = star player)
\end{enumerate}
\textbf{Production :} Chaque binôme invente 2 phrases passives (une sans agent, une avec \eng{by} justifié) sur l'actualité camerounaise (ex. CAN, routes, examens, agriculture).
\end{experience}

\coursec{Comparaison français/anglais et pièges}

Dans la section précédente, nous avons vu \emph{quand} utiliser la voix passive en anglais : agent inconnu ou sans importance, ton objectif des médias et des sciences, mise en avant du résultat plutôt que de l'auteur de l'action. Mais savoir \emph{quand} employer le passif ne suffit pas : il faut aussi savoir le \emph{construire sans faute}. Or c'est précisément ici que les élèves francophones trébuchent, car le français et l'anglais ne fonctionnent pas du tout de la même manière. Cette section est donc consacrée à un face-à-face systématique des deux langues, puis aux cinq pièges classiques qu'il faut apprendre à désamorcer.

\courssub{Le français \eng{on} : le premier grand piège}

\textbf{Observation.} Lis attentivement ces deux annonces que l'on peut voir dans une rédaction de journal à Douala :

\begin{exemplebox}[Observer avant de traduire]
\eng{English is spoken in this office.} — « On parle anglais dans ce bureau. »

\eng{Newspapers are sold here.} — « On vend des journaux ici. »

\eng{This bridge was built in 2010.} — « On a construit ce pont en 2010. »
\end{exemplebox}

Remarque bien : dans les trois phrases françaises, on utilise le pronom indéfini \emph{on}. Dans les trois phrases anglaises, il n'y a \textbf{aucun sujet équivalent à \emph{on}} : l'anglais utilise simplement la \cle{voix passive}. C'est la traduction naturelle, idiomatique, correcte.

\begin{attention}[Ne traduis jamais \emph{on} par \eng{one} dans ce contexte]
La tentation est grande de traduire mot à mot : « On parle anglais ici » $\rightarrow$ *\eng{One speaks English here.} Grammaticalement, cette phrase existe, mais elle sonne \textbf{faux, guindé, artificiel} : \eng{one} comme pronom indéfini est réservé à un anglais très formel ou littéraire (\eng{One never knows.} — « On ne sait jamais. »). Dans 95 \% des cas, le \emph{on} français se rend par le \textbf{passif anglais} :
\begin{itemize}
\item « On parle anglais ici. » $\rightarrow$ \eng{English is spoken here.} (et non *\eng{One speaks English here.})
\item « On vous demande à l'accueil. » $\rightarrow$ \eng{You are wanted at reception.}
\item « On a volé mon téléphone au marché. » $\rightarrow$ \eng{My phone has been stolen at the market.}
\end{itemize}
\end{attention}

\begin{propriete}[Règle de correspondance \emph{on} $\rightarrow$ passif]
Quand le français utilise \emph{on} avec une valeur \textbf{indéfinie} (on ne sait pas qui agit, ou cela n'a pas d'importance), l'anglais préfère la voix passive : \textbf{sujet-patient + be conjugué + participe passé}.

Exception : quand \emph{on} signifie clairement \emph{nous} (« On va au stade samedi » = « Nous allons au stade samedi »), on traduit par \eng{we} : \eng{We are going to the stadium on Saturday.} — et non par le passif !
\end{propriete}

\courssub{Le français \emph{se faire} + infinitif et le passif français}

Deuxième correspondance essentielle : le français possède des tournures comme \emph{se faire voler}, \emph{se faire licencier}, \emph{se faire opérer}. L'anglais n'a pas cette construction avec un infinitif : il utilise directement le passif.

\begin{exemplebox}[Exemples traduits]
\eng{He was robbed.} — « Il s'est fait voler. »

\eng{She was fired last week.} — « Elle s'est fait licencier la semaine dernière. »

\eng{My father was operated on in Buea.} — « Mon père s'est fait opérer à Buea. »

\eng{The journalist was interviewed on CRTV.} — « Le journaliste s'est fait interviewer sur CRTV. »
\end{exemplebox}

De même, le passif français classique (\emph{être} + participe passé \textbf{accordé}) correspond au passif anglais (\emph{be} + participe passé \textbf{invariable}) — et c'est là que se cache le troisième piège, que nous verrons plus bas.

\begin{center}
\begin{tabularx}{\linewidth}{|X|X|X|}
\hline
\rowcolor{popblueL} \textbf{Français} & \textbf{English} & \textbf{Remarque} \\
\hline
On parle anglais ici. & English is spoken here. & \emph{on} indéfini $\rightarrow$ passif \\
\hline
On a construit ce pont en 2010. & This bridge was built in 2010. & passé composé $\rightarrow$ \eng{was/were} + pp \\
\hline
Il s'est fait voler. & He was robbed. & \emph{se faire} + inf. $\rightarrow$ passif simple \\
\hline
Les lettres ont été écrit\textbf{es}. & The letters were written. & accord français, invariabilité anglaise \\
\hline
On (= nous) part demain. & We are leaving tomorrow. & \emph{on} = \emph{nous} $\rightarrow$ \eng{we}, pas de passif \\
\hline
\end{tabularx}
\end{center}

\begin{popfigure}
\begin{tikzpicture}[node distance=0.6cm]
\node[draw=popblue, fill=popblueL, rounded corners, text width=5.6cm, align=center, font=\small\bfseries] (fr) {FRANÇAIS};
\node[draw=poporange, fill=poporangeL, rounded corners, text width=5.6cm, align=center, font=\small\bfseries, right=1.6cm of fr] (en) {ENGLISH};
\node[draw=popline, fill=white, rounded corners, text width=5.6cm, align=center, font=\small, below=0.35cm of fr] (f1) {on + verbe\\ \emph{On vend des journaux.}};
\node[draw=popline, fill=white, rounded corners, text width=5.6cm, align=center, font=\small, below=0.35cm of en] (e1) {passive voice\\ \emph{Newspapers are sold.}};
\node[draw=popline, fill=white, rounded corners, text width=5.6cm, align=center, font=\small, below=0.3cm of f1] (f2) {se faire + infinitif\\ \emph{Il s'est fait voler.}};
\node[draw=popline, fill=white, rounded corners, text width=5.6cm, align=center, font=\small, below=0.3cm of e1] (e2) {passive voice\\ \emph{He was robbed.}};
\node[draw=popline, fill=white, rounded corners, text width=5.6cm, align=center, font=\small, below=0.3cm of f2] (f3) {être + pp \textbf{accordé}\\ \emph{Les filles ont été invitées.}};
\node[draw=popline, fill=white, rounded corners, text width=5.6cm, align=center, font=\small, below=0.3cm of e2] (e3) {be + pp \textbf{invariable}\\ \emph{The girls were invited.}};
\draw[-{Stealth}, thick, popgreen] (f1) -- (e1);
\draw[-{Stealth}, thick, popgreen] (f2) -- (e2);
\draw[-{Stealth}, thick, popgreen] (f3) -- (e3);
\end{tikzpicture}
\legende{Trois constructions françaises, une seule réponse anglaise : la voix passive.}
\end{popfigure}

\courssub{Les cinq pièges classiques des francophones}

\textbf{Piège 1 : oublier \eng{be}.} En français, on peut dire « la maison construite en 1990 » sans auxiliaire. En anglais, une phrase passive complète exige \textbf{toujours} l'auxiliaire \eng{be} conjugué au temps voulu.

\begin{attention}[Sans \eng{be}, pas de passif]
*\eng{The house built in 1990.} n'est pas une phrase : c'est un groupe nominal (ça signifie « la maison construite en 1990 », sans verbe). Il faut écrire : \eng{The house was built in 1990.} — « La maison a été construite en 1990. »

Attention particulière au présent : *\eng{Rice grown in Ndop.} $\rightarrow$ \eng{Rice is grown in Ndop.} — « On cultive le riz à Ndop. »
\end{attention}

\textbf{Piège 2 : mettre le temps sur le participe.} C'est l'auxiliaire \eng{be} qui porte le temps, jamais le participe passé. Le participe passé est figé.

\begin{attention}[Le temps se conjugue sur \eng{be}, pas sur le verbe principal]
*\eng{The letter is wrote.} est une double faute : \eng{is} (présent) + \eng{wrote} (prétérit). Dis : \eng{The letter is written.} (présent passif) ou \eng{The letter was written.} (prétérit passif) — « La lettre est écrite / a été écrite. »

Autre exemple : *\eng{The match was playeded} $\rightarrow$ \eng{The match was played.} — « Le match a été joué. »
\end{attention}

\textbf{Piège 3 : accorder le participe comme en français.} En français : « les filles ont été invité\textbf{es} », « les lettres écrit\textbf{es} ». En anglais, le participe passé est \textbf{invariable} : jamais de \emph{-s}, jamais de \emph{-e}.

\begin{propriete}[Règle d'invariabilité du participe passé anglais]
Le participe passé anglais ne s'accorde \textbf{jamais}, ni en genre ni en nombre, contrairement au participe passé français.
\begin{itemize}
\item \eng{The girls were invited.} — « Les filles ont été invitées. » (jamais *\eng{inviteds})
\item \eng{The letters were written.} — « Les lettres ont été écrites. » (jamais *\eng{writtens})
\item \eng{The roads have been repaired.} — « Les routes ont été réparées. »
\end{itemize}
Cette règle vaut aussi pour l'actif au present perfect : \eng{She has finished.} / \eng{They have finished.} — le participe ne change pas.
\end{propriete}

\textbf{Piège 4 : confondre \eng{by} (agent) et \eng{with} (instrument).} En français, la préposition \emph{par} sert aux deux : « tué par le voleur » et « tué par un couteau ». L'anglais distingue rigoureusement : \eng{by} introduit l'\textbf{agent} (celui qui agit), \eng{with} introduit l'\textbf{instrument} (l'objet avec lequel on agit).

\begin{exemplebox}[\eng{by} = agent, \eng{with} = instrument]
\eng{The bread was cut with a knife.} — « Le pain a été coupé avec un couteau. » (et non *\eng{by a knife})

\eng{The bread was cut by the baker.} — « Le pain a été coupé par le boulanger. »

\eng{The window was broken with a stone by the children.} — « La fenêtre a été cassée avec une pierre par les enfants. »

\eng{The article was written with a pen by the reporter.} — « L'article a été écrit avec un stylo par le journaliste. »
\end{exemplebox}

Test simple : si tu peux remplacer par \emph{avec} en français, utilise \eng{with} ; si tu peux dire « c'est X qui a fait l'action », utilise \eng{by}.

\textbf{Piège 5 : mettre au passif un verbe intransitif.} Seuls les verbes \textbf{transitifs} (qui acceptent un complément d'objet direct) peuvent passer à la voix passive. Les verbes intransitifs comme \eng{happen}, \eng{arrive}, \eng{die}, \eng{sleep}, \eng{occur}, \eng{fall} n'ont pas de passif.

\begin{attention}[Pas de passif sans COD]
*\eng{The accident was happened.} est impossible. Dis : \eng{The accident happened.} — « L'accident s'est produit / est arrivé. »

Autres fautes fréquentes :
\begin{itemize}
\item *\eng{The train was arrived late.} $\rightarrow$ \eng{The train arrived late.} — « Le train est arrivé en retard. »
\item *\eng{My grandfather was died in 2015.} $\rightarrow$ \eng{My grandfather died in 2015.} — « Mon grand-père est mort en 2015. »
\end{itemize}
Astuce : demande-toi « peut-on dire \emph{quelqu'un + verbe + quelque chose} ? ». \eng{Someone stole my phone} (oui, COD) $\rightarrow$ passif possible : \eng{My phone was stolen.} Mais on ne peut pas dire *\eng{someone happened the accident} $\rightarrow$ pas de passif.
\end{attention}

\courssub{Méthode de vérification et entraînement}

\begin{methode}[Vérifie ta phrase passive en 4 questions]
Avant de rendre ta copie ou de parler, passe chaque phrase passive par ce filtre :
\begin{enumerate}
\item \textbf{Y a-t-il \eng{be} conjugué ?} (\eng{is}, \eng{was}, \eng{has been}, \eng{will be}...) Sinon, ta phrase n'a pas de verbe.
\item \textbf{Le verbe principal est-il au participe passé ?} (\eng{written}, \eng{built}, \eng{stolen}...) Pas de prétérit, pas d'infinitif après \eng{be}.
\item \textbf{L'agent est-il vraiment nécessaire ?} S'il est évident ou sans intérêt, supprime le \eng{by}-groupe : \eng{The thief was arrested.} suffit, inutile d'ajouter \eng{by the police}.
\item \textbf{Le verbe est-il transitif ?} Si oui, le passif est possible ; si non (\eng{happen}, \eng{arrive}, \eng{die}), reviens à la voix active.
\end{enumerate}
\end{methode}

\begin{exoresolu}[Corrige les fautes]
Énoncé — Chaque phrase contient une erreur typique de francophone. Corrige-la et justifie :
\begin{enumerate}
\item \eng{One speaks French in Yaoundé.}
\item \eng{The article is wrote by a student.}
\item \eng{The girls were inviteds to the party.}
\item \eng{The meat was cut by a knife.}
\item \eng{The accident was happened near Bafoussam.}
\end{enumerate}
\tcblower
Solution détaillée :
\begin{enumerate}
\item \eng{French is spoken in Yaoundé.} — Le \emph{on} indéfini français se traduit par le passif, pas par \eng{one}.
\item \eng{The article is written by a student.} — Après \eng{is}, il faut le participe passé \eng{written} ; \eng{wrote} est le prétérit actif.
\item \eng{The girls were invited to the party.} — Le participe passé anglais est invariable : jamais de \emph{-s} comme en français (« invitées »).
\item \eng{The meat was cut with a knife.} — Un couteau est un instrument, pas un agent : préposition \eng{with}, pas \eng{by}.
\item \eng{The accident happened near Bafoussam.} — \eng{happen} est intransitif : il n'a pas de voix passive.
\end{enumerate}
\end{exoresolu}

\begin{exercice}[À toi de jouer]
Traduis en anglais en utilisant la voix passive quand c'est nécessaire :
\begin{enumerate}
\item \eng{Newspapers are sold at this kiosk.}
\item \eng{My brother's bag was stolen at Douala station.} (ou \eng{My brother was robbed of his bag at Douala station.})
\item \eng{These houses were built in 2015.}
\item \eng{The letter was written with a pencil.}
\item \eng{The match took place last night in Bamenda.}
\end{enumerate}
\end{exercice}

\begin{corrige}[Exercice « À toi de jouer »]
\begin{enumerate}
\item \eng{Newspapers are sold at this kiosk.} — \emph{on} indéfini $\rightarrow$ passif au présent.
\item \eng{My brother's bag was stolen at Douala station.} (ou \eng{My brother was robbed of his bag...}) — \emph{se faire voler} $\rightarrow$ passif au prétérit.
\item \eng{These houses were built in 2015.} — participe \eng{built} invariable, auxiliaire \eng{were} au pluriel.
\item \eng{The letter was written with a pencil.} — instrument $\rightarrow$ \eng{with}.
\item \eng{The match took place last night in Bamenda.} — attention : \emph{se dérouler} est intransitif (\eng{take place}), donc voix active obligatoire, pas de passif !
\end{enumerate}
\end{corrige}

\begin{aretenir}[L'essentiel de la section]
\begin{itemize}
\item Le \emph{on} indéfini français $\rightarrow$ voix passive anglaise, jamais \eng{one} dans ce contexte.
\item \emph{Se faire} + infinitif $\rightarrow$ passif simple : \eng{He was robbed.}
\item Une phrase passive = \eng{be} conjugué + participe passé \textbf{invariable}. Pas d'accord, pas de temps sur le participe.
\item \eng{by} = l'agent (qui agit) ; \eng{with} = l'instrument (avec quoi on agit).
\item Les verbes intransitifs (\eng{happen}, \eng{arrive}, \eng{die}...) n'ont pas de passif.
\end{itemize}
\end{aretenir}

\coursec{Application guidée et production}

Après avoir analysé les différences structurelles entre le français et l'anglais et identifié les pièges classiques (omission de \eng{be}, confusion entre participe passé et prétérit, mauvais usage de \eng{by}), nous passons à la mise en pratique. Cette section vous propose un entraînement progressif : de la simple identification à la production écrite autonome, en passant par la transformation systématique — compétence évaluée chaque année au Baccalauréat.

\courssub{Exercice 1 : Identification de la voix (actif / passif)}

\begin{exercice}[Voix active ou passive ?]
Lisez les six extraits d'articles de presse ci-dessous. Pour chaque phrase, indiquez si le verbe est à la \textbf{voix active} ou à la \textbf{voix passive}. Justifiez votre réponse en identifiant le sujet et le verbe.
\begin{enumerate}
\item \eng{The Minister of Education inaugurated the new library yesterday.}
\item \eng{A new curriculum was approved by the Parliament last week.}
\item \eng{Heavy rains have caused severe flooding in the North Region.}
\item \eng{The election results will be announced tomorrow evening.}
\item \eng{Our national team is playing the final match this Saturday.}
\item \eng{The contract had been signed before the deadline.}
\end{enumerate}
\end{exercice}

\begin{corrige}[Exercice 1]
\begin{enumerate}
\item \textbf{Active}. Sujet : \eng{The Minister of Education} (agent) ; Verbe : \eng{inaugurated} (prétérit actif). Le sujet \emph{fait} l'action.
\item \textbf{Passive}. Sujet : \eng{A new curriculum} (patient) ; Verbe : \eng{was approved} (\eng{be} au prétérit + participe passé). L'agent (\eng{by the Parliament}) est exprimé.
\item \textbf{Active}. Sujet : \eng{Heavy rains} ; Verbe : \eng{have caused} (present perfect actif).
\item \textbf{Passive}. Sujet : \eng{The election results} ; Verbe : \eng{will be announced} (\eng{be} au futur + participe passé). Agent omis.
\item \textbf{Active}. Sujet : \eng{Our national team} ; Verbe : \eng{is playing} (present continuous actif).
\item \textbf{Passive}. Sujet : \eng{The contract} ; Verbe : \eng{had been signed} (\eng{be} au past perfect + participe passé). Agent omis.
\end{enumerate}
\end{corrige}

\courssub{Exercice 2 : Complétion — forme passive correcte (be + participe passé)}

\begin{exercice}[Conjuguez le verbe à la voix passive au temps indiqué]
Complétez chaque phrase en mettant le verbe entre parenthèses à la \textbf{voix passive} au temps convenable. Attention à l'accord du participe passé (invariable en anglais) et à la forme de \eng{be}.
\begin{enumerate}
\item \eng{The news \_\_\_\_\_\_\_\_ (broadcast) every evening at 8 p.m.} \textit{(present simple)}
\item \eng{The stolen car \_\_\_\_\_\_\_\_ (find) by the police last night.} \textit{(past simple)}
\item \eng{The new bridge \_\_\_\_\_\_\_\_ (build) since January.} \textit{(present perfect)}
\item \eng{The exam papers \_\_\_\_\_\_\_\_ (mark) by the teachers next week.} \textit{(future simple)}
\item \eng{The letters \_\_\_\_\_\_\_\_ (send) before the holiday.} \textit{(past perfect)}
\item \eng{English \_\_\_\_\_\_\_\_ (speak) in many countries.} \textit{(present simple)}
\item \eng{The match \_\_\_\_\_\_\_\_ (cancel) because of the rain.} \textit{(past simple)}
\item \eng{A new hospital \_\_\_\_\_\_\_\_ (open) by the Prime Minister next month.} \textit{(future simple)}
\end{enumerate}
\end{exercice}

\begin{corrige}[Exercice 2]
\begin{enumerate}
\item \eng{is broadcast} (\eng{be} présent singulier + participe passé de \eng{broadcast} — invariable). \textit{Note : \eng{news} est singulier en anglais.}
\item \eng{was found} (\eng{be} prétérit singulier + participe passé irrégulier de \eng{find}).
\item \eng{has been built} (\eng{have been} + participe passé de \eng{build}).
\item \eng{will be marked} (\eng{will be} + participe passé régulier).
\item \eng{had been sent} (\eng{had been} + participe passé irrégulier de \eng{send}).
\item \eng{is spoken} (\eng{be} présent singulier + participe passé irrégulier de \eng{speak}).
\item \eng{was cancelled} (\eng{be} prétérit + participe passé \eng{cancelled} — double \eng{l} en anglais britannique).
\item \eng{will be opened} (\eng{will be} + participe passé régulier).
\end{enumerate}
\end{corrige}

\courssub{Exercice 3 : Transformation active → passive (avec et sans agent)}

\begin{methode}[Réussir la transformation active → passive au Bac]
Suivez ces \textbf{4 étapes} systématiquement :
\begin{enumerate}
\item \textbf{Soulignez le COD} (Complément d'Objet Direct) de la phrase active. C'est lui qui deviendra le sujet de la phrase passive.
\item \textbf{Identifiez le temps du verbe actif}. C'est ce temps que vous devrez reporter sur l'auxiliaire \eng{be}.
\item \textbf{Conjuguez \eng{be} à ce temps} (en l'accordant avec le \emph{nouveau} sujet) et ajoutez le \textbf{participe passé} du verbe principal.
\item \textbf{Décidez de l'agent} : si le sujet de la phrase active est \eng{someone}, \eng{people}, \eng{they} (générique), \eng{we} (générique) ou un agent évident/inconnu $\rightarrow$ \textbf{omettez \eng{by + agent}}. Sinon, conservez-le (\eng{by + agent}).
\end{enumerate}
\end{methode}

\begin{exemplebox}[Exemples traités pas à pas]
\textbf{Exemple 1 (avec agent) :}
Active : \eng{The students cleaned the classroom.}
\begin{itemize}
\item COD : \eng{the classroom} $\rightarrow$ Nouveau sujet.
\item Temps actif : \textit{past simple} ($\rightarrow$ \eng{was/were}).
\item Sujet singulier $\rightarrow$ \eng{was cleaned}.
\item Agent \eng{the students} pertinent $\rightarrow$ \eng{by the students}.
\end{itemize}
Passive : \eng{The classroom was cleaned by the students.}

\textbf{Exemple 2 (sans agent — sujet générique) :}
Active : \eng{Someone has stolen my bicycle.}
\begin{itemize}
\item COD : \eng{my bicycle} $\rightarrow$ Nouveau sujet.
\item Temps actif : \textit{present perfect} ($\rightarrow$ \eng{has/have been}).
\item Sujet singulier $\rightarrow$ \eng{has been stolen}.
\item Agent \eng{Someone} = inconnu/générique $\rightarrow$ \textbf{omission de \eng{by someone}}.
\end{itemize}
Passive : \eng{My bicycle has been stolen.}
\end{exemplebox}

\begin{exoresolu}[Transformation guidée]
Transformez la phrase suivante à la voix passive. Décidez si l'agent doit être conservé.
\eng{The organizers will welcome the visitors at the airport.}

\textbf{Solution détaillée :}
\begin{enumerate}
\item COD : \eng{the visitors} $\rightarrow$ devient sujet.
\item Verbe actif : \eng{will welcome} $\rightarrow$ Futur simple (\eng{will}).
\item \eng{be} au futur : \eng{will be} + participe passé \eng{welcomed}.
\item Sujet \eng{the visitors} (pluriel) $\rightarrow$ \eng{will be welcomed}.
\item Agent : \eng{The organizers} (spécifique, pertinent) $\rightarrow$ \eng{by the organizers}.
\end{enumerate}
\eng{The visitors will be welcomed at the airport by the organizers.}
\end{exoresolu}

\begin{exercice}[Transformez les phrases suivantes à la voix passive. Omettez l'agent quand il est générique ou inconnu.]
\begin{enumerate}
\item \eng{The government has launched a new youth employment programme.}
\item \eng{Someone stole my phone in the market yesterday.}
\item \eng{They are building a new stadium in Douala.}
\item \eng{The teacher corrected the exams last week.}
\item \eng{People speak English all over the world.}
\end{enumerate}
\end{exercice}

\begin{corrige}[Exercice 3]
\begin{enumerate}
\item \eng{A new youth employment programme has been launched by the government.} (Agent spécifique conservé).
\item \eng{My phone was stolen in the market yesterday.} (Agent \eng{Someone} omis).
\item \eng{A new stadium is being built in Douala.} (Agent \eng{They} générique omis ; \textit{present continuous} $\rightarrow$ \eng{is being built}).
\item \eng{The exams were corrected by the teacher last week.} (Agent spécifique conservé).
\item \eng{English is spoken all over the world.} (Agent \eng{People} générique omis).
\end{enumerate}
\end{corrige}

\courssub{Exercice 4 : Correction d'erreurs fréquentes (pièges francophones)}

\begin{attention}[Erreurs sanctionnées au Bac]
\begin{itemize}
\item \textbf{Omission de \eng{be}} : *\eng{The book written by Ngugi} $\rightarrow$ \eng{The book \textbf{was} written by Ngugi}.
\item \textbf{Mauvais participe} (forme base ou prétérit) : *\eng{The match was win} $\rightarrow$ \eng{The match \textbf{was won}} ; *\eng{The letter was wrote} $\rightarrow$ \eng{The letter \textbf{was written}}.
\item \textbf{Accord erroné du participe} : Le participe passé \textbf{ne s'accorde jamais} en anglais. *\eng{The letters were sended} $\rightarrow$ \eng{The letters were sent}.
\item \textbf{Confusion \eng{been}/\eng{being}} : *\eng{The work has being done} $\rightarrow$ \eng{The work \textbf{has been} done} (perfect) ; *\eng{The work is been done} $\rightarrow$ \eng{The work \textbf{is being} done} (continuous).
\item \textbf{\eng{By} + pronom sujet} : *\eng{The cake was eaten by they} $\rightarrow$ \eng{by \textbf{them}}.
\end{itemize}
\end{attention}

\begin{exercice}[Corrigez les erreurs dans les phrases suivantes. Réécrivez la phrase correcte.]
\begin{enumerate}
\item \eng{The match was won by our team yesterday.} \textit{(Cette phrase est-elle correcte ?)}
\item \eng{The novel was write by a famous author.}
\item \eng{The documents have been sent last Monday.}
\item \eng{The house is being painted by they.}
\item \eng{The results will be announce tomorrow.}
\end{enumerate}
\end{exercice}

\begin{corrige}[Exercice 4]
\begin{enumerate}
\item \textbf{Correcte.} \eng{The match was won by our team yesterday.} (Prétérit passif régulier \eng{won} = participe passé de \eng{win}, agent exprimé).
\item \eng{The novel \textbf{was written} by a famous author.} (Participe passé irrégulier de \eng{write} : \eng{wrote} (prétérit) / \eng{written} (participe)).
\item \eng{The documents \textbf{were sent} last Monday.} (Marqueur de temps passé \eng{last Monday} $\rightarrow$ \textit{past simple} passif \eng{were sent}, pas \textit{present perfect}). \textit{Note : si l'intention est antériorité, on dirait \eng{had been sent}}.
\item \eng{The house is being painted by \textbf{them}.} (Pronom objet \eng{them} après \eng{by}, pas \eng{they}).
\item \eng{The results will be \textbf{announced} tomorrow.} (Participe passé régulier \eng{announced}, pas forme base \eng{announce}).
\end{enumerate}
\end{corrige}

\courssub{Exercice 5 : Traduction « on » + passif}

\begin{propriete}[Traduire « on » par la voix passive]
En français, le pronom indéfini \textbf{« on »} (valeur passive, agent non précisé) se traduit en anglais par la \textbf{voix passive sans agent} (\eng{by someone} omis).
\begin{center}
\begin{tabularx}{\linewidth}{|X|X|X|}
\hline
\rowcolor{popblueL} \textbf{Français (on + actif)} & \textbf{Anglais (Passif sans agent)} & \textbf{Temps} \\
\hline
On diffuse les nouvelles à 20 h. & \eng{The news \textbf{is broadcast} at 8 p.m.} & Present simple \\
On a construit cet hôpital en 2020. & \eng{This hospital \textbf{was built} in 2020.} & Past simple \\
On a déjà réparé la route. & \eng{The road \textbf{has been repaired}.} & Present perfect \\
On ouvrira l'école lundi. & \eng{The school \textbf{will be opened} on Monday.} & Future simple \\
\hline
\end{tabularx}
\end{center}
\end{propriete}

\begin{exercice}[Traduisez en anglais en utilisant la voix passive (sans agent).]
\begin{enumerate}
\item On diffuse les nouvelles à 20 h.
\item On a élu un nouveau président la semaine dernière.
\item On parle anglais dans ce bureau.
\item On a déjà nettoyé la salle.
\item On construira un nouveau marché l'an prochain.
\end{enumerate}
\end{exercice}

\begin{corrige}[Exercice 5]
\begin{enumerate}
\item \eng{The news is broadcast at 8 p.m.} (\eng{news} est singulier en anglais).
\item \eng{A new president was elected last week.}
\item \eng{English is spoken in this office.}
\item \eng{The room has already been cleaned.}
\item \eng{A new market will be built next year.}
\end{enumerate}
\end{corrige}

\courssub{Production guidée : Rédaction d'un court article de presse}

\begin{methode}[Rédiger un article de presse scolaire (style journalistique)]
\begin{enumerate}
\item \textbf{Headline (Titre)} : Court, nominal, présent simple ou infinitif (\eng{School Wins Regional Trophy}).
\item \textbf{Lead (Chapeau)} : 1 phrase résumant \eng{Who, What, Where, When}. Souvent au passif (\eng{A trophy was awarded...}).
\item \textbf{Body (Corps)} : Développement chronologique. Alternez actif (actions des élèves) et passif (événements subis, décisions officielles). Utilisez \eng{by + agent} pour les autorités.
\item \textbf{Quote (Citation)} : Discours direct (\eng{"We are proud," said the principal.}) ou indirect.
\item \textbf{Tense consistency} : Prétérit pour les faits passés, présent pour les situations actuelles, passif pour l'institutionnel.
\end{enumerate}
\end{methode}

\begin{experience}[Atelier d'écriture : « Notre école en fête »]
Rédigez un court article de presse (80–100 mots) sur un événement récent de votre établissement (journée culturelle, remise de prix, match sportif, inauguration). \textbf{Contraintes obligatoires :}
\begin{itemize}
\item Utilisez \textbf{au moins 3 voix passives} (différents temps si possible).
\item Incluez \textbf{au moins une citation} (discours direct).
\item Respectez la structure : Titre — Chapeau — Corps — Conclusion.
\end{itemize}
\end{experience}

\begin{exemplebox}[Modèle d'article (environ 90 mots)]
\textbf{Government High School Yaoundé Celebrates 50th Anniversary}

\textit{Yaoundé, May 15} — A spectacular ceremony \textbf{was organised} yesterday to mark the 50th anniversary of Government High School Yaoundé. The school gates \textbf{were decorated} with green and yellow ribbons early in the morning. Over 2,000 alumni and students \textbf{were welcomed} by the principal, Mr. Mbarga. A new science laboratory \textbf{has been funded} by the Parents' Association and \textbf{will be inaugurated} next September. \eng{"This school has shaped generations of leaders,"} \textbf{declared} the Minister of Secondary Education in his speech. The day \textbf{ended} with a friendly football match.
\end{exemplebox}

\begin{savaistu}[Le passif au Baccalauréat Série A]
Les sujets du Baccalauréat A (épreuve écrite) proposent \textbf{régulièrement} un exercice de transformation : \eng{"Rewrite the following sentences in the passive voice."} ou \eng{"Put the verbs in brackets into the correct passive form."}
\begin{itemize}
\item Les temps les plus testés : \textbf{Present Simple, Past Simple, Present Perfect, Future Simple} (\eng{will be + V-en}).
\item Le \textit{Present Perfect Passive} (\eng{has/have been + V-en}) est un classique pour exprimer le résultat présent d'une action passée.
\item Entraînez-vous à transformer des phrases avec \eng{someone, people, they} (sujets génériques) $\rightarrow$ omission systématique de l'agent.
\item \textbf{Piège récurrent} : verbes à deux compléments (\eng{give, send, offer, tell}). Au passif, deux transformations possibles : \eng{He gave me a book} $\rightarrow$ \eng{I was given a book} (sujet = COI) \textbf{OU} \eng{A book was given to me} (sujet = COD). Les deux sont corrects ; la première est plus naturelle.
\end{itemize}
\end{savaistu}

\begin{popfigure}
\begin{tikzpicture}[
    scale=0.9,
    transform shape,
    node distance=0.8cm and 1.2cm,
    timebox/.style={rectangle, rounded corners=4pt, draw=popblue, thick, fill=popblueL, minimum width=3.2cm, minimum height=1.6cm, align=center, font=\small},
    arrow/.style={-Stealth, thick, popdark}
]
% Nœuds : Temps
\node[timebox, align=center] (pres){\textbf{Present Simple Passive}\\\eng{am/is/are + V-en}\\\eng{The news \textbf{is broadcast} daily.}};
\node[timebox, right=of pres, align=center] (past){\textbf{Past Simple Passive}\\\eng{was/were + V-en}\\\eng{The hall \textbf{was cleaned} yesterday.}};
\node[timebox, right=of past, align=center] (perf){\textbf{Present Perfect Passive}\\\eng{have/has been + V-en}\\\eng{The exams \textbf{have been marked}.}};
\node[timebox, right=of perf, align=center] (fut){\textbf{Future Simple Passive}\\\eng{will be + V-en}\\\eng{A new lab \textbf{will be built}.}};

% Flèches
\draw[arrow] (pres) -- (past) node[midway, above, font=\tiny, popmuted] {prétérit};
\draw[arrow] (past) -- (perf) node[midway, above, font=\tiny, popmuted] {antériorité};
\draw[arrow] (perf) -- (fut) node[midway, above, font=\tiny, popmuted] {futur};

% Légende temps
\node[below=0.3cm of pres, font=\tiny\itshape, popmuted] {Habitude / Fait général};
\node[below=0.3cm of past, font=\tiny\itshape, popmuted] {Action datée terminée};
\node[below=0.3cm of perf, font=\tiny\itshape, popmuted] {Passé récent / Résultat présent};
\node[below=0.3cm of fut, font=\tiny\itshape, popmuted] {Prédiction / Décision};
\end{tikzpicture}
\legende{Frise des quatre temps passifs essentiels pour le Bac : structure \eng{be + V-en} et exemple contextuel.}
\end{popfigure}

\begin{aretenir}[Check-list finale avant l'examen]
\begin{itemize}
\item \cle{Structure} : \eng{be} (conjugué au temps requis) + \cle{Participe Passé} (V-en / colonne 3).
\item \cle{Sujet passif} = COD de l'actif. Accord de \eng{be} avec CE sujet.
\item \cle{Agent} : \eng{by + complément} seulement si pertinent (agent spécifique, connu, non générique).
\item \cle{« On »} $\rightarrow$ Passif sans agent.
\item \cle{Verbes à 2 compléments} : 2 passifs possibles (sujet = COI ou COD).
\item \cle{Verbes intransitifs} : \textbf{Jamais de passif} (\eng{happen, sleep, arrive, go, die, occur}...).
\item \cle{Prépositions} : Conservez la préposition du verbe à particule (\eng{look after} $\rightarrow$ \eng{is looked after}).
\end{itemize}
\end{aretenir}

\newpage

\coursec{Activité d'intégration}

\coursec{Lesson 56 — Grammar: The active and passive voices}

\courssub{Observation et découverte}
\begin{textebox}{Authentic Material: News Report Excerpt (Cameroon Tribune Style)}
“Good evening. The new science complex at the Government Bilingual High School Buea \textbf{was officially inaugurated} yesterday by the Minister of Secondary Education. The ceremony \textbf{was attended} by hundreds of stakeholders. In sports, the Lions of Buea \textbf{have been qualified} for the FENASSCO regional final. The match \textbf{will be played} next Friday. Meanwhile, the PTA General Assembly \textbf{has been postponed} due to the ongoing transport strike. A new date \textbf{will be announced} soon. Finally, the Clean Campus Campaign \textbf{is being launched} today by the Student Council.”
\end{textebox}

\begin{experience}[Activité d'observation guidée (10 min)]
\begin{enumerate}
    \item Lisez le texte ci-dessus. Surlignez tous les verbes qui ne sont pas à la forme active habituelle (sujet + verbe + objet).
    \item Pour chaque verbe surligné, identifiez : le \textbf{sujet} (qui subit l'action), l'\textbf{auxiliaire \eng{be}} (sa forme et son temps), le \textbf{participe passé} (V-en), et l'\textbf{agent} (introduit par \eng{by}) s'il est présent.
    \item Comparez : \eng{The Minister inaugurated the building} (Actif) vs \eng{The building was inaugurated by the Minister} (Passif). Que devient l'objet de l'actif ? Que devient le sujet de l'actif ?
\end{enumerate}
\end{experience}

\begin{aretenir}[Découverte : La transformation Actif $\rightarrow$ Passif]
\begin{itemize}
    \item \textbf{Sujet Actif} (Agent) $\rightarrow$ \textbf{Complément d'agent} (\eng{by} + agent) ou suppression.
    \item \textbf{Objet Actif} (Patient) $\rightarrow$ \textbf{Sujet Passif} (place initiale, accord avec l'auxiliaire \eng{be}).
    \item \textbf{Verbe Actif} $\rightarrow$ \textbf{\eng{be} (même temps) + Participe Passé (V-en)}.
    \item \textbf{Règle d'or :} Seuls les verbes \textbf{transitifs directs} (suivis d'un objet direct) peuvent passer à la voix passive. \eng{Happen, die, sleep, go, arrive} $\rightarrow$ \textbf{Jamais de passif}.
\end{itemize}
\end{aretenir}

\courssub{Formation de la voix passive}
\begin{propriete}[Structure fondamentale]
La voix passive se forme avec l'auxiliaire \textbf{\eng{be}} conjugué au temps voulu + le \textbf{participe passé (V-en)} du verbe lexical.
\[
\textbf{Sujet} + \eng{be}_{\text{(temps)}} + \textbf{V-en} + [\eng{by} + \text{Agent}]
\]
\begin{center}
\begin{tabularx}{\linewidth}{|X|X|X|}
\hline
\rowcolor{popblueL} \textbf{Temps / Modalité} & \textbf{Structure \eng{be} + V-en} & \textbf{Exemple contextuel (Bulletin scolaire)} \\
\hline
Present Simple & \eng{am / is / are + V-en} & \eng{The bulletin \textbf{is broadcast} every Friday.} \\
\hline
Present Continuous & \eng{am / is / are being + V-en} & \eng{The labs \textbf{are being inspected} right now.} \\
\hline
Present Perfect & \eng{has / have been + V-en} & \eng{Three prizes \textbf{have been awarded} to our students.} \\
\hline
Past Simple (Prétérit) & \eng{was / were + V-en} & \eng{The meeting \textbf{was postponed} yesterday.} \\
\hline
Past Continuous & \eng{was / were being + V-en} & \eng{The campaign \textbf{was being launched} when we arrived.} \\
\hline
Future Simple & \eng{will be + V-en} & \eng{The final \textbf{will be played} next Friday.} \\
\hline
Future Perfect & \eng{will have been + V-en} & \eng{The work \textbf{will have been finished} by June.} \\
\hline
Modaux (\eng{must, should, can, etc.}) & \eng{modal + be + V-en} & \eng{The report \textbf{must be submitted} tomorrow.} \\
\hline
\eng{Be going to} & \eng{am / is / are going to be + V-en} & \eng{New extinguishers \textbf{are going to be installed}.} \\
\hline
\end{tabularx}
\end{center}
\textbf{Points clés :} Pas d'auxiliaire \eng{do/does/did} à la passive. L'auxiliaire est \textbf{toujours \eng{be}}. Le participe passé \textbf{ne s'accorde jamais} en anglais (ni en genre, ni en nombre).
\end{propriete}

\begin{exemplebox}[Verbes irréguliers fréquents dans les médias (Colonne 3 = Participe Passé)]
\begin{center}
\begin{tabular}{|l|l|l|l|}
\hline
\rowcolor{popblueL} \textbf{Infinitif} & \textbf{Prétérit} & \textbf{Participe Passé (V-en)} & \textbf{Traduction} \\
\hline
\eng{build} & \eng{built} & \eng{built} & construire \\
\eng{write} & \eng{wrote} & \eng{written} & écrire \\
\eng{choose} & \eng{chose} & \eng{chosen} & choisir \\
\eng{put} & \eng{put} & \eng{put} & mettre, placer \\
\eng{send} & \eng{sent} & \eng{sent} & envoyer \\
\eng{broadcast} & \eng{broadcast} & \eng{broadcast} & diffuser \\
\eng{lead} & \eng{led} & \eng{led} & diriger, mener \\
\eng{win} & \eng{won} & \eng{won} & gagner \\
\eng{beat} & \eng{beat} & \eng{beaten} & battre \\
\eng{pay} & \eng{paid} & \eng{paid} & payer \\
\hline
\end{tabular}
\end{center}
\end{exemplebox}

\courssub{Emplois de la voix passive}
\begin{propriete}[Pourquoi choisir la voix passive ? (Valeurs pragmatiques)]
\begin{itemize}
    \item \textbf{Focus sur l'événement / le résultat (Thème = Sujet) :} \eng{« A new library \textbf{has been opened}. »} (L'ouverture compte, pas qui a coupé le ruban).
    \item \textbf{Agent inconnu, évident, générique ou indésirable :} \eng{« The match \textbf{was postponed}. »} (Inutile de dire « by the organizers »). \eng{« English \textbf{is spoken} here. »} (Agent générique = people).
    \item \textbf{Style impersonnel / objectif / administratif :} \eng{« It \textbf{is reported} that... »} / \eng{« Students \textbf{are advised} to... »} / \eng{« The form \textbf{must be filled} in. »}.
    \item \textbf{Agent important / prestigieux / responsable (introduit par \eng{by}) :} \eng{« The building \textbf{was inaugurated by the Minister}. »} (L'agent donne de l'autorité à l'info).
    \item \textbf{Cohésion textuelle (Given-New) :} Placer l'information connue (donnée) en sujet, l'information nouvelle au focus final. \eng{« The Minister arrived. He \textbf{was welcomed} by the Principal. »}
\end{itemize}
\end{propriete}

\begin{attention}[Verbes à double objet (Ditransitifs) : Deux passifs possibles]
\begin{center}
\begin{tabularx}{\linewidth}{|X|X|}
\hline
\rowcolor{popblueL} \textbf{Actif} & \eng{The Principal awarded \textbf{the students} (OI) \textbf{prizes} (OD).} \\
\hline
Passif 1 (Focus destinataire) & \eng{\textbf{The students} \textbf{were awarded} prizes (by the Principal).} \\
\hline
Passif 2 (Focus objet) & \eng{\textbf{Prizes} \textbf{were awarded} to the students (by the Principal).} \\
\hline
\end{tabularx}
\end{center}
\textbf{Note :} Au passif 1, l'OD (\eng{prizes}) reste après le verbe. Au passif 2, l'OI devient \eng{to} + objet. Les deux sont corrects ; le choix dépend du focus.
\end{attention}

\courssub{Comparaison français/anglais et pièges}
\begin{attention}[Check-list des erreurs fréquentes des francophones — \eng{Proofreading Checklist}]
\begin{enumerate}
    \item \textbf{Oublier l'auxiliaire \eng{be} :} \eng{*The library opened} (Actif ! / Adjectif) $\rightarrow$ \eng{The library \textbf{was} opened} (Passif passé). \eng{*The match postponed} $\rightarrow$ \eng{The match \textbf{was} postponed}.
    \item \textbf{Mauvais participe passé (Verbes irréguliers) :} \eng{*builded, *writed, *choosed, *putted, *sended, *broadcasted}. $\rightarrow$ \eng{built, written, chosen, put, sent, broadcast}. \textbf{Apprenez la colonne 3 par cœur.}
    \item \textbf{Accord du participe passé (Transfert FR $\rightarrow$ EN) :} \eng{*The prizes were awarded\textbf{e}s}. $\rightarrow$ \textbf{JAMAIS d'accord en anglais}. \eng{The prizes were awarded}. \eng{The letter was written}. \eng{The letters were written}.
    \item \textbf{Mauvaise place de l'agent \eng{by} :} \eng{*By the Minister the library was opened}. $\rightarrow$ \eng{The library was opened \textbf{by the Minister}} (en fin de proposition).
    \item \textbf{Confusion \eng{been} (Perfect) / \eng{being} (Continu) :} \eng{has \textbf{been} done} (Perfect) vs \eng{is \textbf{being} done} (Continu en cours). Prononciation : \eng{been} [bɪn], \eng{being} [biːɪŋ].
    \item \textbf{Passif impossible avec verbes intransitifs :} \eng{*The accident was happened}, \eng{*He was died}, \eng{*The meeting was occurred}. $\rightarrow$ \eng{The accident happened}, \eng{He died}, \eng{The meeting took place / occurred}.
    \item \textbf{Verbes à particule / prépositionnels (Phrasal / Prepositional verbs) :} Ne pas oublier la particule/préposition. \eng{They looked after the children} $\rightarrow$ \eng{The children were looked \textbf{after}}. \eng{They dealt with the problem} $\rightarrow$ \eng{The problem was dealt \textbf{with}}.
    \item \textbf{\eng{Get} vs \eng{Be} (Passif « get » : informel, souvent événementiel/involontaire) :} \eng{He got fired} (familier) vs \eng{He was fired} (standard). \textbf{Au Bac / écrit formel : privilégiez \eng{be}.}
\end{enumerate}
\end{attention}

\begin{center}
\begin{tabularx}{\linewidth}{|X|X|X|}
\hline
\rowcolor{popblueL} \textbf{Français (Accord PP)} & \textbf{English (No Agreement)} & \textbf{Remarque} \\
\hline
La lettre \textbf{est écrite}. & \eng{The letter \textbf{is written}.} & Singulier \\
Les lettres \textbf{sont écrites}. & \eng{The letters \textbf{are written}.} & \textbf{Pas de \eng{-s} à \eng{written}} \\
La maison \textbf{a été construite}. & \eng{The house \textbf{has been built}.} & \eng{built} invariable \\
Les maisons \textbf{ont été construites}. & \eng{The houses \textbf{have been built}.} & \textbf{Pas de \eng{-es} à \eng{built}} \\
\hline
\end{tabularx}
\end{center}

\courssub{Application guidée et production}

\courssub{Objectifs de l'activité}
\begin{itemize}
    \item \textbf{Mobiliser} la voix passive dans un contexte de communication authentique : la rédaction et la présentation d'un bulletin d'information scolaire (\eng{School News Bulletin}).
    \item \textbf{Construire} des phrases correctes à la voix passive à au moins quatre temps différents (présent, passé composé, prétérit, futur, modalités) dont une avec agent introduit par \eng{by}.
    \item \textbf{Identifier} à l'oral les structures passives entendues et \textbf{évaluer} leur correction grammaticale (choix de l'auxiliaire \eng{be}, forme du participe passé, place de l'agent).
    \item \textbf{Adopter} le registre formel et le style journalistique (concision, objectivité, vocabulaire des médias : \eng{to inaugurate, to award, to postpone, to launch, stakeholders, coverage, statement, anchor, source, to release}).
    \item \textbf{Travailler} en équipe : répartition des rôles (rédacteur, relecteur, présentateur, secrétaire), négociation du contenu, respect du temps de parole.
\end{itemize}

\courssub{Situation}
\begin{textebox}{Contexte : The Weekly Bulletin — Lycée Bilingue de Buea}
Vous êtes membres du club de presse \eng{« Young Reporters Club »} du Lycée Bilingue de Buea. Chaque vendredi, le club produit \eng{The Weekly Bulletin}, un journal lu de cinq minutes diffusé sur le système de sonorisation de l'établissement pendant la pause de 10 h, et publié sur le tableau d'affichage numérique de l'administration. Cette semaine, vous devez couvrir les événements suivants (réels ou plausibles) :
\begin{enumerate}
    \item L'inauguration officielle du nouveau bâtiment des sciences (laboratoires de physique, chimie, biologie) par le Ministre des Enseignements Secondaires ou son représentant.
    \item La qualification de l'équipe de football du lycée (les \eng{Lions of Buea}) pour la finale régionale des \eng{FENASSCO} (Fédération Nationale des Sports Scolaires) qui se jouera vendredi prochain au stade municipal.
    \item L'attribution de trois prix d'excellence au concours national de dissertation en anglais (organisé par l'Ambassade des États-Unis à Yaoundé) à trois élèves de Terminale A.
    \item Le report de l'assemblée générale des parents d'élèves (PTA) prévue initialement mardi, à cause de la grève des transporteurs (\eng{bendskin} et taxis) à Douala et Buea.
    \item Le lancement d'une campagne de nettoyage (\eng{Clean Campus Campaign}) par le bureau des élèves (\eng{Student Council}) en partenariat avec une ONG locale (\eng{Green Cameroon}).
    \item La visite surprise du Proviseur dans les dortoirs pour vérifier l'installation des nouveaux extincteurs.
\end{enumerate}
Votre mission : rédiger le script complet du bulletin (6 à 8 phrases), le répéter, puis le présenter en direct devant la classe (simulation studio CRTV).
\end{textebox}

\courssub{Consignes}
\begin{enumerate}
    \item \textbf{Analyse du corpus (10 min) :} En groupe, lisez la situation. Sélectionnez 6 à 8 informations à traiter. Pour chaque information, décidez si la voix passive est pertinente (focus sur l'action/la victime/le résultat plutôt que sur l'acteur). Cochez les temps verbaux envisagés.
    \item \textbf{Rédaction collaborative (25 min) :} Rédigez le script final sur une grande feuille (format A3) ou sur tablette.
        \begin{itemize}
            \item \textbf{Contrainte grammaticale obligatoire :} Au minimum \textbf{4 phrases à la voix passive} à des temps \textbf{différents} (ex. \eng{Present Perfect}, \eng{Simple Past}, \eng{Future Simple}, \eng{Modal + be + V-en}). L'une d'elles \textbf{doit} inclure un agent introduit par \eng{by}.
            \item \textbf{Contrainte lexicale :} Utilisez au moins 5 mots du lexique « Media \& Communication » vu en Chapter 5 (\eng{headline, to broadcast, coverage, anchor, source, to release, statement, audience}).
            \item \textbf{Style :} Ton neutre, phrases complètes, pas de langage SMS. Commencez par un \eng{lead} (accroche) type : \eng{« Good morning, this is The Weekly Bulletin, I am [Name] reporting from Buea. »}
        \end{itemize}
    \item \textbf{Relecture croisée (10 min) :} Échangez votre script avec un autre groupe. Vérifiez : (a) Formation correcte : \eng{be} (temps) + \eng{V-en} (participe passé). (b) Forme du participe passé (irréguliers : \eng{built, written, chosen, broadcast}). (c) Place de \eng{by} + agent (en fin de proposition). (d) Ponctuation et majuscules (noms propres : \eng{Buea, FENASSCO, Minister, Green Cameroon}). Signez la relecture.
    \item \textbf{Répétition \& Attribution des rôles (10 min) :} Désignez deux \eng{anchors} (présentateurs) qui liront le bulletin debout, micro en main (stylo), avec une diction claire, le bon rythme (chunking) et l'intonation descendante sur les affirmations. Les deux autres sont \eng{floor managers} (chronomètre, signes « slower », « louder »).
    \item \textbf{Présentation \& Écoute active (15 min) :} Chaque groupe passe au « studio » (devant le tableau). La classe écoute avec une grille d'évaluation : relever \textbf{toutes} les phrases passives entendues, noter le temps, l'agent (s'il y a), et cocher \eng{Correct / Erreur}. Notez aussi la prononciation de \eng{been} [bɪn], \eng{awarded}, \eng{postponed}.
    \item \textbf{Débat \& Affichage (10 min) :} Retour collectif. Le professeur corrige les erreurs majeures au tableau. La classe vote pour le « Best Bulletin » (critères : exactitude grammaticale, richesse lexicale, fluidité orale, respect du temps). Le script gagnant est affiché au tableau d'affichage de la classe pour la semaine.
\end{enumerate}

\courssub{Ressources à mobiliser}
\begin{textebox}{Lexique utile — Media \& School Events (Glossaire)}
\begin{center}
\begin{tabularx}{\linewidth}{|l|X|l|}
\hline
\rowcolor{popblueL} \textbf{English} & \textbf{Context / Collocation} & \textbf{Français} \\
\hline
\eng{to inaugurate} & \eng{to inaugurate a building} & inaugurer \\
\eng{to award} & \eng{to award a prize / a medal} & décerner, attribuer (prix) \\
\eng{to postpone} & \eng{to postpone a meeting to Friday} & reporter \\
\eng{to launch} & \eng{to launch a campaign} & lancer \\
\eng{to release} & \eng{to release a statement / a report} & publier, diffuser (communiqué) \\
\eng{coverage} & \eng{media coverage, live coverage} & couverture (médiatique) \\
\eng{anchor (n)} & \eng{the news anchor} & présentateur(trice) JT \\
\eng{statement} & \eng{to issue a statement} & communiqué, déclaration \\
\eng{source} & \eng{reliable source, anonymous source} & source \\
\eng{stakeholder} & \eng{key stakeholders (parents, admin)} & partie prenante \\
\eng{outbreak} & \eng{outbreak of violence / disease} & déclenchement, épidémie \\
\eng{feat} & \eng{a remarkable feat} & exploit, prouesse \\
\eng{premises} & \eng{school premises} & locaux, enceinte \\
\hline
\end{tabularx}
\end{center}
\end{textebox}

\courssub{Proposition de production et corrigé commenté}
\begin{exoresolu}{Modèle de bulletin — \eng{The Weekly Bulletin, Friday 17th May}}
\textbf{Script lu par les deux \eng{anchors} (alternance) :}

\eng{“Good morning everyone. This is \textbf{The Weekly Bulletin}, broadcasting live from the Lycée Bilingue de Buea. I am \textbf{Ngwa}, and I am \textbf{Mbolle}.}

\eng{\textbf{Headline one:} The new science complex \textbf{has been officially inaugurated} by the Minister of Secondary Education. The ceremony \textbf{was attended} by hundreds of students and parents yesterday.}

\eng{\textbf{Headline two:} Our football team, the Lions of Buea, \textbf{have qualified} for the FENASSCO regional final. The match \textbf{will be played} next Friday at the Municipal Stadium.}

\eng{\textbf{Headline three:} Three excellence prizes \textbf{were awarded} to our Terminale A students in the US Embassy essay competition. The winners \textbf{have been congratulated} by the Principal this morning.}

\eng{\textbf{Headline four:} Due to the transport strike in the Southwest Region, the PTA General Assembly \textbf{has been postponed} until further notice. A new date \textbf{will be announced} next week.}

\eng{\textbf{Headline five:} Finally, the \eng{Clean Campus Campaign} \textbf{is being launched} today by the Student Council in partnership with \eng{Green Cameroon}. All students \textbf{are requested} to participate actively.}

\eng{That is all for this week. Stay informed, stay disciplined. This was \textbf{The Weekly Bulletin}. Thank you.”}

\tcblower
\textbf{Commentaire linguistique détaillé (Corrigé commenté pour l'enseignant)}

\begin{enumerate}
    \item \textbf{\eng{has been officially inaugurated} (Present Perfect Passive) :} Temps choisi pour un événement récent (hier) avec conséquence présente (le bâtiment est ouvert). \eng{Officially} (adverbe) bien placé entre auxiliaires. \textbf{Agent :} \eng{by the Minister of Secondary Education} (agent important, autorité).
    \item \textbf{\eng{was attended} (Past Simple Passive) :} Récit d'un événement passé terminé (la cérémonie d'hier). Sujet \eng{The ceremony} (singulier) $\rightarrow$ \eng{was}. Participe régulier \eng{attended}.
    \item \textbf{\eng{have qualified} (Present Perfect Active) :} \textbf{Correction pédagogique :} \eng{to qualify} (se qualifier) est intransitif dans ce sens. La passive \eng{have been qualified} est possible mais lourde (implique une action externe de la fédération). L'actif \eng{have qualified} est l'usage standard et naturel ici. Le modèle corrige ce point pour refléter l'anglais authentique.
    \item \textbf{\eng{will be played} (Future Simple Passive) :} Prédiction programmée. \eng{will be + V-en}. Pas de \eng{will be being played}.
    \item \textbf{\eng{were awarded} (Past Simple Passive) :} Fait accompli dans le passé (cérémonie de remise des prix antérieure). \eng{Three prizes} (sujet pluriel) $\rightarrow$ \eng{were}. \eng{Awarded} (régulier).
    \item \textbf{\eng{have been congratulated} (Present Perfect Passive) :} Action récente (ce matin) lien avec le présent. Sujet \eng{The winners} (pluriel) $\rightarrow$ \eng{have}.
    \item \textbf{\eng{has been postponed} (Present Perfect Passive) :} Annonce d'une décision prise récemment, valable maintenant. \eng{The PTA General Assembly} (singulier) $\rightarrow$ \eng{has}. \eng{Postponed} (régulier).
    \item \textbf{\eng{will be announced} (Future Simple Passive) :} Engagement futur. Agent omis (évident : l'administration).
    \item \textbf{\eng{is being launched} (Present Continuous Passive) :} Action en cours au moment de parler (aujourd'hui). Structure \eng{is being + V-en}. \textbf{Piège :} ne pas écrire \eng{*is launched} (habitude) ni \eng{*has been launched} (terminé).
    \item \textbf{\eng{are requested} (Present Simple Passive) :} Ordre / Consigne formelle (style administratif). \eng{All students} (pluriel) $\rightarrow$ \eng{are}. Verbe \eng{request} + \eng{to} + infinitif (\eng{to participate}).
\end{enumerate}

\textbf{Bilan grammatical du modèle :} 8 phrases passives sur 10 phrases totales (hors intro/conclusion). Temps utilisés : Present Perfect (3), Past Simple (2), Future Simple (2), Present Continuous (1), Present Simple (1). Agent \eng{by} explicite : 2 fois (Minister, Principal, Student Council/Green Cameroon). Respect total des contraintes (variété temps, agent \eng{by}, lexique médias).
\end{exoresolu}

\courssub{Bilan de l'activité}
\begin{aretenir}[Ce qu'il faut retenir de l'activité d'intégration]
\begin{itemize}
    \item La voix passive n'est pas un exercice de grammaire isolé : c'est un \textbf{outil rhétorique majeur du journalisme} (écrit et oral) pour hiérarchiser l'information (mettre le fait nouveau en sujet).
    \item La maîtrise exige \textbf{deux compétences distinctes} : (1) La \textbf{formation morphologique} parfaite (\eng{be} + V-en, irréguliers, pas d'accord) et (2) Le \textbf{choix pragmatique} du temps (aspect) et de la présence/absence de l'agent \eng{by}.
    \item À l'oral (\eng{speaking}), la passive demande une \textbf{prononciation soignée} des formes faibles de \eng{been} [bɪn], \eng{being} [biːɪŋ], \eng{was} [wəz], \eng{were} [wə] et du \eng{-ed} final ([d], [t], [ɪd]) pour ne pas brouiller le message.
    \item Le travail en groupe (rédaction $\rightarrow$ relecture $\rightarrow$ oral) reproduit la \textbf{chaîne de production médiatique réelle} (journaliste $\rightarrow$ rédacteur en chef $\rightarrow$ présentateur) et développe l'autonomie et l'esprit critique.
\end{itemize}
\end{aretenir}

\begin{methode}[Comment évaluer cette activité (Grille professeur — 20 pts)]
\begin{enumerate}
    \item \textbf{Exactitude grammaticale (6 pts) :} 4 passifs obligatoires corrects (formation, temps, irréguliers) = 4 pts. Agent \eng{by} correct = 1 pt. Variété des temps (min 4 différents) = 1 pt.
    \item \textbf{Qualité lexicale et style (4 pts) :} 5 mots du chapitre utilisés correctement = 2 pts. Registre journalistique (lead, transitions, ton neutre) = 2 pts.
    \item \textbf{Expression orale / Prononciation (5 pts) :} Fluidité, chunking, intonation descendante, articulation des formes passives (\eng{been/being}, \eng{-ed}) = 3 pts. Respect du temps (3-4 min) = 1 pt. Contact visuel / posture \eng{anchor} = 1 pt.
    \item \textbf{Compétence collaborative (5 pts) :} Répartition rôles, respect délais, qualité relecture croisée (grille remplie), implication = 5 pts.
    \item \textbf{Total : 20 pts.} Note de groupe (15 pts) + Note individuelle oral (5 pts).
\end{enumerate}
\end{methode}

\begin{experience}[Prolongement possible — Projet « School TV »]
Si l'équipement le permet (smartphone + haut-parleur), enregistrez les bulletins (vidéo ou audio). Montez un \eng{« Best of »} mensuel diffusé sur l'écran de la salle des profs ou le groupe WhatsApp de la classe. Cela ancre la grammaire dans une production durable et valorisante.
\end{experience}

\begin{exercice}[Entraînement individuel — Transformation Actif $\rightarrow$ Passif]
Transformez les phrases actives suivantes en voix passive. Conservez le temps. N'ajoutez l'agent \eng{by} que s'il apporte une information essentielle.
\begin{enumerate}
    \item The Minister inaugurated the new science building yesterday.
    \item The federation has qualified our team for the final.
    \item The strike postponed the PTA meeting.
    \item The Student Council is launching the Clean Campus Campaign today.
    \item The Principal will announce the new date next week.
    \item Green Cameroon sponsors the campaign.
    \item They have installed new fire extinguishers in the dormitories.
    \item The jury awarded three prizes to our students.
    \item The anchor broadcasts the bulletin every Friday.
    \item You must submit the report before Monday.
\end{enumerate}
\end{exercice}

\begin{corrige}[Exercice — Transformation]
\begin{enumerate}
    \item The new science building \textbf{was inaugurated} by the Minister yesterday. (Agent important)
    \item Our team \textbf{has been qualified} for the final (by the federation). (Agent optionnel, souvent omis)
    \item The PTA meeting \textbf{was postponed} (due to the strike). (Agent omis, cause exprimée autrement)
    \item The Clean Campus Campaign \textbf{is being launched} today by the Student Council. (Agent important)
    \item The new date \textbf{will be announced} next week. (Agent omis, évident)
    \item The campaign \textbf{is sponsored} by Green Cameroon. (Agent important)
    \item New fire extinguishers \textbf{have been installed} in the dormitories. (Agent omis)
    \item Three prizes \textbf{were awarded} to our students (by the jury). / Our students \textbf{were awarded} three prizes (by the jury). (Deux passifs possibles)
    \item The bulletin \textbf{is broadcast} every Friday. (Agent omis, générique)
    \item The report \textbf{must be submitted} before Monday. (Agent omis, impersonnel)
\end{enumerate}
\end{corrige}

\begin{savaistu}[Le saviez-vous ? — La passive au Cameroun et dans les médias]
Au Cameroun anglophone (Northwest, Southwest), l'usage de la passive dans les médias (\eng{CRTV}, \eng{The Post}, \eng{Cameroon Tribune}) suit les normes de l'anglais standard britannique. On observe cependant une tolérance accrue pour le \textbf{\eng{get}-passive} à l'oral informel (\eng{« My phone got stolen »}). À l'écrit formel (administration, examens, journalisme), seul le \textbf{\eng{be}-passive} est attendu. Le programme de Terminale A exige la maîtrise des temps simples, parfaits, modaux et du \eng{be going to} à la passive, ainsi que la gestion des verbes à double objet (\eng{award, give, send, offer}) et des verbes prépositionnels (\eng{look after, deal with}). Le \eng{Baccalauréat} teste souvent la transformation de phrase (Item type : \eng{Rewrite in the passive voice}) et la détection d'erreurs (\eng{Error correction} : accord fantôme, \eng{been/being}, verbe intransitif).
\end{savaistu}

\newpage

\coursec{Méthodes et exercices résolus}

\courssub{Observation et découverte : Voix active vs Voix passive}

\begin{textebox}[Dialogue : Au marché de Mokolo (Yaoundé)]
-- \eng{Have you heard? The new market hall \textbf{was inaugurated} yesterday by the Mayor.}
-- \eng{Really? Who \textbf{built} it?}
-- \eng{A Chinese company \textbf{built} it. The works \textbf{were funded} by the State budget.}
-- \eng{And the old stalls?}
-- \eng{They \textbf{were demolished} last month. But the traders \textbf{were relocated} to a temporary site.}
-- \eng{Good. Let's hope the new hall \textbf{will be maintained} properly.}
\end{textebox}

\begin{definition}[Voix Active et Voix Passive]
La \cle{voix active} présente le sujet comme \textbf{l'agent} qui accomplit l'action : \eng{Subject + Verb + Object}.\\
La \cle{voix passive} présente le sujet comme \textbf{le patient} qui subit l'action : \eng{Subject + be (conjugué) + Past Participle [+ by Agent]}.
\end{definition}

\begin{exemplebox}[Paire active / passive (Contexte camerounais)]
\begin{itemize}
    \item \textbf{Active :} \eng{The Minister \textbf{launched} the curriculum.} (Le Ministre \emph{fait} l'action.)
    \item \textbf{Passive :} \eng{The curriculum \textbf{was launched} by the Minister.} (Le curriculum \emph{subit} l'action.)
\end{itemize}
\eng{by the Minister} = \textbf{Complément d'Agent} (introduit par \eng{by}). Il est souvent omis s'il est vague (\eng{people, someone, they}) ou connu du contexte.
\end{exemplebox}

\begin{methode}[Identifier la voix active ou passive dans une phrase]
\textbf{Objectif :} Repérer instantanément si le sujet \emph{fait} l'action (voix active) ou \emph{subit} l'action (voix passive).
\begin{enumerate}
    \item \textbf{Trouver le verbe conjugué} (le noyau de la prédication).
    \item \textbf{Identifier le sujet} (qui/quoi + verbe).
    \item \textbf{Analyser la relation Sujet--Verbe} :
    \begin{itemize}
        \item Sujet = \textbf{Agent} (fait l'action) $\rightarrow$ \cle{Voix Active} (\eng{Subject + Verb + Object}).
        \item Sujet = \textbf{Patient} (subit l'action) $\rightarrow$ \cle{Voix Passive} (\eng{Subject + be + Past Participle [+ by Agent]}).
    \end{itemize}
    \item \textbf{Test de transformation} : Essayer d'ajouter \eng{by [quelqu'un]} après le verbe. Si c'est possible sans changer le sens, c'est une passive.
    \item \textbf{Attention aux verbes d'état (linking verbs)} (\eng{be, seem, become, look, feel, smell, taste, sound, appear, stay}) : ils ne sont \textbf{jamais} passifs. \eng{He is happy} / \eng{She looks tired} = voix active.
\end{enumerate}
\textbf{Réflexe Bac :} Souligne le verbe, entoure le sujet, trace une flèche : Sujet $\rightarrow$ Verbe (Actif) ou Verbe $\rightarrow$ Sujet (Passif).
\end{methode}

\begin{exoresolu}[Identification des voix : Contexte camerounais]
\textbf{Énoncé :} Pour chaque phrase ci-dessous, indique si le verbe est à la \textbf{voix active} ou à la \textbf{voix passive}. Justifie en identifiant le Sujet (S) et le Complément d'Agent (le cas échéant).
\begin{enumerate}
    \item The Minister of Secondary Education \textbf{launched} the new curriculum yesterday in Yaoundé.
    \item The new curriculum \textbf{was launched} by the Minister yesterday.
    \item Many students \textbf{prefer} practical subjects like Computer Science.
    \item Practical subjects \textbf{are preferred} by many students.
    \item The bendskin driver \textbf{looked} tired after the long journey to Buea.
    \item The GCE results \textbf{will be released} next August.
    \item A Chinese company \textbf{is building} the new bridge over the Wouri.
    \item The match \textbf{was won} by Canon Yaoundé in the final.
\end{enumerate}
\tcblower
\textbf{Solution détaillée :}
\begin{enumerate}
    \item \textbf{Active.} S = \eng{The Minister} (Agent). Verbe = \eng{launched} (Past Simple). Structure : S + V + O.
    \item \textbf{Passive.} S = \eng{The new curriculum} (Patient). Verbe = \eng{was launched} (be + Past Participle). Agent = \eng{by the Minister}.
    \item \textbf{Active.} S = \eng{Many students} (Agent). Verbe = \eng{prefer} (Present Simple). Verbe d'action, sujet agent.
    \item \textbf{Passive.} S = \eng{Practical subjects} (Patient). Verbe = \eng{are preferred} (be + PP). Agent = \eng{by many students}.
    \item \textbf{Active (Verbe d'état).} S = \eng{The bendskin driver}. Verbe = \eng{looked} (Past Simple de \eng{look} = sembler/avoir l'air). \textbf{Piège :} \eng{looked} n'est pas \eng{was looked}. Pas de \eng{by} possible. C'est un verbe d'état (linking verb), voix active obligatoire.
    \item \textbf{Passive.} S = \eng{The GCE results} (Patient). Verbe = \eng{will be released} (Future Passive). Agent omis (évident : le GCE Board).
    \item \textbf{Active.} S = \eng{A Chinese company} (Agent). Verbe = \eng{is building} (Present Continuous). Sujet agent, action en cours.
    \item \textbf{Passive.} S = \eng{The match} (Patient). Verbe = \eng{was won} (Past Simple Passive). Agent = \eng{by Canon Yaoundé}.
\end{enumerate}
\textbf{Conclusion :} La présence de \eng{be} + \emph{Participe Passé} + \eng{by} (explicite ou implicite) est le marqueur formel de la passive. L'absence d'agent ou la nature du verbe (état, intransitif) signalent l'active.
\end{exoresolu}

\courssub{Formation de la voix passive : La règle du \eng{be} + Participe Passé}

\begin{propriete}[Formule générale de la voix passive]
\textbf{Voix Passive = Sujet (Patient) + \eng{be} (conjugué au temps requis + accord) + Participe Passé (PP) du verbe principal [+ \eng{by} Agent].}
\begin{itemize}
    \item L'auxiliaire \eng{be} se conjugue \textbf{exactement au même temps} que le verbe actif original.
    \item Le \textbf{Participe Passé (3\textsuperscript{e} colonne des verbes irréguliers)} est \textbf{invariant} (ne s'accorde pas en anglais).
    \item Les modaux (\eng{can, must, should, might, could}) et semi-modaux (\eng{have to, be going to, used to}) sont suivis de \eng{be} (forme base) + PP.
\end{itemize}
\end{propriete}

\begin{center}
\begin{tabularx}{\linewidth}{|X|X|X|}
\hline
\rowcolor{popblueL} \textbf{Temps / Mode (Actif)} & \textbf{Formule Passive (\eng{be} + PP)} & \textbf{Exemples (to write / to do / to build)} \\
\hline
Present Simple & \eng{am / is / are + PP} & \eng{Letters are written. / The work is done.} \\
\hline
Past Simple & \eng{was / were + PP} & \eng{The letter was written. / The bridge was built.} \\
\hline
Present Continuous & \eng{am / is / are being + PP} & \eng{The book is being written. / The road is being repaired.} \\
\hline
Past Continuous & \eng{was / were being + PP} & \eng{The match was being played. / The hall was being renovated.} \\
\hline
Present Perfect & \eng{has / have been + PP} & \eng{The report has been written. / The fees have been paid.} \\
\hline
Past Perfect & \eng{had been + PP} & \eng{The cocoa had been harvested. / The decision had been taken.} \\
\hline
Future (\eng{will}) & \eng{will be + PP} & \eng{The stadium will be inaugurated. / The law will be applied.} \\
\hline
Future (\eng{be going to}) & \eng{am / is / are going to be + PP} & \eng{The school is going to be renovated.} \\
\hline
Modals (\eng{can, must, should...}) & \eng{Modal + be + PP} & \eng{The form must be filled in. / It can be done.} \\
\hline
\eng{have to} (semi-modal) & \eng{have / has / had to be + PP} & \eng{The visitors have to be picked up.} \\
\hline
Infinitif & \eng{to be + PP} & \eng{The work needs to be done. / He wants to be respected.} \\
\hline
Gérondif & \eng{being + PP} & \eng{He hates being interrupted. / The children like being looked after.} \\
\hline
\end{tabularx}
\end{center}

\cle{Point clé :} Le participe passé \textbf{ne bouge jamais} (invariant), seul \eng{be} conjugue et s'accorde avec le \textbf{nouveau sujet} (ex: \eng{The letter \textbf{was} sent} vs \eng{The letters \textbf{were} sent}).

\begin{attention}[Pièges de forme fréquents]
\begin{enumerate}
    \item \textbf{Oublier \eng{being} aux temps Continuous :} *\eng{The road is repaired} (Present Simple) au lieu de \eng{The road \textbf{is being} repaired} (Present Continuous Passive).
    \item \textbf{Oublier \eng{been} aux temps Perfect :} *\eng{The report has published} au lieu de \eng{The report \textbf{has been} published}.
    \item \textbf{Conjuguer \eng{be} après un modal :} *\eng{The fee must \textbf{is} paid} $\rightarrow$ FAUX. \eng{Must be paid} (base verbale).
    \item \textbf{Participe passé irrégulier incorrect :} *\eng{was did} (PP = \eng{done}), *\eng{was stole} (PP = \eng{stolen}), *\eng{was wrote} (PP = \eng{written}), *\eng{was chose} (PP = \eng{chosen}), *\eng{was broke} (PP = \eng{broken}).
\end{enumerate}
\end{attention}

\begin{exoresolu}[Mise en forme passive : Tous les temps et modaux]
\textbf{Énoncé :} Mettez les verbes entre parenthèses à la \textbf{voix passive} au temps indiqué (ou conservé). Traduisez la phrase obtenue.
\begin{enumerate}
    \item (Present Simple) Cocoa beans \_\_\_\_\_\_ (export) from Kribi port every year.
    \item (Past Continuous) The road \_\_\_\_\_\_ (repair) when the accident happened.
    \item (Present Perfect) The results \_\_\_\_\_\_ just \_\_\_\_\_\_ (publish) on the MINESEC website.
    \item (Modal \eng{must}) Application forms \_\_\_\_\_\_ (submit) before Friday.
    \item (Future \eng{will}) The new stadium \_\_\_\_\_\_ (inaugurate) by the President next month.
    \item (\eng{be going to}) The old market \_\_\_\_\_\_ (demolish) soon.
    \item (Past Simple) The match \_\_\_\_\_\_ (win) by Canon Yaoundé in 2023.
    \item (Present Continuous) English \_\_\_\_\_\_ (teach) in all secondary schools currently.
    \item (Infinitive after \eng{need}) The classroom \_\_\_\_\_\_ (clean) before Monday.
    \item (Gerund after \eng{avoid}) He hates \_\_\_\_\_\_ (interrupt) when he speaks.
\end{enumerate}
\tcblower
\textbf{Solution détaillée :}
\begin{enumerate}
    \item \eng{are exported} (\eng{are} + PP). \textit{« Les fèves de cacao sont exportées du port de Kribi chaque année. »}
    \item \eng{was being repaired} (\eng{was being} + PP). \textit{« La route était en train d'être réparée quand l'accident est survenu. »} (Passif de \eng{were repairing}).
    \item \eng{have just been published} (\eng{have been} + PP). \textit{« Les résultats viennent d'être publiés sur le site du MINESEC. »}
    \item \eng{must be submitted} (\eng{must be} + PP). \textit{« Les formulaires d'inscription doivent être déposés avant vendredi. »}
    \item \eng{will be inaugurated} (\eng{will be} + PP). \textit{« Le nouveau stade sera inauguré par le Président le mois prochain. »}
    \item \eng{is going to be demolished} (\eng{is going to be} + PP). \textit{« Le vieux marché va être démoli bientôt. »}
    \item \eng{was won} (\eng{was} + PP). \textit{« Le match a été gagné par Canon Yaoundé en 2023. »}
    \item \eng{is being taught} (\eng{is being} + PP). \textit{« L'anglais est en train d'être enseigné dans tous les lycées actuellement. »}
    \item \eng{needs to be cleaned} / \eng{needs cleaning} (\eng{need + gérondif passif} ou \eng{need + infinitif passif}). \textit{« La salle de classe doit être nettoyée avant lundi. »}
    \item \eng{being interrupted} (\eng{being} + PP). \textit{« Il déteste être interrompu quand il parle. »}
\end{enumerate}
\textbf{Justifications grammaticales :}
\begin{itemize}
    \item Phrases 2 \& 8 : Continuous $\rightarrow$ \eng{being} obligatoire (\eng{was being repaired}, \eng{is being taught}). Erreur fréquente : *\eng{is taught} (Present Simple) au lieu de \eng{is being taught}.
    \item Phrase 3 : Perfect $\rightarrow$ \eng{been} (PP de \eng{be}) + PP du verbe (\eng{been published}). Ne pas écrire *\eng{have been publish}.
    \item Phrase 4 : Modal $\rightarrow$ pas de conjugaison sur \eng{be} (forme base). *\eng{must are submitted} $\rightarrow$ FAUX.
    \item Phrase 9 : \eng{need} + gérondif (sens passif) = \eng{needs cleaning} OU \eng{needs to be cleaned}.
    \item Phrase 10 : Après préposition (\eng{avoid}) $\rightarrow$ gérondif passif \eng{being interrupted}.
\end{itemize}
\end{exoresolu}

\courssub{Emplois de la voix passive et transformations syntaxiques}

\begin{propriete}[Quand utiliser la voix passive ?]
La voix passive est privilégiée quand :
\begin{enumerate}
    \item \textbf{L'agent est inconnu, évident ou sans importance :} \eng{My phone was stolen at the market.} (On ne sait pas qui / "Someone" est vague).
    \item \textbf{On veut mettre l'accent sur l'action, le résultat ou l'objet concerné :} \eng{The stadium was inaugurated yesterday.} (Le stade est le sujet important, pas le Ministre).
    \item \textbf{Registre formel, journalistique, scientifique, administratif :} \eng{The experiment was conducted under strict conditions.} \eng{The form must be completed in capital letters.}
    \item \textbf{Pour lier les phrases (cohésion) :} Le sujet passif reprend souvent l'objet de la phrase précédente. \eng{The government launched a new policy. It was welcomed by teachers.}
\end{enumerate}
\end{propriete}

\begin{methode}[Transformer une phrase active en passive : Étapes et cas particuliers]
\textbf{Objectif :} Opérer la transformation syntaxique complète sans perte de sens.
\begin{enumerate}
    \item \textbf{Identifier le COD (Objet Direct) de la phrase active.} Il deviendra le \textbf{Nouveau Sujet} de la passive.
    \item \textbf{Conjuguer \eng{be} au même temps} que le verbe actif, \textbf{accordé avec ce nouveau sujet}.
    \item \textbf{Mettre le verbe principal au Participe Passé (PP)}.
    \item \textbf{Gérer l'ancien sujet (Agent) :} Le placer après \eng{by} (Complément d'Agent). \textbf{L'omettre} s'il est générique (\eng{people, someone, they, one, we} vague).
    \item \textbf{Conserver les compléments circonstanciels} (lieu, temps, manière) à leur place habituelle (fin de phrase).
\end{enumerate}
\textbf{Cas particuliers (obligatoires au Bac) :}
\begin{itemize}
    \item \textbf{Verbes à particule (Phrasal Verbs transitifs) :} \textbf{Ne jamais séparer le verbe de sa particule.} La particule reste collée au PP à la fin. \eng{They put off the meeting} $\rightarrow$ \eng{The meeting was put off}. (Pas *\eng{put the meeting off}).
    \item \textbf{Verbes prépositionnels (Prepositional Verbs) :} La préposition \textbf{reste collée au verbe}. \eng{They looked after the children} $\rightarrow$ \eng{The children were looked \textbf{after}}.
    \item \textbf{Verbes à double objet (give, send, offer, tell, ask, show, teach, pay, promise, refuse...) :} \textbf{Deux passives possibles.}
    \begin{enumerate}
        \item \textbf{COD devient Sujet (standard) :} \eng{He gave me a book} $\rightarrow$ \eng{A book was given \textbf{to} me}. (Ajouter \eng{to} devant l'ancien COI).
        \item \textbf{COI devient Sujet (fréquent, plus naturel si personne) :} \eng{He gave me a book} $\rightarrow$ \eng{I was given a book}. \cle{Préférer cette forme si le COI est une personne.}
    \end{enumerate}
    \item \textbf{Verbes suivis de \eng{to}-infinitif (ask, tell, order, advise, allow, forbid, encourage, expect, want...) :} \eng{They asked him to leave} $\rightarrow$ \eng{He was asked \textbf{to leave}}. (L'infinitif \eng{to} est conservé).
\end{itemize}
\textbf{Ordre des mots Passif :} Nouveau Sujet + \eng{be} (conjugué) + PP + (\eng{by} Agent) + Compléments Circonstanciels.
\end{methode}

\begin{exoresolu}[Transformation Active $\rightarrow$ Passive : Cas complexes]
\textbf{Énoncé :} Transformez les phrases actives suivantes en voix passive. Choisissez le sujet le plus naturel (personne si possible pour les verbes à double objet). Omettez l'agent (\eng{by...}) s'il est générique (\eng{people, someone, they}).
\begin{enumerate}
    \item Someone stole my phone at the market in Douala.
    \item They are building a new bridge over the Wouri river.
    \item The teacher gave the students difficult homework.
    \item You must fill in this application form carefully.
    \item People speak English and French in Cameroon.
    \item The police are looking for the suspect in Buea.
    \item They will ask the candidates several questions during the interview.
    \item We have to pick up the visitors at the airport.
    \item The headmaster advised the students to work harder.
    \item They offered the winner a scholarship.
\end{enumerate}
\tcblower
\textbf{Solution détaillée :}
\begin{enumerate}
    \item \eng{My phone was stolen at the market in Douala.} (Agent \eng{someone} omis. \eng{stole} $\rightarrow$ \eng{was stolen}).
    \item \eng{A new bridge is being built over the Wouri river.} (Present Continuous Passif : \eng{are building} $\rightarrow$ \eng{is being built}. Agent \eng{they} omis).
    \item \eng{The students were given difficult homework (by the teacher).} \textbf{Meilleur choix :} COI (\eng{students}) devient Sujet. \eng{were given} + COD (\eng{homework}) conservé. \textit{Alternative :} \eng{Difficult homework was given to the students.}
    \item \eng{This application form must be filled in carefully.} (Verbe à particule \eng{fill in} : particule \eng{in} reste à la fin. Modal \eng{must be} + PP).
    \item \eng{English and French are spoken in Cameroon.} (Agent \eng{people} omis. Verbe \eng{speak} $\rightarrow$ \eng{are spoken}).
    \item \eng{The suspect is being looked for in Buea.} (Verbe prépositionnel \eng{look for} : \eng{for} reste. Present Continuous Passif : \eng{are looking} $\rightarrow$ \eng{is being looked for}).
    \item \eng{The candidates will be asked several questions (during the interview).} (Double objet : COI \eng{candidates} devient Sujet. \eng{will be asked} + COD \eng{several questions}).
    \item \eng{The visitors have to be picked up at the airport.} (Semi-modal \eng{have to} $\rightarrow$ \eng{have to be} + PP. Phrasal verb \eng{pick up} : \eng{up} reste).
    \item \eng{The students were advised to work harder (by the headmaster).} (Verbe + \eng{to}-infinitif : \eng{to work} conservé).
    \item \eng{The winner was offered a scholarship.} (Double objet : COI \eng{winner} devient Sujet. \eng{was offered} + COD \eng{a scholarship}). \textit{Alternative :} \eng{A scholarship was offered to the winner.}
\end{enumerate}
\textbf{Points de vigilance :} Phrase 4 (\eng{filled in} pas *\eng{filled}), Phrase 6 (\eng{looked for} pas *\eng{looked}), Phrase 8 (\eng{picked up} pas *\eng{picked}), Phrase 9 (\eng{advised to work} pas *\eng{advised work}).
\end{exoresolu}

\begin{experience}[Atelier oral : "Passive Voice News Report"]
\textbf{Consigne :} En binôme. L'élève A est journaliste, l'élève B est témoin / officiel.
\textbf{Scénario :} Un incendie s'est déclaré au marché central de Bamenda hier soir.
\begin{enumerate}
    \item \textbf{Préparation (5 min) :} L'élève A prépare 5 questions au passif (\eng{When was the fire noticed? How many shops were destroyed? Were the firemen called immediately? What will be done for the victims?}). L'élève B prépare les réponses (mixte actif/passif).
    \item \textbf{Interview (3 min) :} Jeu de rôle face à la classe.
    \item \textbf{Feedback :} La classe note les formes passives entendues et vérifie l'accord \eng{be} + PP.
\end{enumerate}
\textbf{Objectif :} Automatiser l'oral \eng{be} + PP dans un contexte de communication réel (reportage).
\end{experience}

\courssub{Comparaison Français / Anglais et pièges de traduction}

\begin{attention}[Les 5 interférences majeures Français $\rightarrow$ Anglais (Passif / Actif / Réfléchi)]
\begin{enumerate}
    \item \textbf{Verbes intransitifs français "être + PP" = ACTIF en anglais.}
    \begin{center}
    \begin{tabular}{|l|l|l|}
    \hline
    \rowcolor{poporangeL} \textbf{Français} & \textbf{Anglais (ACTIF)} & \textbf{Erreur fréquente} \\
    \hline
    \fr{Il est arrivé / Il est survenu} & \eng{It \textbf{happened} / \textbf{occurred}} & *\eng{It was happened} \\
    \fr{Il est mort} & \eng{He \textbf{died}} & *\eng{He was died} \\
    \fr{Il est né} & \eng{He \textbf{was born}} (Passif lexicalisé, exception) & -- \\
    \fr{Il est venu / Il est allé} & \eng{He \textbf{came} / \textbf{went}} & *\eng{He was come} \\
    \fr{Il s'est produit} & \eng{It \textbf{happened} / \textbf{took place}} & *\eng{It was produced} \\
    \fr{La loi est entrée en vigueur} & \eng{The law \textbf{came into force} / \textbf{took effect}} & *\eng{was entered into force} \\
    \fr{Il manque} & \eng{It \textbf{lacks} / There \textbf{is a lack of}} & *\eng{It is lacked} \\
    \hline
    \end{tabular}
    \end{center}
    \cle{Règle :} \eng{Happen, occur, die, arrive, come, go, sleep, rain, snow, lack, resemble, appear, seem, consist} $\rightarrow$ \textbf{JAMAIS PASSIF}.
    \item \textbf{"On / Se faire / Se voir" + Infinitif $\neq$ Passif anglais standard.}
    \begin{itemize}
        \item \fr{On m'a dit de partir} $\rightarrow$ \eng{I \textbf{was told to leave}} (Passif + Infinitif \textbf{OK} car \eng{tell} est transitif).
        \item \fr{Je me suis fait voler mon portable} $\rightarrow$ \eng{I \textbf{had my phone stolen}} (Structure \textbf{causative} \eng{have + objet + PP}, \textbf{PAS} *\eng{I was stolen my phone}).
        \item \fr{Il s'est fait arrêter} $\rightarrow$ \eng{He \textbf{got arrested}} (Passif avec \eng{get} pour action subie involontaire / malchance) ou \eng{He was arrested}.
        \item \fr{Elle s'est fait opérer} $\rightarrow$ \eng{She \textbf{had an operation}} / \eng{She \textbf{had her appendix removed}}.
    \end{itemize}
    \item \textbf{Verbes d'état / opinion français = Actif en anglais.}
    \begin{itemize}
        \item \fr{Être d'accord} = \eng{agree} (Actif). *\eng{be agreed} n'existe pas (sauf adjectif : \eng{the agreed price}).
        \item \fr{S'intéresser à} = \eng{be interested \textbf{in}} (Adjectif + préposition, pas passif de \eng{interest}).
        \item \fr{S'occuper de} = \eng{look \textbf{after}} / \eng{take \textbf{care of}} (Verbes prépositionnels, passif possible : \eng{The children were looked after}).
    \end{itemize}
    \item \textbf{Passif impersonnel "Il est..." : deux structures anglaises.}
    \begin{itemize}
        \item \fr{Il est dit que... / On dit que...} $\rightarrow$ \eng{It \textbf{is said} that...} (Passif impersonnel) OU \eng{He \textbf{is said to be}...} (Sujet personnel + infinitif passif). \cle{Très formel, journalistique.}
        \item \fr{Il faut que... / Il est nécessaire de...} $\rightarrow$ \eng{It \textbf{is necessary} to...} / \eng{... \textbf{must be done}} / \eng{... \textbf{has to be done}}.
    \end{itemize}
    \item \textbf{Prépositions après adjectifs participiaux (faux amis) :}
    \begin{center}
    \begin{tabular}{|l|l|l|}
    \hline
    \rowcolor{popgreenL} \textbf{Français} & \textbf{Anglais correct} & \textbf{Erreur (calque)} \\
    \hline
    \fr{Intéressé par} & \eng{interested \textbf{in}} & *\eng{interested \textbf{by} / \textbf{to}} \\
    \fr{Marié à} & \eng{married \textbf{to}} & *\eng{married \textbf{with}} \\
    \fr{Déçu par} & \eng{disappointed \textbf{with / by}} & *\eng{disappointed \textbf{of}} \\
    \fr{Fatigué de} & \eng{tired \textbf{of / from}} & *\eng{tired \textbf{by}} \\
    \fr{Rempli de} & \eng{filled \textbf{with}} & *\eng{filled \textbf{of}} \\
    \fr{Connu pour} & \eng{known \textbf{for}} & *\eng{known \textbf{as}} (sauf "known as a liar") \\
    \hline
    \end{tabular}
    \end{center}
\end{enumerate}
\cle{Règle d'or :} Traduire le \textbf{sens} et la \textbf{valeur syntaxique}, pas la morphologie. Si le verbe anglais est intransitif, pas de passif possible.
\end{attention}

\begin{exoresolu}[Correction d'erreurs de traduction (Français $\rightarrow$ Anglais)]
\textbf{Énoncé :} Les phrases anglaises ci-dessous contiennent des erreurs dues à l'interférence du français. Corrigez-les et expliquez l'erreur.
\begin{enumerate}
    \item The accident \textbf{was happened} yesterday near the Rond-Point Deïdo.
    \item The students \textbf{were agreed} with the teacher's decision.
    \item My brother \textbf{was died} in 2010.
    \item This problem \textbf{will be resolved} by the committee next week. (Phrase correcte ? Justifiez).
    \item The match \textbf{was played} by the players with great courage.
    \item I \textbf{was stolen} my wallet in the bus.
    \item The new law \textbf{is entered into force} last month.
    \item English \textbf{is spoken} in many countries. (Correct ? Justifiez).
    \item The thief \textbf{was arrested} by the police yesterday. (Correct ? Justifiez).
    \item This book \textbf{was written} by Mongo Beti. (Correct ? Justifiez).
\end{enumerate}
\tcblower
\textbf{Solution détaillée et corrections :}
\begin{enumerate}
    \item \textbf{Erreur :} *\eng{was happened}. \fr{Arriver / Se produire} = \eng{happen / occur} (Verbes intransitifs $\rightarrow$ \textbf{JAMAIS de passif}).
    \textbf{Correction :} \eng{The accident \textbf{happened} / \textbf{occurred} yesterday near the Rond-Point Deïdo.}
    \item \textbf{Erreur :} *\eng{were agreed}. \fr{Être d'accord} = \eng{agree} (Verbe intransitif $\rightarrow$ Actif). \eng{Agreed} n'existe qu'en participe passé adjectival (\eng{the agreed price}).
    \textbf{Correction :} \eng{The students \textbf{agreed} with the teacher's decision.}
    \item \textbf{Erreur :} *\eng{was died}. \fr{Mourir} = \eng{die} (Intransitif). \eng{Dead} est adjectif (\eng{He is dead}), \eng{died} est prétérit.
    \textbf{Correction :} \eng{My brother \textbf{died} in 2010.}
    \item \textbf{Correct.} \eng{Resolve} est transitif. Passif \eng{will be resolved} parfait. Agent \eng{by the committee} pertinent ici (responsabilité identifiée).
    \item \textbf{Style lourd / Inélégant.} \eng{Was played by the players} : l'agent (\eng{players}) est l'évidence même (qui d'autre joue ?). La passive est inutile.
    \textbf{Amélioration :} \eng{The players \textbf{played} the match with great courage.} (Voix Active).
    \item \textbf{Erreur majeure :} *\eng{I was stolen my wallet}. \fr{Se faire voler} $\neq$ \eng{be stolen} (qui signifie "être volé" pour l'objet). Structure causative : \eng{have + objet + PP}.
    \textbf{Correction :} \eng{I \textbf{had my wallet stolen} in the bus.} (Ou : \eng{My wallet \textbf{was stolen} in the bus} -- focus sur l'objet).
    \item \textbf{Erreur :} *\eng{is entered into force}. \fr{Entrer en vigueur} = \eng{come into force / take effect} (Verbes intransitifs). Temps : \eng{last month} $\rightarrow$ Past Simple.
    \textbf{Correction :} \eng{The new law \textbf{came into force} / \textbf{took effect} last month.}
    \item \textbf{Correct.} Verbe transitif \eng{speak} (langue). Passif standard pour énoncer un fait général, agent générique omis. Registre neutre/encyclopédique.
    \item \textbf{Correct.} \eng{Arrest} est transitif. Agent \eng{by the police} pertinent (identification de l'autorité). Passif naturel au style journalistique / rapport de police.
    \item \textbf{Correct.} \eng{Write} est transitif. Agent \eng{by Mongo Beti} pertinent (information nouvelle, focus sur l'auteur). Passif standard pour présenter une œuvre.
\end{enumerate}
\end{exoresolu}

\begin{savaistu}[Le passif avec \eng{get} : Registre familier / parlé]
En anglais parlé, \eng{get} remplace souvent \eng{be} pour former la passive, surtout pour des actions \textbf{involontaires, soudaines ou négatives} :
\begin{itemize}
    \item \eng{He \textbf{got fired} last week.} (Il s'est fait virer / a été viré).
    \item \eng{My bike \textbf{got stolen}.} (Mon vélo s'est fait voler).
    \item \eng{Don't \textbf{get confused}.} (Ne te mélange pas les pinceaux).
    \item \eng{They \textbf{got married} in Buea.} (Ils se sont mariés -- passif lexicalisé).
\end{itemize}
\cle{Attention :} \eng{get} + PP est \textbf{interdit à l'écrit formel} (composition, lettre administrative, article de presse). Utilisez \eng{be} + PP.
\end{savaistu}

\courssub{Application guidée et production écrite}

\begin{methode}[Utiliser la voix passive à l'écrit : Registre formel, journalistique et scientifique]
\textbf{Objectif :} Choisir la voix passive stratégiquement pour produire un texte formel (article de presse, rapport, lettre administrative, compte-rendu d'expérience).
\begin{enumerate}
    \item \textbf{Identifier l'intention communicative :}
    \begin{itemize}
        \item \textbf{Focus sur l'action/résultat} (l'agent est inconnu, évident ou secondaire) $\rightarrow$ \cle{Passive}.
        \item \textbf{Focus sur l'agent/responsable} $\rightarrow$ \cle{Active}.
        \item \textbf{Neutralité / Objectivité scientifique / Administrative} $\rightarrow$ \cle{Passive} (\eng{The experiment was conducted...}, \eng{The form must be completed...}).
    \end{itemize}
    \item \textbf{Éviter la surcharge} : Ne pas mettre \textbf{tous} les verbes au passif. Alterner Actif (pour la clarté, le dynamisme) et Passif (pour le focus).
    \item \textbf{Structurer le paragraphe :} \eng{Topic Sentence} (souvent Passif pour annoncer le thème) $\rightarrow$ Développement (Mixte) $\rightarrow$ Conclusion.
    \item \textbf{Vocabulaire de liaison formel (Passifs impersonnels très utiles) :}
    \eng{Furthermore, Moreover, Consequently, Therefore, It is reported that, It is estimated that, It is believed that, It is expected that, It has been decided that.}
\end{enumerate}
\textbf{Modèle "It is said/believed/expected/reported that..."} (Passif de verbe de parole/opinion) : Permet de rapporter une rumeur ou une info officielle sans s'engager personnellement. Ex: \eng{It is expected that the GCE results will be released in August.} $\rightarrow$ \eng{The GCE results are expected to be released in August.} (Sujet personnel + infinitif passif).
\end{methode}

\begin{exoresolu}[Rédaction d'un article de presse (Journalistic Writing)]
\textbf{Énoncé :} Rédigez un court article de presse (environ 60--80 mots) annonçant la rénovation du stade de la Réunification à Douala. Utilisez \textbf{au moins 5 formes passives} (différents temps/modaux). Mots imposés : \eng{renovate, fund, complete, expect, inaugurate}.
\textbf{Consignes :} Ton neutre, objectif. Structure : Titre -- Lead (premier paragraphe) -- Détails.
\tcblower
\textbf{Proposition de rédaction (Modèle corrigé) :}

\begin{textebox}[Article de presse : Stade de la Réunification]
\textbf{Historic Stadium to Undergo Major Renovation}

The iconic Stade de la Réunification in Douala \textbf{is being renovated} to meet international standards. The project, \textbf{which is funded} by the Ministry of Sports, \textbf{was launched} last Monday. According to officials, the first phase \textbf{will be completed} by December. It \textbf{is expected} that the stadium \textbf{will be inaugurated} before the next AFCON qualifiers. Local residents hope the works \textbf{will not disrupt} daily traffic excessively.
\end{textebox}

\textbf{Analyse des formes passives utilisées (6 passifs + 1 actif choisi) :}
\begin{enumerate}
    \item \eng{is being renovated} (Present Continuous Passif -- Focus sur l'action en cours).
    \item \eng{is funded} (Present Simple Passif -- Fait établi, source du financement).
    \item \eng{was launched} (Past Simple Passif -- Événement passé, agent omis).
    \item \eng{will be completed} (Future Passif -- Prévision officielle).
    \item \eng{is expected that... will be inaugurated} (Passif impersonnel \eng{It is expected} + Future Passif \eng{will be inaugurated} -- Style journalistique standard).
    \item \eng{will not disrupt} (Active volontaire pour le dernier verbe : les travaux sont l'agent, on garde l'actif pour varier). \textit{Note :} Si on voulait passif : \eng{daily traffic will not be disrupted}.
\end{enumerate}
\textbf{Conclusion :} La passive permet de mettre en avant le stade (patrimoine), le projet, les phases, sans répéter "The Ministry", "The workers", "The contractors".
\end{exoresolu}

\begin{exercice}[Entraînement complet : Voix Active / Passive]
\textbf{Exercice 1 : Identification} -- Soulignez le verbe et cochez Active (A) ou Passive (P).
\begin{enumerate}
    \item The new curriculum has been implemented in all schools. \hfill $\square$ A \quad $\square$ P
    \item The students are writing their exams this week. \hfill $\square$ A \quad $\square$ P
    \item The bridge over the Wouri was built by a Chinese firm. \hfill $\square$ A \quad $\square$ P
    \item Cocoa prices have increased recently. \hfill $\square$ A \quad $\square$ P
    \item The results will be published on the official website. \hfill $\square$ A \quad $\square$ P
\end{enumerate}

\textbf{Exercice 2 : Mise en forme} -- Mettez le verbe à la voix passive au temps convenable.
\begin{enumerate}
    \item The letters (send) \_\_\_\_\_\_ yesterday by express mail. (Past Simple)
    \item The new hospital (equip) \_\_\_\_\_\_ with modern machines currently. (Present Continuous)
    \item All candidates (inform) \_\_\_\_\_\_ of the date change by SMS. (Present Perfect)
    \item The contract (sign) \_\_\_\_\_\_ by both parties tomorrow. (Future \eng{will})
    \item This form (fill in) \_\_\_\_\_\_ in capital letters. (Modal \eng{must})
    \item The old building (demolish) \_\_\_\_\_\_ next month. (\eng{be going to})
    \item The thief (catch) \_\_\_\_\_\_ by the police last night. (Past Simple)
    \item English (teach) \_\_\_\_\_\_ in primary schools since 2018. (Present Perfect)
\end{enumerate}

\textbf{Exercice 3 : Transformation Active $\rightarrow$ Passive} -- Transformez. Omettez l'agent si générique.
\begin{enumerate}
    \item Someone has stolen the goalkeeper's gloves.
    \item They are repairing the main road to Bafoussam.
    \item The teacher gave each student a textbook.
    \item You must switch off your phones during the exam.
    \item People speak French and English in Cameroon.
    \item The committee will choose the best project.
    \item We have to look after the young children.
    \item The headmaster asked the prefects to organize the party.
\end{enumerate}

\textbf{Exercice 4 : Correction d'erreurs} -- Corrigez les phrases fautives (grammaire ou traduction).
\begin{enumerate}
    \item The meeting was happened in the conference hall.
    \item I was agreed with his opinion.
    \item The baby was borned at 3 a.m.
    \item My father was died in 2005.
    \item The match was played by the players very well.
    \item I was stolen my bag at the market.
    \item The law is entered into force yesterday.
    \item The work was did by the students.
    \item The money was stealed from the drawer.
    \item The report was wrote by the secretary.
\end{enumerate}

\textbf{Exercice 5 : Production écrite} -- Rédigez un court rapport (80--100 mots) sur la cérémonie de remise des prix de votre établissement. Utilisez au moins 6 formes passives variées (temps, modaux, infinitif, gérondif). Mots suggérés : \eng{award, congratulate, encourage, organize, attend, deliver, announce}.
\end{exercice}

\begin{corrige}[Exercice 1 : Identification]
\begin{enumerate}
    \item \textbf{P} (\eng{has been implemented} -- Present Perfect Passive).
    \item \textbf{A} (\eng{are writing} -- Present Continuous Active).
    \item \textbf{P} (\eng{was built} -- Past Simple Passive + Agent).
    \item \textbf{A} (\eng{have increased} -- Present Perfect Active, \eng{increase} intransitif ici).
    \item \textbf{P} (\eng{will be published} -- Future Passive).
\end{enumerate}
\end{corrige}

\begin{corrige}[Exercice 2 : Mise en forme]
\begin{enumerate}
    \item \eng{were sent}
    \item \eng{is being equipped}
    \item \eng{have been informed}
    \item \eng{will be signed}
    \item \eng{must be filled in}
    \item \eng{is going to be demolished}
    \item \eng{was caught}
    \item \eng{has been taught}
\end{enumerate}
\end{corrige}

\begin{corrige}[Exercice 3 : Transformation]
\begin{enumerate}
    \item \eng{The goalkeeper's gloves have been stolen.}
    \item \eng{The main road to Bafoussam is being repaired.}
    \item \eng{Each student was given a textbook.} (COI $\rightarrow$ Sujet) \textit{OU} \eng{A textbook was given to each student.}
    \item \eng{Phones must be switched off during the exam.} (Phrasal verb \eng{switch off}).
    \item \eng{French and English are spoken in Cameroon.}
    \item \eng{The best project will be chosen (by the committee).}
    \item \eng{The young children have to be looked after.} (Verbe prépositionnel \eng{look after}).
    \item \eng{The prefects were asked to organize the party.} (Verbe + \eng{to}-infinitif).
\end{enumerate}
\end{corrige}

\begin{corrige}[Exercice 4 : Correction d'erreurs]
\begin{enumerate}
    \item \eng{The meeting \textbf{took place} / \textbf{happened} in the conference hall.} (\eng{happen} intransitif).
    \item \eng{I \textbf{agreed} with his opinion.} (\eng{agree} intransitif).
    \item \eng{The baby \textbf{was born} at 3 a.m.} (\eng{born} = PP lexicalisé, pas \eng{borned}).
    \item \eng{My father \textbf{died} in 2005.} (\eng{die} intransitif).
    \item \eng{The players \textbf{played} the match very well.} (Actif préférable, agent évident).
    \item \eng{I \textbf{had my bag stolen} at the market.} (Causative \eng{have + objet + PP}). Ou \eng{My bag was stolen at the market.}
    \item \eng{The law \textbf{came into force} / \textbf{took effect} yesterday.} (\eng{come into force} intransitif + Past Simple).
    \item \eng{The work \textbf{was done} by the students.} (PP de \eng{do} = \eng{done}).
    \item \eng{The money \textbf{was stolen} from the drawer.} (PP de \eng{steal} = \eng{stolen}).
    \item \eng{The report \textbf{was written} by the secretary.} (PP de \eng{write} = \eng{written}).
\end{enumerate}
\end{corrige}

\begin{corrige}[Exercice 5 : Production écrite (Proposition de rapport)]
\begin{textebox}[Rapport : Cérémonie de remise des prix -- Lycée de Buea]
\textbf{Prize-Giving Ceremony Held at Buea High School}

The annual prize-giving ceremony \textbf{was organized} last Saturday in the school hall. The event \textbf{was attended} by the Regional Delegate of Secondary Education, parents, and former students. Deserving students \textbf{were awarded} prizes for academic excellence and discipline. In his speech, the Principal \textbf{congratulated} the winners and \textbf{encouraged} others to work harder. It \textbf{was announced} that a new library \textbf{will be built} next year. The ceremony \textbf{was closed} with a cultural display. Everyone agreed it \textbf{had been} a memorable day.
\end{textebox}
\textbf{Formes passives utilisées (7) :} \eng{was organized, was attended, were awarded, was announced, will be built, was closed, had been}. Verbes actifs conservés pour le dynamisme : \eng{congratulated, encouraged, agreed}.
\end{corrige}

\begin{aretenir}[Synthèse : Voix Active vs Passive -- Check-list Bac]
\begin{itemize}
    \item \textbf{Formule :} Sujet + \eng{be} (temps + accord) + PP [+ \eng{by} Agent].
    \item \textbf{Marqueurs de temps :} \eng{being} (Continu), \eng{been} (Perfect), \eng{be} (Modaux/Infinitif/Gérondif).
    \item \textbf{Verbes intransitifs = JAMAIS passif :} \eng{happen, occur, die, arrive, come, go, sleep, rain, lack, resemble, appear, seem}.
    \item \textbf{Verbes d'état = Actif :} \eng{be, look, feel, seem, become, agree, consist}.
    \item \textbf{Particules / Prépositions :} Restent collées au PP (\eng{put off, looked after, filled in, picked up}).
    \item \textbf{Double objet :} COI $\rightarrow$ Sujet (naturel) \eng{I was given a book} OU COD $\rightarrow$ Sujet \eng{A book was given to me}.
    \item \textbf{Français "Se faire" :} \eng{have + objet + PP} (\eng{I had my phone repaired}), \textbf{PAS} *\eng{I was repaired my phone}.
    \item \textbf{Passif impersonnel :} \eng{It is said/reported/expected that...} / \eng{He is said to be...}
    \item \textbf{Registre :} Passif = formel, objectif, focus sur résultat. Actif = clair, direct, focus sur agent.
    \item \textbf{Vérification finale :} Sujet pluriel $\rightarrow$ \eng{were/are/have been} ; Sujet singulier $\rightarrow$ \eng{was/is/has been}. PP irrégulier correct ?
\end{itemize}
\end{aretenir}

\newpage

\coursec{Exercices}

\courssub{Je vérifie mes connaissances}

\begin{exercice}[QCM : Forme passive correcte]
Choisissez la forme passive correcte pour compléter chaque phrase.
\begin{enumerate}
    \item English \_\_\_\_\_ in many countries around the world.
    \begin{enumerate}
        \item is spoke
        \item is spoken
        \item is speaking
        \item has spoke
    \end{enumerate}
    \item The new stadium \_\_\_\_\_ next year.
    \begin{enumerate}
        \item will build
        \item will be built
        \item is building
        \item has built
    \end{enumerate}
    \item The match \_\_\_\_\_ live on CRTV when the power cut happened.
    \begin{enumerate}
        \item was broadcasted
        \item was broadcasting
        \item was being broadcast
        \item has been broadcast
    \end{enumerate}
    \item My phone \_\_\_\_\_ yesterday at the market.
    \begin{enumerate}
        \item was stolen
        \item was stole
        \item has stolen
        \item had stole
    \end{enumerate}
    \item The cocoa beans \_\_\_\_\_ to Europe every month.
    \begin{enumerate}
        \item are exported
        \item are exporting
        \item have exported
        \item export
    \end{enumerate}
\end{enumerate}
\end{exercice}

\begin{exercice}[Vrai ou Faux justifié]
Déterminez si les affirmations suivantes sont vraies (V) ou fausses (F). Justifiez votre réponse en français en vous basant sur la règle de formation de la voix passive.
\begin{enumerate}
    \item La phrase "The accident was happened yesterday" est correcte car \eng{happen} est un verbe d'action.
    \item Dans la phrase "The letter was written by John", l'agent est introduit par la préposition \eng{by}.
    \item On peut transformer n'importe quelle phrase active en phrase passive en anglais.
    \item La forme passive du \eng{Present Perfect} de \eng{write} est \eng{has been written}.
    \item "English is spoke here" est la forme passive correcte du \eng{Present Simple}.
\end{enumerate}
\end{exercice}

\begin{exercice}[Vocabulaire : Mots-outils de la voix passive]
Complétez le tableau suivant en associant chaque temps verbal à sa structure passive (auxiliaire \eng{be} + participe passé) et donnez un exemple original pour chacun.
\begin{center}
\begin{tabularx}{\linewidth}{|X|X|X|}
\hline
\rowcolor{popblueL} \textbf{Temps} & \textbf{Structure passive (auxiliaire + V3)} & \textbf{Exemple personnel (contexte camerounais si possible)} \\
\hline
Present Simple &  &  \\
\hline
Past Simple &  &  \\
\hline
Present Continuous &  &  \\
\hline
Present Perfect &  &  \\
\hline
Future Simple (will) &  &  \\
\hline
Modal (must/can) &  &  \\
\hline
\end{tabularx}
\end{center}
\end{exercice}

\begin{exercice}[Conjugaison à compléter : Formes passives]
Mettez le verbe entre parenthèses à la voix passive au temps indiqué. Attention aux verbes irréguliers.
\begin{enumerate}
    \item The Indomitable Lions \_\_\_\_\_ (defeat) in the final last Sunday. (Past Simple)
    \item A new bridge over the Wouri River \_\_\_\_\_ (build) at the moment. (Present Continuous)
    \item All the examination papers \_\_\_\_\_ (mark) by next Friday. (Future Simple)
    \item French and English \_\_\_\_\_ (speak) officially in Cameroon. (Present Simple)
    \item The thief \_\_\_\_\_ (catch) by the police just now. (Present Perfect)
    \item The meeting \_\_\_\_\_ (postpone) because of the heavy rain. (Past Simple)
    \item This novel \_\_\_\_\_ (write) by a famous Nigerian author. (Present Simple)
    \item The results \_\_\_\_\_ (announce) tomorrow morning. (Future Simple / \eng{be going to})
\end{enumerate}
\end{exercice}

\courssub{Je m'entraîne}

\begin{exercice}[Transformation : Active $\rightarrow$ Passive]
Transformez les phrases suivantes de la voix active à la voix passive. Conservez le même temps. N'incluez l'agent (\eng{by}...) que s'il est pertinent (sujet précis, non indéfini).
\begin{enumerate}
    \item The headmaster punished the boys for making noise.
    \item Someone has broken the window of the staff room.
    \item They will broadcast the AFCON match live on national television.
    \item The students are cleaning the classroom right now.
    \item A famous journalist wrote this article about the cocoa industry.
    \item The government has announced new measures for youth employment.
    \item You must submit your application before June 30th.
    \item They sell fresh fish and plantains at the Mokolo market every morning.
\end{enumerate}
\end{exercice}

\begin{exercice}[Traduction : Français $\rightarrow$ Anglais (Voix passive)]
Traduisez les phrases suivantes en anglais en utilisant impérativement la voix passive.
\begin{enumerate}
    \item On a volé mon téléphone portable hier dans le bendskin.
    \item On vend des journaux et des magazines ici.
    \item Ce pont a été construit en 2010 par une entreprise chinoise.
    \item La finale de la Coupe du Cameroun sera jouée au stade d'Olembé.
    \item Les résultats du Baccalauréat ont été proclamés la semaine dernière.
    \item On ne doit pas jeter les déchets par terre.
    \item Un nouveau lycée est en train d'être construit à Buea.
    \item L'anglais est parlé dans de nombreux pays du Commonwealth.
\end{enumerate}
\end{exercice}

\begin{exercice}[Texte à trous : Article de presse]
Complétez le texte suivant avec la forme passive correcte des verbes entre parenthèses.
\begin{textebox}[Cameroon Hosts International Tech Summit]
Yaoundé, March 15 -- The 5th International Technology Summit \_\_\_\_\_\_\_\_\_ (open) by the Minister of Posts and Telecommunications yesterday morning. Over 500 delegates from 20 countries \_\_\_\_\_\_\_\_\_ (expect) to attend the three-day event. The main conference hall \_\_\_\_\_\_\_\_\_ (renovate) specially for this occasion last month.

During the opening ceremony, it \_\_\_\_\_\_\_\_\_ (announce) that a new digital innovation hub \_\_\_\_\_\_\_\_\_ (create) in Douala next year. Several start-ups \_\_\_\_\_\_\_\_\_ (select) to present their projects. The best project \_\_\_\_\_\_\_\_\_ (award) a grant of 10 million FCFA on the final day.

Organisers say that special attention \_\_\_\_\_\_\_\_\_ (give) to cybersecurity and artificial intelligence this year. All sessions \_\_\_\_\_\_\_\_\_ (stream) live on the official website.
\end{textebox}
\end{exercice}

\begin{exercice}[Correction d'erreurs : Pièges fréquents]
Les phrases suivantes contiennent des erreurs typiques des francophones. Identifiez l'erreur, expliquez-la brièvement en français et réécrivez la phrase correcte.
\begin{enumerate}
    \item *The house was build in 2000.
    \item *English is spoke in this shop.
    \item *The accident was happened near the roundabout.
    \item *My bag has been stole in the taxi.
    \item *The match will played tomorrow afternoon.
    \item *This book was wrote by Mongo Beti.
    \item *The letters are send every Monday.
    \item *He was died in a car crash.
\end{enumerate}
\end{exercice}

\courssub{Je me prépare au Bac}

\begin{exercice}[Type Bac : Compréhension et Grammaire -- Texte original]
Lisez attentivement le texte ci-dessous et répondez aux questions qui suivent.
\begin{textebox}[Digital Cameroon: A New Era for Youth]
The "Digital Cameroon 2025" strategy was launched by the Head of State three years ago. Since then, significant investments have been made in fibre optic infrastructure across the ten regions. Thousands of young Cameroonians have been trained in coding, digital marketing, and data analysis through government-sponsored programmes.

Last week, a major tech hub was inaugurated in Buea. It was named "Silicon Mountain Hub" by the Minister of Youth Affairs. The centre is equipped with high-speed internet and modern workstations. It is hoped that local start-ups will be incubated there and that innovative solutions for agriculture and health will be developed.

However, challenges remain. Electricity cuts are still experienced in many rural areas. Moreover, the cost of internet data is considered too high by many students. If these issues are not addressed quickly, the digital divide between urban and rural youth will be widened.
\end{textebox}

\textbf{Questions de compréhension (en anglais) :}
\begin{enumerate}
    \item Who launched the "Digital Cameroon 2025" strategy?
    \item What has been done for young Cameroonians since the launch?
    \item What was inaugurated in Buea last week?
    \item What two main challenges are mentioned in the last paragraph?
\end{enumerate}

\textbf{Questions de langue :}
\begin{enumerate}
    \setcounter{enumi}{4}
    \item Identify \textbf{five} verbs in the passive voice in the text and give their tense (e.g., \eng{was launched -- Past Simple}).
    \item Rewrite the following sentence from the text in the active voice: "It is hoped that local start-ups will be incubated there." (Start with \eng{We hope...})
    \item Transform this active sentence into the passive voice: "The government sponsors training programmes for youth."
    \item Find in the text the passive equivalent of: "People experience electricity cuts..." (Paragraph 3).
\end{enumerate}
\end{exercice}

\begin{exercice}[Type Bac : Traduction Thématique]
Traduisez le paragraphe suivant en anglais. Utilisez la voix passive chaque fois que c'est naturel ou nécessaire (au moins 4 occurrences).
\begin{textebox}[Source française]
La Coupe d'Afrique des Nations a été organisée avec succès au Cameroun en 2022. Les stades ont été rénovés et de nouvelles routes ont été construites pour l'événement. Des milliers de touristes ont été accueillis dans les villes de Yaoundé, Douala, Buea et Limbé. On dit que cet événement a boosté l'économie locale. Cependant, certains projets n'ont pas été terminés à temps. Il faut que des efforts soient faits pour entretenir ces infrastructures.
\end{textebox}
\end{exercice}

\begin{exercice}[Type Bac : Sujet de Composition -- Expression Écrite]
\textbf{Consigne :} Vous êtes rédacteur pour le journal scolaire \eng{"The High School Echo"}. Écrivez un article d'environ \textbf{120 à 150 mots} sur un événement récent qui s'est déroulé dans votre établissement (journée culturelle, finale sportive, visite d'une personnalité, remise de prix, plantation d'arbres, etc.).

\textbf{Contraintes obligatoires :}
\begin{itemize}
    \item Utilisez la voix passive \textbf{au moins 4 fois} (soulignez-les dans votre copie).
    \item Utilisez au moins 3 temps différents (ex: Past Simple, Present Perfect, Present Simple, Future).
    \item Structurez votre article : Titre accrocheur, introduction (quoi, quand, où), développement (déroulement, détails), conclusion (impact, avis).
    \item Vocabulaire soigné : \eng{inaugurated, attended, awarded, organized, performed, planted, delivered, celebrated}.
\end{itemize}
\end{exercice}

\newpage

\coursec{Corrigés des exercices}

\begin{corrige}[Exercice 1 — QCM : Forme passive correcte]
\begin{enumerate}
    \item \textbf{b) is spoken} — \eng{English is spoken in many countries around the world.} (L'anglais est parlé dans de nombreux pays du monde.) Présent simple passif : \cle{is} + participe passé \cle{spoken} (verbe irrégulier \eng{speak -- spoke -- spoken}).
    \item \textbf{b) will be built} — \eng{The new stadium will be built next year.} (Le nouveau stade sera construit l'année prochaine.) Futur passif : \cle{will be} + participe passé \cle{built} (\eng{build -- built -- built}).
    \item \textbf{c) was being broadcast} — \eng{The match was being broadcast live on CRTV when the power cut happened.} (Le match était en train d'être diffusé en direct sur la CRTV quand la coupure de courant a eu lieu.) Past Continuous passif : \cle{was being} + participe passé ; \eng{broadcast} est invariable (\eng{broadcast -- broadcast -- broadcast}), donc \eng{broadcasted} est fautif.
    \item \textbf{a) was stolen} — \eng{My phone was stolen yesterday at the market.} (Mon téléphone a été volé hier au marché.) Past Simple passif ; participe passé de \eng{steal} : \cle{stolen} (\eng{steal -- stole -- stolen}).
    \item \textbf{a) are exported} — \eng{The cocoa beans are exported to Europe every month.} (Les fèves de cacao sont exportées vers l'Europe chaque mois.) Présent simple passif, sujet pluriel : \cle{are} + \cle{exported}.
\end{enumerate}
\textbf{Point de méthode :} la voix passive exige toujours l'auxiliaire \eng{be} conjugué + le \textbf{participe passé} (3\textsuperscript{e} colonne), jamais le prétérit (2\textsuperscript{e} colonne).
\end{corrige}

\begin{corrige}[Exercice 2 — Vrai ou Faux justifié]
\begin{enumerate}
    \item \textbf{FAUX.} \eng{Happen} est un verbe \textbf{intransitif} : il n'a pas de complément d'objet direct, donc il ne peut pas se mettre au passif. On dit : \eng{The accident happened yesterday.} (L'accident s'est produit hier.)
    \item \textbf{VRAI.} Le complément d'agent (celui qui fait l'action) est introduit par \eng{by} : \eng{The letter was written by John.} (La lettre a été écrite par John.)
    \item \textbf{FAUX.} Seuls les verbes \textbf{transitifs} (suivis d'un COD) peuvent se conjuguer au passif. Les verbes intransitifs (\eng{happen, arrive, die, sleep, go}) ne le peuvent pas.
    \item \textbf{VRAI.} Present Perfect passif : \cle{have/has been} + participe passé. Avec un sujet singulier : \eng{It has been written.} (\eng{write -- wrote -- written}.)
    \item \textbf{FAUX.} Après \eng{is}, il faut le participe passé \eng{spoken}, pas le prétérit \eng{spoke}. Forme correcte : \eng{English is spoken here.} (On parle anglais ici.)
\end{enumerate}
\end{corrige}

\begin{corrige}[Exercice 3 — Vocabulaire : Mots-outils de la voix passive]
Tableau complété (les exemples sont des propositions ; toute phrase correcte et originale est acceptée) :
\begin{center}
\begin{tabularx}{\linewidth}{|X|X|X|}
\hline
\rowcolor{popblueL} \textbf{Temps} & \textbf{Structure passive} & \textbf{Exemple} \\
\hline
Present Simple & am / is / are + V3 & \eng{Newspapers are sold at the post office.} \\
\hline
Past Simple & was / were + V3 & \eng{The trophy was won by our school team.} \\
\hline
Present Continuous & am / is / are being + V3 & \eng{A new road is being built in Bamenda.} \\
\hline
Present Perfect & have / has been + V3 & \eng{The results have been published online.} \\
\hline
Future Simple (will) & will be + V3 & \eng{The exams will be corrected next week.} \\
\hline
Modal (must/can) & must / can + be + V3 & \eng{The rules must be respected by everyone.} \\
\hline
\end{tabularx}
\end{center}
\textbf{Traductions :} Les journaux sont vendus à la poste. / Le trophée a été remporté par l'équipe de notre lycée. / Une nouvelle route est en train d'être construite à Bamenda. / Les résultats ont été publiés en ligne. / Les examens seront corrigés la semaine prochaine. / Les règles doivent être respectées par tous.
\end{corrige}

\begin{attention}[Point de vigilance]
Après un modal (\eng{must, can, should, may}), l'auxiliaire \eng{be} reste \textbf{invariable} : \eng{must be done}, jamais \eng{must is done} ni \eng{must been done}.
\end{attention}


\begin{corrige}[Exercice 4 — Conjugaison à compléter : Formes passives]
\begin{enumerate}
    \item \eng{The Indomitable Lions \textbf{were defeated} in the final last Sunday.} (Les Lions Indomptables ont été battus en finale dimanche dernier.) Sujet pluriel $\rightarrow$ \eng{were} + participe passé régulier \eng{defeated}.
    \item \eng{A new bridge over the Wouri River \textbf{is being built} at the moment.} (Un nouveau pont sur la Wouri est en train d'être construit en ce moment.) Present Continuous passif : \eng{is being} + \eng{built} (\eng{build -- built -- built}).
    \item \eng{All the examination papers \textbf{will be marked} by next Friday.} (Toutes les copies d'examen seront corrigées d'ici vendredi prochain.)
    \item \eng{French and English \textbf{are spoken} officially in Cameroon.} (Le français et l'anglais sont parlés officiellement au Cameroun.) Sujet pluriel ; participe passé \eng{spoken}.
    \item \eng{The thief \textbf{has just been caught} by the police.} (Le voleur vient d'être attrapé par la police.) Present Perfect passif ; \eng{catch -- caught -- caught}. \eng{Just} appelle le Present Perfect.
    \item \eng{The meeting \textbf{was postponed} because of the heavy rain.} (La réunion a été reportée à cause des fortes pluies.)
    \item \eng{This novel \textbf{is written} by a famous Nigerian author.} (Ce roman est écrit par un célèbre auteur nigérian.) \eng{write -- wrote -- written}.
    \item \eng{The results \textbf{are going to be announced} tomorrow morning.} (Les résultats vont être annoncés demain matin.) Futur avec \eng{be going to} : \eng{are going to be} + participe passé.
\end{enumerate}
\end{corrige}

\begin{corrige}[Exercice 5 — Transformation : Active $\rightarrow$ Passive]
\begin{enumerate}
    \item \eng{The boys were punished by the headmaster for making noise.} (Les garçons ont été punis par le proviseur pour avoir fait du bruit.) Agent pertinent (sujet précis) $\rightarrow$ on garde \eng{by the headmaster}.
    \item \eng{The window of the staff room has been broken.} (La fenêtre de la salle des professeurs a été cassée.) \eng{Someone} est indéfini $\rightarrow$ on supprime l'agent.
    \item \eng{The AFCON match will be broadcast live on national television.} (Le match de la CAN sera diffusé en direct à la télévision nationale.) \eng{They} indéfini $\rightarrow$ pas d'agent ; \eng{broadcast} invariable.
    \item \eng{The classroom is being cleaned right now.} (La salle de classe est en train d'être nettoyée en ce moment.) Present Continuous passif ; agent \eng{the students} omis car évident.
    \item \eng{This article about the cocoa industry was written by a famous journalist.} (Cet article sur l'industrie du cacao a été écrit par un journaliste célèbre.) Agent précis $\rightarrow$ conservé.
    \item \eng{New measures for youth employment have been announced by the government.} (De nouvelles mesures pour l'emploi des jeunes ont été annoncées par le gouvernement.) Sujet pluriel $\rightarrow$ \eng{have been}.
    \item \eng{Your application must be submitted before June 30th.} (Votre candidature doit être déposée avant le 30 juin.) Passif avec modal : \cle{must be} + participe passé.
    \item \eng{Fresh fish and plantains are sold at the Mokolo market every morning.} (Du poisson frais et des plantains sont vendus au marché Mokolo chaque matin.) \eng{They} indéfini $\rightarrow$ supprimé.
\end{enumerate}
\textbf{Rappel de méthode :} le COD de la phrase active devient le sujet de la phrase passive ; le verbe prend la forme \eng{be} + participe passé au même temps ; l'ancien sujet devient complément d'agent avec \eng{by} (ou disparaît s'il est vague : \eng{someone, they, people}).
\end{corrige}

\begin{corrige}[Exercice 6 — Traduction : Français $\rightarrow$ Anglais (Voix passive)]
\begin{enumerate}
    \item \eng{My mobile phone was stolen yesterday on the bendskin.} (ou \eng{on the motorbike taxi})
    \item \eng{Newspapers and magazines are sold here.}
    \item \eng{This bridge was built in 2010 by a Chinese company.}
    \item \eng{The final of the Cameroon Cup will be played at the Olembe Stadium.}
    \item \eng{The Baccalaureate results were proclaimed last week.}
    \item \eng{Rubbish must not be thrown on the ground.} (ou \eng{Litter must not be dropped on the ground.})
    \item \eng{A new high school is being built in Buea.}
    \item \eng{English is spoken in many Commonwealth countries.}
\end{enumerate}
\end{corrige}

\begin{attention}[Piège du « on » français]
Le français « on » se rend très souvent par la voix passive en anglais : « on a volé mon téléphone » $\rightarrow$ \eng{my phone was stolen} (et non \eng{someone has stolen...}, lourd et moins naturel). De même « on dit que... » $\rightarrow$ \eng{it is said that...}.
\end{attention}


\begin{corrige}[Exercice 7 — Texte à trous : Article de presse]
Texte complété :
\begin{enumerate}
    \item \eng{was opened} — Past Simple passif (\eng{yesterday morning}) ; \eng{The 5th International Technology Summit was opened by the Minister...}
    \item \eng{were expected} — Past Simple passif, sujet pluriel (\eng{over 500 delegates}).
    \item \eng{was renovated} — Past Simple passif (\eng{last month}).
    \item \eng{was announced} — tournure impersonnelle \eng{it was announced that...} (il a été annoncé que...).
    \item \eng{will be created} — futur passif (\eng{next year}).
    \item \eng{were selected} — Past Simple passif, sujet pluriel (\eng{several start-ups}).
    \item \eng{will be awarded} — futur passif ; \eng{award a grant} = attribuer une bourse.
    \item \eng{is being given} / \eng{will be given} — les deux sont acceptables ; \eng{this year} suggère le Present Continuous passif (\eng{is being given}) ou le futur.
    \item \eng{will be streamed} — futur passif ; \eng{stream} = diffuser en ligne.
\end{enumerate}
\textbf{Traduction du début :} « Yaoundé, 15 mars — Le 5\textsuperscript{e} Sommet international des technologies a été ouvert par le ministre des Postes et Télécommunications hier matin. Plus de 500 délégués venus de 20 pays étaient attendus pour cet événement de trois jours. »
\end{corrige}

\begin{corrige}[Exercice 8 — Correction d'erreurs : Pièges fréquents]
\begin{enumerate}
    \item \eng{The house was \textbf{built} in 2000.} — Erreur : le participe passé \textbf{built} a été remplacé par la base verbale \eng{build}. Après \eng{was}, il faut le participe passé ; ici \eng{build -- built -- built} (prétérit et participe identiques).
    \item \eng{English is \textbf{spoken} in this shop.} — Erreur : confusion prétérit/participe passé. \eng{speak -- spoke -- spoken} ; au passif, toujours la 3\textsuperscript{e} colonne.
    \item \eng{The accident \textbf{happened} near the roundabout.} — Erreur : \eng{happen} est intransitif, il n'a pas de passif. Le français « s'est produit » tente à la voix pronominale, mais l'anglais reste à l'actif.
    \item \eng{My bag has been \textbf{stolen} in the taxi.} — Erreur : participe passé de \eng{steal} = \eng{stolen}, pas \eng{stole} (prétérit).
    \item \eng{The match will \textbf{be played} tomorrow afternoon.} — Erreur : auxiliaire \eng{be} oublié après \eng{will}. Futur passif : \eng{will be} + participe passé.
    \item \eng{This book was \textbf{written} by Mongo Beti.} — Erreur : prétérit \eng{wrote} utilisé à la place du participe passé \eng{written}.
    \item \eng{The letters are \textbf{sent} every Monday.} — Erreur : base verbale \eng{send} au lieu du participe passé \eng{sent}. \eng{send -- sent -- sent}.
    \item \eng{He \textbf{died} in a car crash.} — Erreur : \eng{die} est intransitif (comme « mourir ») ; pas de passif possible. Ne pas confondre avec \eng{he was killed} (il a été tué), qui est correct car \eng{kill} est transitif.
\end{enumerate}
\end{corrige}

\begin{corrige}[Exercice 9 — Type Bac : Compréhension et Grammaire]
\textbf{Barème indicatif (10 pts) :} compréhension 4 pts (1 pt par question) ; langue 6 pts (Q5 : 2,5 pts ; Q6 : 1 pt ; Q7 : 1,5 pts ; Q8 : 1 pt).

\textbf{Questions de compréhension — réponses attendues :}
\begin{enumerate}
    \item \eng{The "Digital Cameroon 2025" strategy was launched by the Head of State.} (Le chef de l'État.)
    \item \eng{They have been trained in coding, digital marketing and data analysis through government-sponsored programmes.} (Ils ont été formés au codage, au marketing numérique et à l'analyse de données.)
    \item \eng{A major tech hub, named "Silicon Mountain Hub", was inaugurated in Buea last week.} (Un grand pôle technologique.)
    \item \eng{The two main challenges are electricity cuts in rural areas and the high cost of internet data.} (Les coupures d'électricité et le coût élevé des données internet.)
\end{enumerate}

\textbf{Questions de langue :}
\begin{enumerate}
    \setcounter{enumi}{4}
    \item Cinq verbes au passif (au choix) : \eng{was launched} (Past Simple) ; \eng{have been made} (Present Perfect) ; \eng{have been trained} (Present Perfect) ; \eng{was inaugurated} (Past Simple) ; \eng{was named} (Past Simple) ; \eng{is equipped} (Present Simple) ; \eng{is hoped} (Present Simple) ; \eng{will be incubated / will be developed / will be widened} (Future Simple) ; \eng{are experienced} (Present Simple) ; \eng{is considered} (Present Simple) ; \eng{are not addressed} (Present Simple).
    \item \eng{We hope that local start-ups will be incubated there.} (Nous espérons que des start-ups locales y seront incubées.) La tournure impersonnelle \eng{it is hoped that...} devient \eng{we hope that...} à l'actif.
    \item \eng{Training programmes for youth are sponsored by the government.} (Des programmes de formation pour les jeunes sont financés par le gouvernement.) Present Simple passif, sujet pluriel.
    \item \eng{Electricity cuts are still experienced in many rural areas.} (Paragraphe 3.)
\end{enumerate}

\textbf{Points de vigilance :} ne pas traduire mot à mot ; vérifier l'accord de \eng{be} avec le nouveau sujet (\eng{programmes are}, pas \eng{is}) ; conserver le temps d'origine dans les transformations.
\end{corrige}

\begin{corrige}[Exercice 10 — Type Bac : Traduction Thématique]
\textbf{Barème indicatif (5 pts) :} exactitude grammaticale 2 pts (dont 1 pt pour les 4 passifs corrects) ; vocabulaire 1,5 pt ; respect du sens 1 pt ; fluidité 0,5 pt.

\textbf{Proposition de traduction :}
\begin{textebox}[Modèle de traduction]
The Africa Cup of Nations was successfully organised in Cameroon in 2022. Stadiums were renovated and new roads were built for the event. Thousands of tourists were welcomed in the cities of Yaoundé, Douala, Buea and Limbe. It is said that this event boosted the local economy. However, some projects were not completed on time. Efforts must be made to maintain this infrastructure.
\end{textebox}

\textbf{Occurrences passives :} \eng{was organised, were renovated, were built, were welcomed, it is said, were not completed, must be made} — 7 occurrences, bien au-delà des 4 exigées.

\textbf{Points de vigilance :}
\begin{itemize}
    \item « On dit que » $\rightarrow$ \eng{it is said that} (passif impersonnel), pas \eng{they say that} si l'on veut garder la voix passive.
    \item « Il faut que des efforts soient faits » $\rightarrow$ \eng{efforts must be made} (modal + \eng{be} + participe passé).
    \item \eng{boosted} reste à l'actif car le sujet \eng{this event} fait l'action.
    \item Orthographe (britannique) : \eng{successfully}, \eng{organised}, \eng{maintain}, \eng{infrastructure}.
\end{itemize}
\end{corrige}

\begin{corrige}[Exercice 11 — Type Bac : Sujet de Composition — Expression Écrite]
\textbf{Barème indicatif (10 pts) :} respect de la consigne et du format article 2 pts ; utilisation correcte d'au moins 4 passifs 3 pts ; variété des temps (3 minimum) 2 pts ; vocabulaire et richesse 2 pts ; organisation et cohérence 1 pt.

\textbf{Proposition de rédaction modèle :}
\begin{textebox}[A Memorable Prize Day at Our School]
Last Friday, a colourful prize-giving ceremony was organised in the main courtyard of our high school. The event was attended by the Divisional Officer, parents, teachers and all the students.

The ceremony was opened by the principal, who delivered an inspiring speech about hard work and discipline. Then, prizes were awarded to the best students of each class. Traditional dances were performed by the school cultural club, and songs were sung by the choir. Everyone was impressed by the talent displayed on stage.

Since last year, many improvements have been made in our school, and this celebration proved it. At the end of the day, trees were planted in the school garden to mark the occasion.

It was a wonderful day that will be remembered for a long time. Such events should be organised every year to encourage excellence among students.
\end{textebox}

\textbf{Analyse du modèle :}
\begin{itemize}
    \item \textbf{Passifs soulignés (11 occurrences) :} \eng{was organised, was attended, was opened, were awarded, were performed, were sung, was impressed, have been made, were planted, will be remembered, should be organised}.
    \item \textbf{Temps utilisés :} Past Simple (\eng{was organised, delivered}), Present Perfect (\eng{have been made}), Future Simple (\eng{will be remembered}), modal (\eng{should be organised}).
    \item \textbf{Structure :} titre accrocheur, introduction (quoi, quand, où), développement (déroulement), conclusion (impact et souhait).
\end{itemize}

\textbf{Points de vigilance pour les candidats :}
\begin{itemize}
    \item Souligner les passifs comme demandé dans la consigne.
    \item Compter les mots (120--150) : le modèle en contient environ 150.
    \item Éviter les passifs sur verbes intransitifs (\eng{happen, occur}).
    \item Soigner la date et le lieu dès la première phrase (technique journalistique : \emph{who, what, when, where}).
\end{itemize}
\end{corrige}

\newpage

\coursec{Fiche bilan}

\begin{popfigure}
\begin{tikzpicture}[
    mindmap,
    concept color=popblue,
    level 1 concept/.append style={level distance=3.5cm, sibling angle=360/7},
    level 2 concept/.append style={level distance=2.5cm},
    every node/.style={font=\small, align=center},
    branch1/.style={concept color=poporange},
    branch2/.style={concept color=popgreen},
    branch3/.style={concept color=poppurple},
    branch4/.style={concept color=poppink},
    branch5/.style={concept color=popgold},
    branch6/.style={concept color=popblue!70!popgreen},
    branch7/.style={concept color=popdark}
]
\node[concept, text=white, scale=1.2, align=center]{Active \& Passive\\Voices}
    child[branch1] { node[concept, text=white, align=center]{Formation\\ \texttt{be + V-en}} }
    child[branch2] { node[concept, text=white, align=center]{Tenses\\ Table} }
    child[branch3] { node[concept, text=white, align=center]{Agent\\ \texttt{by} / \texttt{with}} }
    child[branch4] { node[concept, text=white, align=center]{Uses\\ Focus / Formal} }
    child[branch5] { node[concept, text=white, align=center]{Double Object\\ Verbs} }
    child[branch6] { node[concept, text=white, align=center]{Modals\\ \texttt{be done}} }
    child[branch7] { node[concept, text=white, align=center]{Traps\\ FR vs EN} };
\end{tikzpicture}
\legende{Carte mentale : voix active et passive (EN56)}
\end{popfigure}

\begin{bilan}
\end{bilan}

\courssub{Lexique clé — Key Vocabulary}
\begin{center}
\begin{tabularx}{\linewidth}{|X|X|X|}\hline
\rowcolor{popblueL} \textbf{English} & \textbf{Français} & \textbf{Remarque} \\ \hline
\cle{Active voice} & Voix active & Sujet = agent (doer) \\ \hline
\cle{Passive voice} & Voix passive & Sujet = patient (receiver) \\ \hline
\cle{Agent} & Agent & Introduit par \eng{by} \\ \hline
\cle{Instrument} & Instrument & Introduit par \eng{with} \\ \hline
\cle{Past participle (V-en)} & Participe passé & Forme \cle{-ed} ou irrégulière \\ \hline
\cle{Transitive verb} & Verbe transitif & Prend un COD (seul convertible) \\ \hline
\cle{Intransitive verb} & Verbe intransitif & Pas de COD $\rightarrow$ pas de passif \\ \hline
\cle{Ditransitive verb} & Verbe ditransitif & Deux COD (ex. \eng{give, send}) \\ \hline
\cle{Impersonal passive} & Passif impersonnel & \eng{It is said that...} \\ \hline
\cle{Causative} & Causatif & \eng{have/get sth done} \\ \hline
\end{tabularx}
\end{center}

\courssub{Règles de grammaire à connaître par cœur}
\begin{center}
\begin{tabularx}{\linewidth}{|>{\bfseries}l|X|X|}\hline
\rowcolor{popblueL} \textbf{Temps / Modalité} & \textbf{Active} & \textbf{Passive (\texttt{be + V-en})} \\ \hline
Present Simple & \eng{They clean the streets.} & \eng{The streets \cle{are cleaned}.} \\ \hline
Past Simple & \eng{They cleaned...} & \eng{The streets \cle{were cleaned}.} \\ \hline
Present Perfect & \eng{They have cleaned...} & \eng{The streets \cle{have been cleaned}.} \\ \hline
Past Perfect & \eng{They had cleaned...} & \eng{The streets \cle{had been cleaned}.} \\ \hline
Future (will) & \eng{They will clean...} & \eng{The streets \cle{will be cleaned}.} \\ \hline
Modals (can/must...) & \eng{They must clean...} & \eng{The streets \cle{must be cleaned}.} \\ \hline
Present Continuous & \eng{They are cleaning...} & \eng{The streets \cle{are being cleaned}.} \\ \hline
Past Continuous & \eng{They were cleaning...} & \eng{The streets \cle{were being cleaned}.} \\ \hline
\eng{going to} & \eng{They are going to clean...} & \eng{The streets \cle{are going to be cleaned}.} \\ \hline
Gerund / Infinitive & \eng{cleaning / to clean} & \eng{being cleaned / to be cleaned} \\ \hline
\end{tabularx}
\end{center}

\courssub{Méthodes express}
\begin{methode}[Transformer Actif $\rightarrow$ Passif]
\begin{enumerate}
\item Identifier le \textbf{sujet}, le \textbf{verbe} (temps) et le \textbf{COD} de la phrase active.
\item Le \textbf{COD devient sujet} de la phrase passive.
\item Conjuguer l'auxiliaire \textbf{\texttt{be}} au \textbf{même temps} que le verbe actif.
\item Ajouter le \textbf{participe passé (V-en)} du verbe principal.
\item (Optionnel) Ajouter \eng{by + agent} si l'agent est pertinent / connu.
\end{enumerate}
\end{methode}

\begin{methode}[Choisir \texttt{by} ou \texttt{with}]
\begin{enumerate}
\item \texttt{\cle{by}} + \textbf{agent vivant / personne} : \eng{The book was written \cle{by} Chimamanda.}
\item \texttt{\cle{with}} + \textbf{instrument / matière / outil inanimé} : \eng{The door was opened \cle{with} a key.}
\item Verbes d'état / sentiment (\eng{known, loved, filled, covered}) $\rightarrow$ souvent \texttt{with} ou \texttt{by} selon contexte.
\end{enumerate}
\end{methode}

\begin{methode}[Verbes à deux objets (Ditransitifs)]
\begin{itemize}
\item Active : \eng{She gave \textbf{me} (IO) \textbf{a book} (DO).}
\item Passif 1 (standard) : \eng{\textbf{I} was given \textbf{a book}.} (IO $\rightarrow$ Sujet)
\item Passif 2 (moins fréquent) : \eng{\textbf{A book} was given \textbf{to me}.} (DO $\rightarrow$ Sujet + \texttt{to})
\item \textbf{Règle Bac} : Privilégier la transformation de l'\textbf{objet indirect (personne)} en sujet.
\end{itemize}
\end{methode}


\begin{aretenir}[Checklist avant le Bac]
\begin{itemize}
\item $\square$ Je sais identifier le COD pour en faire le sujet du passif.
\item $\square$ Je conjugue correctement \texttt{be} à tous les temps (surtout \texttt{been being} au perfect continuous -- rare mais possible).
\item $\square$ Je maîtrise les \textbf{participles passés irréguliers} (\eng{done, written, given, taken, made, seen}...).
\item $\square$ Je distingue \texttt{\cle{by}} (agent) et \texttt{\cle{with}} (instrument/matière).
\item $\square$ Je ne transforme \textbf{jamais} un verbe intransitif (\eng{happen, sleep, arrive, die, occur}) en passif.
\item $\square$ Je gère les \textbf{verbes à particule} : la particule reste collée au verbe (\eng{The meeting was \cle{put off}}).
\item $\square$ Je sais former les \textbf{deux passifs} des verbes ditransitifs (\eng{give, send, offer, tell}) et je choisis le sujet personne.
\item $\square$ J'utilise le passif \textbf{impersonnel} (\eng{It is reported that... / He is said to be...}) pour le style journalistique / académique.
\item $\square$ Je distingue \textbf{passif} (\eng{The car was repaired}) et \textbf{causatif} (\eng{I had the car repaired} = j'ai fait réparer).
\item $\square$ Je n'oublie pas \eng{\cle{to be born}} (toujours passif au sens propre : \eng{I \cle{was born} in Buea}).
\item $\square$ Je repère les \textbf{faux amis} : \eng{actually} (\emph{en fait}), \eng{eventually} (\emph{finalement}), \eng{sensible} (\emph{raisonnable}).
\end{itemize}
\end{aretenir}

\findecours

\end{document}
